\documentclass[aps,preprint]{revtex4}%
\usepackage{amsfonts}
\usepackage{amsmath}
\usepackage{amssymb}
\usepackage{graphicx}%
\setcounter{MaxMatrixCols}{30}
%TCIDATA{OutputFilter=latex2.dll}
%TCIDATA{Version=4.10.0.2363}
%TCIDATA{CSTFile=revtex4.cst}
%TCIDATA{Created=Monday, September 20, 2004 09:58:21}
%TCIDATA{LastRevised=Wednesday, June 29, 2005 10:54:05}
%TCIDATA{<META NAME="GraphicsSave" CONTENT="32">}
%TCIDATA{<META NAME="DocumentShell" CONTENT="Articles\SW\REVTeX 4">}
%TCIDATA{Language=American English}
\newtheorem{theorem}{Theorem}
\newtheorem{acknowledgement}[theorem]{Acknowledgement}
\newtheorem{algorithm}[theorem]{Algorithm}
\newtheorem{axiom}[theorem]{Axiom}
\newtheorem{claim}[theorem]{Claim}
\newtheorem{conclusion}[theorem]{Conclusion}
\newtheorem{condition}[theorem]{Condition}
\newtheorem{conjecture}[theorem]{Conjecture}
\newtheorem{corollary}[theorem]{Corollary}
\newtheorem{criterion}[theorem]{Criterion}
\newtheorem{definition}[theorem]{Definition}
\newtheorem{example}[theorem]{Example}
\newtheorem{exercise}[theorem]{Exercise}
\newtheorem{lemma}[theorem]{Lemma}
\newtheorem{notation}[theorem]{Notation}
\newtheorem{problem}[theorem]{Problem}
\newtheorem{proposition}[theorem]{Proposition}
\newtheorem{remark}[theorem]{Remark}
\newtheorem{solution}[theorem]{Solution}
\newtheorem{summary}[theorem]{Summary}
\newenvironment{proof}[1][Proof]{\noindent\textbf{#1.} }{\ \rule{0.5em}{0.5em}}
\begin{document}
\preprint{ }
\title[Janaks Theorem]{Generalized Janaks Theorem}
\author{B. Weiner}
\affiliation{Pennsylvania State University}
\keywords{one two three}
\pacs{PACS number}

\begin{abstract}
Shell document for REV\TeX{} 4.

\end{abstract}
\date{12/21/2004}
\startpage{1}
\maketitle
\tableofcontents


\section{Introduction\label{Intro}}

\bigskip

\section{Standard Janaks Theorem\label{Standard}}

\subsection{Hohenberg-Kohn DFT}

\bigskip Density Functional Theory is based on the existence of a universal
functional, $\mathcal{\tilde{F}}\left(  \rho\right)  $ \cite{HK} of the charge
density $\rho\left(  \mathbf{r}\right)  ,$ where
\begin{equation}
\rho\left(  \mathbf{r}\right)  =\int D^{1}\left(  \mathbf{r,\xi;r,\xi}\right)
d\mathbf{\xi}%
\end{equation}
and $D^{1}\left(  \mathbf{r,\xi;r}^{\prime}\mathbf{,\xi}^{\prime}\right)  $ is
the integral kernel of the First Order Reduced Operator (FORDO) $D^{1}$
\cite{Coleman} (and see Appendix \ref{DenMat} for details of density and
reduced density operators and matrices), which can be expanded as
\begin{equation}
D^{1}=\int D^{1}\left(  \mathbf{r,\xi;r}^{\prime}\mathbf{,\xi}^{\prime
}\right)  \left\vert \mathbf{r,\xi}\right\rangle \left\langle \mathbf{r}%
^{\prime}\mathbf{,\xi}^{\prime}\right\vert d\mathbf{r}d\mathbf{r}^{\prime
}d\mathbf{\xi}d\mathbf{\xi}^{\prime}%
\end{equation}
in the basis $\left\{  \left\vert \mathbf{r,\xi}\right\rangle \right\}  $
formed by the tensor product of basis of generalized eigenvectors, $\left\{
\left\vert \mathbf{r}\right\rangle \right\}  ,$ of the position operator and
the family of $SU(2)$-coherent spin states $\left\{  \left\vert \mathbf{\xi
}\right\rangle \right\}  $. In order to consider spin dependent effects it is
convenient to extend the original formulation of DFT in terms of $\rho\left(
\mathbf{r}\right)  $ and consider a universal functional, $\mathcal{F}\left(
\gamma\right)  $, of the space-spin density $\gamma\left(  \mathbf{r,\xi
}\right)  =$ $D^{1}\left(  \mathbf{r,\xi;r,\xi}\right)  $
\cite{weiner-trickey} and local external potentials of the position variable
$\mathbf{r}$ and spin space variable $\mathbf{\xi}$ that describe the
interaction with external electric and magnetic fields (the integral kernel of
the magnetic dipole operator is local (i.e. diagonal) in the coherent state
representation but not over the spectrum of the spin operator)$.$

For a system that contains on average $N$-electrons the space spin density
must satisfy the condition
\begin{equation}
\int\gamma\left(  \mathbf{r,\xi}\right)  d\mathbf{r}d\mathbf{\xi}=N,
\label{Number}%
\end{equation}
which does not restrict the density to be the contraction of a pure
$N$-electron state and one could consider incoherent and coherent mixtures of
sates describing different number of electrons. Focusing on pure $N$-electron
states one can define convex sets, $\left\{  \left[  \gamma\right]
_{N}\right\}  ,$ of $N$-electron state density operators by
\begin{equation}
\left[  \gamma\right]  _{N}=\left\{  \left.  D^{N}\right\vert \Xi\left(
D^{N}\right)  =\gamma;D^{N}\in S_{N}\subset\mathcal{B}_{1}\left(
\mathcal{H}^{N}\right)  \right\}
\end{equation}
(see Appendix \ref{DenMat} for the definitions of the linear map $\Xi$ and the
domain $S_{N}$). If condition eq. $\left(  \ref{Number}\right)  $ is satisfied
and $D^{N}\in\mathcal{B}_{1}\left(  \mathcal{H}^{N}\right)  $ then the state
normalization
\begin{equation}
\operatorname{Tr}D^{N}=1
\end{equation}
is implied by this condition and $D^{N}$ needs only to be constrained to lie
in the cone of positive operators of $\mathcal{B}_{1}\left(  \mathcal{H}%
^{N}\right)  $. $\lceil$If $D^{N}\in\mathcal{B}_{2}\left(  \mathcal{H}%
^{N}\right)  $ then
\begin{equation}
\operatorname{Tr}\left\{  \mathcal{N}D^{N}\right\}  =N\operatorname{Tr}%
\left\{  D^{N}\right\}
\end{equation}
as
\begin{equation}
\mathcal{N}\left\vert \Psi\right\rangle =N\left\vert \Psi\right\rangle
\quad\forall\left\vert \Psi\right\rangle \in\mathcal{H}^{N}%
\end{equation}
and
\begin{equation}
\operatorname{Tr}\left\{  \mathcal{N}D^{N}\right\}  =\int\gamma\left(
\mathbf{r,\xi}\right)  d\mathbf{r}d\mathbf{\xi.}%
\end{equation}
Thus
\begin{equation}
N\operatorname{Tr}\left\{  D^{N}\right\}  =N\Rightarrow\operatorname{Tr}%
\left\{  D^{N}\right\}  =1.\rfloor
\end{equation}
Thus $D^{N}$ can always be expressed in the form
\begin{equation}
D^{N}=QQ^{\dagger},
\end{equation}
where $\left(  \text{Appendix \ref{Operators}}\right)  $
\begin{equation}
Q\in\mathcal{B}_{2}\left(  \mathcal{H}^{N}\right)  .
\end{equation}
However this decomposition is not unique as in particular
\begin{equation}
D^{N}=QUU^{\dagger}Q^{\dagger}%
\end{equation}
for any unitary operator $U$. One should also note that $Q$ can always be
chosen to self adjoint and even positive.

Even though the explicit form of the functional $\mathcal{F}_{HK}\left(
\gamma\right)  $ is not known one can define it implicitly \cite{LevyLieb}
\begin{equation}
\mathcal{F}_{HK}\left(  \gamma\right)  =\operatorname{Tr}\left\{  \overset
{}{\left(  K+V_{2}\right)  D_{\ast}^{N}\left(  \gamma\right)  }\right\}
=\min_{D^{N}\in\left[  \mathbf{\gamma}\right]  _{N}}\operatorname{Tr}\left\{
\overset{}{\left(  K+V_{2}\right)  D^{N}}\right\}  , \label{HK}%
\end{equation}
where $K$ is the kinetic energy operator and $V_{2}$ the coulomb operator. The
Lagrangian associated with the constrained linear optimization problem eq.
$\left(  \ref{HK}\right)  $ is
\begin{equation}
\mathcal{L}\left(  D^{N},\lambda,\gamma\right)  =\mathcal{E}_{U}\left(
D^{N}\right)  -\left\langle \left.  \lambda\right\vert \left\{  \Xi\left(
D^{N}\right)  -\gamma\right\}  \right\rangle ,
\end{equation}
where%
\begin{equation}
\mathcal{E}_{U}\left(  D^{N}\right)  =\operatorname{Tr}\left\{  \overset
{}{\left(  K+V_{2}\right)  D^{N}}\right\}  ,
\end{equation}
\begin{equation}
\left\langle \left.  \lambda\right\vert \left\{  \Xi\left(  D^{N}\right)
-\gamma\right\}  \right\rangle =\int\lambda\left(  \mathbf{r,\xi}\right)
\left\{  \Xi\left(  D^{N}\right)  \left(  \mathbf{r,\xi}\right)
-\gamma\left(  \mathbf{r,\xi}\right)  \right\}  d\mathbf{r}d\mathbf{\xi}%
\end{equation}
and $\gamma$ is chosen to belong to the convex set $\mathcal{P}_{N}$ given by
\begin{equation}
\mathcal{P}_{N}=\left\{  f\left\vert f\in L_{1}\left(  \mathbb{R}^{3}%
\times\mathbb{C}\right)  ;\,f\left(  \mathbf{r,}z\right)  \geq0;\int f\left(
\mathbf{r,}z\right)  d\mathbf{r}dz=N\right.  \right\}  .
\end{equation}


$\lceil$Let $D\in\mathcal{B}_{2}\left(  \mathcal{H}_{F}\right)  $ then
\begin{align}
\operatorname{Tr}\left\{  \left(  \mathcal{N}-\operatorname{Tr}\left\{
\mathcal{ND}\right\}  \right)  ^{2}\right\}   &  =\operatorname{Tr}\left\{
\left(  \mathcal{N}^{2}-2\mathcal{N}\operatorname{Tr}\left\{  \mathcal{N}%
D\right\}  +\left(  \operatorname{Tr}\mathcal{N}D\right)  ^{2}\right)
D\right\} \nonumber\\
&  =\operatorname{Tr}\left\{  \mathcal{N}^{2}D\right\}  -2\left(
\operatorname{Tr}\left\{  \mathcal{N}D\right\}  \right)  ^{2}+\left(
\operatorname{Tr}\left\{  \mathcal{N}D\right\}  \right)  ^{2}\operatorname{Tr}%
\left\{  D\right\}  .
\end{align}
If
\begin{equation}
\operatorname{Tr}\left\{  D\right\}  =1
\end{equation}
then
\begin{equation}
\operatorname{Tr}\left\{  \left(  \mathcal{N}-\operatorname{Tr}\left\{
\mathcal{ND}\right\}  \right)  ^{2}\right\}  =\operatorname{Tr}\left\{
\mathcal{N}^{2}D\right\}  -\left(  \operatorname{Tr}\left\{  \mathcal{N}%
D\right\}  \right)  ^{2}.
\end{equation}
If further
\begin{equation}
\operatorname{Tr}\left\{  \left(  \mathcal{N}-\operatorname{Tr}\left\{
\mathcal{ND}\right\}  \right)  ^{2}\right\}  =0
\end{equation}
then
\begin{equation}
\operatorname{Tr}\left\{  \mathcal{N}^{2}D\right\}  =\left(  \operatorname{Tr}%
\left\{  \mathcal{N}D\right\}  \right)  ^{2}.
\end{equation}
and one can show that $\left[  \mathcal{N},D\right]  =0.$ This coupled with
$\operatorname{Tr}D^{1}=N$ gives that $D\in\mathcal{B}_{1}\left(
\mathcal{H}^{N}\right)  .\rfloor$

The constrained linear optimization problem eq. $\left(  \ref{HK}\right)  $
can be expressed as a constrained quadratic optimization problem
\begin{equation}
\mathcal{F}_{HK}\left(  \gamma\right)  =\operatorname{Tr}\left\{  \overset
{}{\left(  K+V_{2}\right)  Q_{\ast}\left(  \gamma\right)  Q_{\ast}^{\dagger
}\left(  \gamma\right)  }\right\}  =\min_{QQ^{\dagger}\in\left[
\mathbf{\gamma}\right]  _{N}}\operatorname{Tr}\left\{  \overset{}{\left(
K+V_{2}\right)  QQ^{\dagger}}\right\}  ,
\end{equation}
that has Lagrangian
\begin{equation}
\mathcal{L}\left(  Q,Q^{\dagger},\lambda,\gamma\right)  =\mathcal{E}%
_{U}\left(  Q,Q^{\dagger}\right)  -\left\langle \left.  \lambda\right\vert
\left\{  \Xi\left(  Q,Q^{\dagger}\right)  -\gamma\right\}  \right\rangle ,
\end{equation}
where
\begin{equation}
\Xi\left(  Q,Q^{\dagger}\right)  =\Xi\left(  QQ^{\dagger}\right)  .
\end{equation}
The normalization constraint
\begin{equation}
\operatorname{Tr}\left\{  QQ^{\dagger}\right\}  =1
\end{equation}
is implied by $Q\in\mathcal{B}_{2}\left(  \mathcal{H}^{N}\right)  $ and
$\int\gamma\left(  \mathbf{x}\right)  d\mathbf{x}=N.$ It can be useful to note
that $\mathcal{E}_{U}$ can be recast as in the form of a inner product
\begin{equation}
\mathcal{E}_{U}\left(  Q,Q^{\dagger}\right)  =\left(  \left.  Q\right\vert
\left(  \widehat{K+V_{2}}\right)  Q\right)  =\operatorname{Tr}\left\{
Q^{\dagger}\left(  K+V_{2}\right)  Q\right\}  ,
\end{equation}
where $\widehat{K+V_{2}}$ denotes the super operator $:\mathcal{B}_{2}\left(
\mathcal{H}^{N}\right)  \rightarrow\mathcal{B}_{2}\left(  \mathcal{H}%
^{N}\right)  $ given by
\begin{equation}
\widehat{K+V_{2}}\left\vert Q\right)  =\mathcal{P}_{\mathcal{B}_{2}\left(
\mathcal{H}^{N}\right)  }\left\vert \left(  K+V_{2}\right)  Q\right)  .
\end{equation}
$.$ The contraction map $\Xi$ can similarly be expressed as
\begin{equation}
\Xi\left(  Q,Q^{\dagger}\right)  \left(  \mathbf{x}\right)  =\left(
Q\left\vert \widehat{\Phi\left(  \mathbf{x}\right)  ^{\dagger}\Phi\left(
\mathbf{x}\right)  }\right.  Q\right)  =\operatorname{Tr}\left\{  Q^{\dagger
}\Phi\left(  \mathbf{x}\right)  ^{\dagger}\Phi\left(  \mathbf{x}\right)
Q\right\}  ,
\end{equation}
which gives
\begin{equation}
\mathcal{L}\left(  Q,Q^{\dagger},\lambda,\gamma\right)  =\left(  Q\left\vert
\left(  \widehat{K+V_{2}}\right)  Q\right.  \right)  -\left\langle
\lambda\left\vert \left\{  \left(  Q\left\vert \widehat{\Phi^{\dagger}\Phi
}Q\right.  \right)  -\gamma\right\}  \right.  \right\rangle .
\end{equation}
The first order KKT conditions \cite{KKT} are%
\begin{equation}
D_{Q,\bar{Q},\lambda}\mathcal{L}\left(  Q,Q^{\dagger},\lambda,\gamma\right)
=\mathbf{0}%
\end{equation}
where
\begin{align}
&  D_{Q,\bar{Q},\lambda}\mathcal{L}\left(  Q,Q^{\dagger},\lambda
,\gamma\right)  =\left(  D_{Q}\mathcal{L},D_{\bar{Q}}\mathcal{L},D_{\lambda
}\mathcal{L}\right) \nonumber\\
D_{Q,\bar{Q},\lambda}\mathcal{L}  &  :\mathbb{C}^{2r_{N}^{2}}\times L_{\infty
}\left(  \mathbb{R}^{3}\times\mathbb{C}\right)  \times L_{1}\left(
\mathbb{R}^{3}\times\mathbb{C}\right)  \rightarrow\mathbb{C}^{2r_{N}^{2}%
}\times L_{1}\left(  \mathbb{R}^{3}\times\mathbb{C}\right)
\end{align}
and%
\begin{subequations}
\begin{align}
D_{Q}\mathcal{L}\left(  Q,Q^{\dagger},\lambda,\gamma\right)   &
=D_{Q}\mathcal{E}_{U}\left(  Q,Q^{\dagger}\right)  -\left\langle
\lambda\left\vert D_{Q}\Xi\left(  Q,Q^{\dagger}\right)  \right.  \right\rangle
\\
D_{\bar{Q}}\mathcal{L}\left(  Q,Q^{\dagger},\lambda,\gamma\right)   &
=D_{\bar{Q}}\mathcal{E}_{U}\left(  Q,Q^{\dagger}\right)  -\left\langle
\lambda\left\vert D_{\bar{Q}}\Xi\left(  Q,Q^{\dagger}\right)  \right.
\right\rangle \\
D_{\lambda}\mathcal{L}\left(  Q,Q^{\dagger},\lambda,\gamma\right)   &
=\Xi\left(  Q,Q^{\dagger}\right)  -\gamma.
\end{align}
In particular this gives
\end{subequations}
\begin{equation}
D_{Q}\mathcal{L}\left(  Q,Q^{\dagger},\lambda,\gamma\right)  =\left(
\begin{array}
[c]{cc}%
D_{Q}\mathcal{E}_{U} & D_{\bar{Q}}\mathcal{E}_{U}%
\end{array}
\right)  -\int\lambda\left(  \mathbf{x}\right)  \left(
\begin{array}
[c]{cc}%
D_{Q}\Xi\left(  \mathbf{x}\right)  & D_{\bar{Q}}\Xi\left(  \mathbf{x}\right)
\end{array}
\right)  d\mathbf{x.}%
\end{equation}
Each derivative is given by
\begin{equation}
D_{Q}\mathcal{E}_{U}=\left(  \partial_{Q_{11}}\mathcal{E}_{U},\cdots
,\partial_{Q_{r_{N}r_{N}}}\mathcal{E}_{U}\right)
\end{equation}
where we have utilized the expansion
\begin{equation}
Q=\sum_{1\leq k,l\leq r_{N}}Q_{kl}\left\vert \Phi_{k}\right\rangle
\left\langle \Phi_{l}\right\vert
\end{equation}
in terms of the basis $\left\{  \left.  \left\vert \Phi_{k}\right\rangle
\left\langle \Phi_{l}\right\vert \right\vert 1\leq k,l\leq r_{N}\right\}  $ of
$\mathcal{B}_{2}\left(  \mathcal{H}^{N}\right)  ,$ where $\left\{  \left.
\left\vert \Phi_{k}\right\rangle \right\vert 1\leq k\leq r_{N}\right\}  $ is a
conb for $\mathcal{H}^{N}.$ This gives
\begin{align}
\mathcal{E}_{U}  &  :\mathbb{C}^{r_{N}^{2}}\times\mathbb{C}^{r_{N}^{2}%
}\rightarrow\mathbb{R}\nonumber\\
\mathcal{E}_{U}\left(  Q,\bar{Q}\right)   &  =\sum_{1\leq k,l,m,n\leq r_{N}%
}\operatorname{Tr}\left\{  \bar{Q}_{kl}\left\vert \Phi_{l}\right\rangle
\left\langle \Phi_{k}\right\vert H\sum_{1\leq m,n\leq r_{N}}Q_{mn}\left\vert
\Phi_{m}\right\rangle \left\langle \Phi_{n}\right\vert \right\} \nonumber\\
&  =\sum_{1\leq k,l,m\leq r_{N}}\bar{Q}_{kl}H_{km}Q_{ml}%
\end{align}
and
\begin{align}
D_{Q\bar{Q}}\mathcal{E}_{U}\left(  Q,\bar{Q}\right)   &  :\mathbb{C}%
^{r_{N}^{2}}\times\mathbb{C}^{r_{N}^{2}}\rightarrow\mathbb{R}\nonumber\\
\partial_{Q_{ml}}\mathcal{E}_{U}\left(  Q,\bar{Q}\right)   &  =\sum_{1\leq
k\leq r_{N}}\bar{Q}_{kl}H_{km}=\sum_{1\leq k\leq r_{N}}\bar{H}_{mk}\bar
{Q}_{kl}\nonumber\\
\partial_{\bar{Q}_{ml}}\mathcal{E}_{U}\left(  Q,\bar{Q}\right)   &
=\sum_{1\leq m\leq r_{N}}H_{mk}Q_{kl}=\overline{\partial_{Q_{ml}}%
\mathcal{E}_{U}\left(  Q,\bar{Q}\right)  }.
\end{align}
The constraint functional $\Xi\left(  QQ^{\dagger}\right)  $ is then given by
\begin{align}
\Xi &  :\mathbb{C}^{r_{N}^{2}}\times\mathbb{C}^{r_{N}^{2}}\rightarrow
L_{1}\left(  \mathbb{R}^{3}\mathbb{\times C}\right) \nonumber\\
\Xi\left(  QQ^{\dagger}\right)  \left(  \mathbf{x}\right)   &  =\sum_{1\leq
k,l,m,n\leq r_{N}}\bar{Q}_{kl}Q_{mn}\operatorname{Tr}\left\{  \left\vert
\Phi_{l}\right\rangle \left\langle \Phi_{k}\right\vert \Phi\left(
\mathbf{x}\right)  ^{\dagger}\Phi\left(  \mathbf{x}\right)  \left\vert
\Phi_{m}\right\rangle \left\langle \Phi_{n}\right\vert \right\} \nonumber\\
&  =\sum_{1\leq k,l,m\leq r_{N}}\bar{Q}_{kl}D_{km}^{1}\left(  \mathbf{x}%
\right)  Q_{ml}%
\end{align}
leading to
\begin{align}
D_{Q\bar{Q}}\Xi\left(  QQ^{\dagger}\right)   &  :\mathbb{C}^{r_{N}^{2}}%
\times\mathbb{C}^{r_{N}^{2}}\rightarrow L_{1}\left(  \mathbb{R}^{3}%
\mathbb{\times C}\right) \nonumber\\
\partial_{Q_{ml}}\Xi\left(  QQ^{\dagger}\right)  \left(  \mathbf{x}\right)
&  =\sum_{1\leq k\leq r_{N}}\bar{Q}_{kl}D_{km}^{1}\left(  \mathbf{x}\right)
=\sum_{1\leq k\leq r_{N}}\bar{D}_{mk}^{1}\left(  \mathbf{x}\right)  \bar
{Q}_{kl}\nonumber\\
\partial_{\bar{Q}_{ml}}\Xi\left(  QQ^{\dagger}\right)  \left(  \mathbf{x}%
\right)   &  =\sum_{1\leq m\leq r_{N}}D_{mk}^{1}\left(  \mathbf{x}\right)
Q_{kl}=\overline{\partial_{Q_{ml}}\Xi\left(  QQ^{\dagger}\right)  \left(
\mathbf{x}\right)  }.
\end{align}
The remaining components of the gradient are given by
\begin{align}
\lambda &  :\mathbb{C}^{r_{N}^{2}}\times\mathbb{C}^{r_{N}^{2}}\rightarrow
\mathbb{R}\nonumber\\
\left\langle \left.  \lambda\right\vert \left\{  \Xi\left(  Q,Q^{\dagger
}\right)  -\gamma\right\}  \right\rangle  &  =\int\lambda\left(
\mathbf{x}\right)  \left\{  \Xi\left(  Q,Q^{\dagger}\right)  \left(
\mathbf{x}\right)  -\gamma\left(  \mathbf{x}\right)  \right\}  d\mathbf{x}%
\end{align}
and
\begin{align}
D_{\lambda}\left\langle \left.  \lambda\right\vert \left\{  \Xi\left(
Q,Q^{\dagger}\right)  -\gamma\right\}  \right\rangle  &  :L_{\infty}\left(
\mathbb{R}^{3}\mathbb{\times C}\right)  \times\mathbb{C}^{r_{N}^{2}}%
\times\mathbb{C}^{r_{N}^{2}}\rightarrow L_{1}\left(  \mathbb{R}^{3}%
\mathbb{\times C}\right) \nonumber\\
\frac{\delta\left\langle \left.  \lambda\right\vert \left\{  \Xi\left(
Q,Q^{\dagger}\right)  -\gamma\right\}  \right\rangle }{\delta\lambda\left(
\mathbf{x}\right)  }  &  =\sum_{1\leq k,l,m\leq r_{N}}\bar{Q}_{kl}D_{km}%
^{1}\left(  \mathbf{x}\right)  Q_{ml}-\gamma\left(  \mathbf{x}\right)  .
\end{align}
Hence
\begin{align}
&  \left[  D_{Q}\mathcal{L}\left(  Q,Q^{\dagger},\lambda,\gamma\right)
\right]  _{ml}=\sum_{1\leq k\leq r_{N}}\left\{  \bar{H}_{mk}-\int
\lambda\left(  \mathbf{x}\right)  \bar{D}_{mk}^{1}\left(  \mathbf{x}\right)
d\mathbf{x}\right\}  \bar{Q}_{kl}=0\nonumber\\
&  \left[  D_{\bar{Q}}\mathcal{L}\left(  Q,Q^{\dagger},\lambda,\gamma\right)
\right]  _{ml}=\sum_{1\leq k\leq r_{N}}\left\{  H_{mk}-\int\lambda\left(
\mathbf{x}\right)  D_{mk}^{1}\left(  \mathbf{x}\right)  \right\}
Q_{kl}=0\nonumber\\
&  D_{\lambda}\mathcal{L}\left(  Q,Q^{\dagger},\lambda,\gamma\right)  \left(
\mathbf{x}\right)  =\sum_{1\leq k,l,m\leq r_{N}}\bar{Q}_{kl}D_{km}^{1}\left(
\mathbf{x}\right)  Q_{ml}-\gamma\left(  \mathbf{x}\right)  \overset{a.e}{=}0.
\end{align}
$\lceil$Compare to the discrete constraint case
\begin{equation}
\mathcal{L}\left(  Q,Q^{\dagger},\lambda,\gamma\right)  =\left\{  \sum_{1\leq
k,l,m\leq r_{N}}\bar{Q}_{kl}H_{km}Q_{ml}-\sum_{1\leq j\leq r}\lambda
_{j}\left\{  \bar{Q}_{kl}D_{kmj}^{1}Q_{ml}-\gamma_{j}\right\}  \right\}
\end{equation}
then
\begin{align}
&  D_{Q,\bar{Q},\lambda}\mathcal{L}\left(  Q,Q^{\dagger},\lambda
,\gamma\right)  =\nonumber\\
&  \left(  \left(  \partial_{Q_{11}}\mathcal{L}\left(  Q,Q^{\dagger}%
,\lambda,\gamma\right)  ,\cdots,\partial_{Q_{r_{N}r_{N}}}\mathcal{L}\left(
Q,Q^{\dagger},\lambda,\gamma\right)  \right)  ,\right. \nonumber\\
&  \left(  \partial_{\bar{Q}_{11}}\mathcal{L}\left(  Q,Q^{\dagger}%
,\lambda,\gamma\right)  ,\cdots,\partial_{\bar{Q}_{r_{N}r_{N}}}\mathcal{L}%
\left(  Q,Q^{\dagger},\lambda,\gamma\right)  \right)  ,\left.  \partial
_{\lambda_{1}}\mathcal{L}\left(  Q,Q^{\dagger},\lambda,\gamma\right)
,\cdots,\partial_{\lambda_{r}}\mathcal{L}\left(  Q,Q^{\dagger},\lambda
,\gamma\right)  \right)
\end{align}
which gives in the continuous case
\begin{align}
&  D_{Q,\bar{Q},\lambda}\mathcal{L}\left(  Q,Q^{\dagger},\lambda
,\gamma\right)  =\nonumber\\
&  \left(  \left(  \partial_{Q_{11}}\mathcal{L}\left(  Q,Q^{\dagger}%
,\lambda,\gamma\right)  ,\cdots,\partial_{Q_{r_{N}r_{N}}}\mathcal{L}\left(
Q,Q^{\dagger},\lambda,\gamma\right)  \right)  \right.  ,\nonumber\\
&  \left(  \partial_{\bar{Q}_{11}}\mathcal{L}\left(  Q,Q^{\dagger}%
,\lambda,\gamma\right)  ,\cdots,\partial_{\bar{Q}_{r_{N}r_{N}}}\mathcal{L}%
\left(  Q,Q^{\dagger},\lambda,\gamma\right)  \right)  ,\left.  \cdots
,\frac{\delta\mathcal{L}\left(  Q,Q^{\dagger},\lambda,\gamma\right)  }%
{\delta\lambda\left(  \mathbf{x}\right)  },\cdots\right)  \rfloor
\end{align}


The Hessian of the Lagrangian is
\begin{equation}
D_{Q\bar{Q}\lambda}^{2}\mathcal{L}\left(  Q,\bar{Q},\lambda,\gamma\right)
=\left(
\begin{array}
[c]{ccc}%
\mathbb{O} & D_{Q\bar{Q}}^{2}\mathcal{E}_{U}-\left\langle \lambda\left\vert
\partial_{Q\bar{Q}}^{2}\Xi\right.  \right\rangle  & \left(  \partial_{Q}%
\Xi\right)  ^{t}\\
D_{\bar{Q}Q}^{2}\mathcal{E}_{U}-\left\langle \lambda\left\vert \partial
_{\bar{Q}Q}^{2}\Xi\right.  \right\rangle  & \mathbb{O} & \left(
\partial_{\bar{Q}}\Xi\right)  ^{t}\\
\partial_{Q}\Xi & \partial_{\bar{Q}}\Xi & \mathbb{O}%
\end{array}
\right)  ,
\end{equation}
which is formed from the Hessian
\begin{equation}
D_{Q}^{2}\mathcal{L}\left(  Q,\bar{Q},\lambda,\gamma\right)  =\left(
\begin{array}
[c]{cc}%
\mathbb{O} & D_{Q\bar{Q}}^{2}\mathcal{E}_{U}-\left\langle \lambda\left\vert
\partial_{Q\bar{Q}}^{2}\Xi\right.  \right\rangle \\
D_{\bar{Q}Q}^{2}\mathcal{E}_{U}-\left\langle \lambda\left\vert \partial
_{\bar{Q}Q}^{2}\Xi\right.  \right\rangle  & \mathbb{O}%
\end{array}
\right)
\end{equation}
and the gradient terms.

In terms of the inner product and super operators we can expand the
Lagrangian
\begin{equation}
\mathcal{L}\left(  Q,Q^{\dagger},\lambda,\gamma\right)  =\left(  Q\left\vert
\left(  \widehat{K+V_{2}}\right)  Q\right.  \right)  -\left\langle
\lambda\left\vert \left\{  \left(  Q\left\vert \widehat{\Phi^{\dagger}\Phi
}Q\right.  \right)  -\gamma\right\}  \right.  \right\rangle
\end{equation}
as
\begin{align}
\mathcal{L}\left(  Q+\delta Q,Q^{\dagger}+\delta Q^{\dagger},\lambda
+\delta\lambda,\gamma\right)   &  =\left(  Q+\delta Q\left\vert \left(
\widehat{K+V_{2}}\right)  Q+\delta Q\right.  \right) \nonumber\\
&  -\left\langle \lambda+\delta\lambda\left\vert \left\{  \left(  Q+\delta
Q\left\vert \widehat{\Phi^{\dagger}\Phi}Q+\delta Q\right.  \right)
-\gamma\right\}  \right.  \right\rangle
\end{align}
giving
\begin{align}
&  \left(  Q\left\vert \left(  \widehat{K+V_{2}}\right)  Q\right.  \right)
+\left(  \delta Q\left\vert \left(  \widehat{K+V_{2}}\right)  Q\right.
\right)  +\left(  Q\left\vert \left(  \widehat{K+V_{2}}\right)  \delta
Q\right.  \right)  +\left(  \delta Q\left\vert \left(  \widehat{K+V_{2}%
}\right)  \delta Q\right.  \right) \nonumber\\
&  -\left\{  \left\langle \lambda\left\vert \left\{  \left(  Q\left\vert
\widehat{\Phi^{\dagger}\Phi}Q\right.  \right)  +\left(  \delta Q\left\vert
\widehat{\Phi^{\dagger}\Phi}Q\right.  \right)  +\left(  Q\left\vert
\widehat{\Phi^{\dagger}\Phi}\delta Q\right.  \right)  +\left(  \delta
Q\left\vert \widehat{\Phi^{\dagger}\Phi}\delta Q\right.  \right)
-\gamma\right\}  \right.  \right\rangle \right. \nonumber\\
&  \left.  +\left\langle \delta\lambda\left\vert \left\{  \left(  Q\left\vert
\widehat{\Phi^{\dagger}\Phi}Q\right.  \right)  +\left(  \delta Q\left\vert
\widehat{\Phi^{\dagger}\Phi}Q\right.  \right)  +\left(  Q\left\vert
\widehat{\Phi^{\dagger}\Phi}\delta Q\right.  \right)  +\left(  \delta
Q\left\vert \widehat{\Phi^{\dagger}\Phi}\delta Q\right.  \right)
-\gamma\right\}  \right.  \right\rangle \right\}  .
\end{align}
Thus
\begin{align}
&  D\mathcal{L}\left(  Q,Q^{\dagger},\lambda,\gamma\right)  \left(  \delta
Q,\delta Q^{\dagger},\delta\lambda\right)  =\left(  \delta Q\left\vert \left(
\widehat{K+V_{2}}\right)  Q\right.  \right)  +\left(  Q\left\vert \left(
\widehat{K+V_{2}}\right)  \delta Q\right.  \right) \nonumber\\
&  -\left\{  \left\langle \lambda\left\vert \left\{  \left(  \delta
Q\left\vert \widehat{\Phi^{\dagger}\Phi}Q\right.  \right)  +\left(
Q\left\vert \widehat{\Phi^{\dagger}\Phi}\delta Q\right.  \right)  \right\}
\right.  \right\rangle +\left\langle \delta\lambda\left\vert \left\{  \left(
Q\left\vert \widehat{\Phi^{\dagger}\Phi}Q\right.  \right)  -\gamma\right\}
\right.  \right\rangle \right\} \nonumber\\
&  =\overline{\left(  Q\left\vert \left(  \widehat{K+V_{2}}\right)  \delta
Q\right.  \right)  }+\left(  Q\left\vert \left(  \widehat{K+V_{2}}\right)
\delta Q\right.  \right) \nonumber\\
&  -\left\{  \left\langle \lambda\left\vert \left\{  \overline{\left(
Q\left\vert \widehat{\Phi^{\dagger}\Phi}\delta Q\right.  \right)  }+\left(
Q\left\vert \widehat{\Phi^{\dagger}\Phi}\delta Q\right.  \right)  \right\}
\right.  \right\rangle +\left\langle \delta\lambda\left\vert \left\{  \left(
Q\left\vert \widehat{\Phi^{\dagger}\Phi}Q\right.  \right)  -\gamma\right\}
\right.  \right\rangle \right\}
\end{align}
$\lceil$ Note
\begin{equation}
\overline{\left(  Q\left\vert \left(  \widehat{A}\right)  \delta Q\right.
\right)  }=\overline{\operatorname{Tr}\left\{  Q^{\dagger}A\delta Q\right\}
}=\operatorname{Tr}\left\{  \delta Q^{\dagger}AQ\right\}  =\left(  \delta
Q\left\vert \widehat{A}Q\right.  \right)  \rfloor
\end{equation}
by considering $\delta Q$ and $i\delta Q$ one obtains the first order KKT
conditions
\begin{subequations}
\begin{align}
\left(  Q\left\vert \left(  \widehat{K+V_{2}}\right)  -\left\langle
\lambda\left\vert \left\{  \widehat{\Phi^{\dagger}\Phi}-\gamma\right\}
\right.  \right\rangle \right.  \delta Q\right)   &  =0\quad\forall\delta
Q\in\mathcal{B}_{2}\left(  \mathcal{H}^{N}\right) \\
\left\langle \delta\lambda\left\vert \left\{  \left(  Q\left\vert
\widehat{\Phi^{\dagger}\Phi}Q\right.  \right)  -\gamma\right\}  \right.
\right\rangle  &  =0\quad\forall\delta\lambda\in L_{\infty}\left(
\mathbb{R}^{3}\times\mathbb{C}\right)
\end{align}%
\end{subequations}
\[
\Updownarrow
\]
%

\begin{subequations}
\begin{align}
\mathcal{P}_{\mathcal{B}_{2}\left(  \mathcal{H}^{N}\right)  }\left\{  \left(
\widehat{K+V_{2}}\right)  -\left\langle \lambda\left\vert \left\{
\widehat{\Phi^{\dagger}\Phi}-\gamma\right\}  \right.  \right\rangle \right\}
\left\vert \overset{}{Q}\right)   &  =0\\
\left(  Q\left\vert \widehat{\Phi^{\dagger}\Phi}Q\right.  \right)  -\gamma &
=0
\end{align}
The operator $Q+\delta Q$ belongs to the level set if
\end{subequations}
\begin{align}
&  \left(  Q+\delta Q\left\vert \widehat{\Phi^{\dagger}\Phi}\right.  \left(
Q+\delta Q\right)  \right)  -\gamma=0\\
&  \Leftrightarrow\left(  Q\left\vert \widehat{\Phi^{\dagger}\Phi}\right.
\delta Q\right)  +\left(  \delta Q\left\vert \widehat{\Phi^{\dagger}\Phi
}\right.  Q\right)  +\left(  \delta Q\left\vert \widehat{\Phi^{\dagger}\Phi
}\right.  \delta Q\right)  =0
\end{align}
which is satisfied if $\delta Q\in\ker\widehat{\Phi^{\dagger}\Phi\left(
\mathbf{x}\right)  }\quad\forall\mathbf{x}$.

We know that for any $\gamma$ there exists at least one IPS\ operator
$\left\vert \Psi\left(  \gamma\right)  \right\rangle \left\langle \Psi\left(
\gamma\right)  \right\vert $ $\in\mathcal{B}_{2}\left(  \mathcal{H}%
^{N}\right)  $ and at least one $Q_{0}\left(  \gamma\right)  \in
\mathcal{B}_{2}\left(  \mathcal{H}^{N}\right)  $ such that
\begin{align}
\left\vert \Psi\left(  \gamma\right)  \right\rangle \left\langle \Psi\left(
\gamma\right)  \right\vert  &  =Q_{0}\left(  \gamma\right)  Q_{0}\left(
\gamma\right)  ^{\dagger}\nonumber\\
\Xi\left(  Q_{0}\left(  \gamma\right)  Q_{0}\left(  \gamma\right)  ^{\dagger
}\right)   &  =\gamma,
\end{align}
showing that $\Xi$ is surjective.

\subsection{Real Variables}

Any positive element $D$ of $\mathcal{B}_{1}\left(  \mathcal{H}^{N}\right)  $
can be expressed as
\begin{equation}
D=a^{2},
\end{equation}
where
\begin{align}
a  &  \in\mathcal{B}_{2}\left(  \mathcal{H}^{N}\right) \nonumber\\
a  &  =a^{\dagger}\nonumber\\
\dim\mathcal{B}_{2}\left(  \mathcal{H}^{N}\right)   &  =r_{N}^{2}.
\end{align}
One can even choose $a$ to be positive but that will be too restrictive for
our purposes as the set of positive Hilbert Schmidt operators form a cone but
do not form a linear space. Thus we can consider \emph{real} Hilbert space
$\mathcal{B}_{2sa}\left(  \mathcal{H}^{N}\right)  $ of the same dimension as
$\mathcal{B}_{2}\left(  \mathcal{H}^{N}\right)  ,$ which has an orthonormal
basis
\begin{align}
\Lambda_{jk}  &  =\frac{1}{\sqrt{2}}\left\{  \left\vert \Psi_{j}\right\rangle
\left\langle \Psi_{k}\right\vert +\left\vert \Psi_{k}\right\rangle
\left\langle \Psi_{j}\right\vert \right\}  ;\quad1\leq j<k\leq r_{N}%
\nonumber\\
\Lambda_{jj}  &  =\left\vert \Psi_{j}\right\rangle \left\langle \Psi
_{j}\right\vert ;\quad1\leq j\leq r_{N}\nonumber\\
\Lambda_{jk}  &  =\frac{i}{\sqrt{2}}\left\{  \left\vert \Psi_{j}\right\rangle
\left\langle \Psi_{k}\right\vert -\left\vert \Psi_{k}\right\rangle
\left\langle \Psi_{j}\right\vert \right\}  ;\quad1\leq j>k\leq r_{N}.
\end{align}
The Lagrangian can then be formulated as
\begin{equation}
\mathcal{L}\left(  a,\lambda,\gamma\right)  =E_{U}\left(  a\right)
-\left\langle \lambda\left\vert \left\{  \Xi\left(  a\right)  -\gamma\right\}
\right.  \right\rangle ,
\end{equation}
where the universal energy, $E_{U}\left(  a\right)  ,$ of a normalized state
$a$ is
\begin{equation}
E_{U}\left(  a\right)  =\left(  a\left\vert \left(  \widehat{K+V_{2}}\right)
\right.  a\right)  =\operatorname{Tr}\left\{  a\left(  K+V_{2}\right)
a\right\}
\end{equation}
and
\begin{equation}
\Xi\left(  a\right)  \left(  \mathbf{x}\right)  =\left(  a\left\vert
\widehat{\Phi\left(  \mathbf{x}\right)  ^{\dagger}\Phi\left(  \mathbf{x}%
\right)  }\right.  a\right)  =\operatorname{Tr}\left\{  a\Phi\left(
\mathbf{x}\right)  ^{\dagger}\Phi\left(  \mathbf{x}\right)  a\right\}  .
\end{equation}
We have not included the normalization term even as it is implied by the
normalization of the density constraint i.e.
\begin{equation}
\int\left(  a\left\vert \widehat{\Phi\left(  \mathbf{x}\right)  ^{\dagger}%
\Phi\left(  \mathbf{x}\right)  }\right.  a\right)  d\mathbf{x}=N\left(
a\left\vert \overset{}{a}\right.  \right)
\end{equation}
and if we choose $\gamma\left(  \mathbf{x}\right)  \geq0$ such that
\begin{equation}
\int\gamma\left(  \mathbf{x}\right)  d\mathbf{x}=N
\end{equation}
then as
\begin{equation}
N\left(  a\left\vert \overset{}{a}\right.  \right)  =\int\gamma\left(
\mathbf{x}\right)  d\mathbf{x}%
\end{equation}
one has
\begin{equation}
\left(  a\left\vert \overset{}{a}\right.  \right)  =1.
\end{equation}


The function $\Xi$ is neither surjective nor injective as the non linear
constraint is
\begin{equation}
\Xi:\mathcal{B}_{2sa}\left(  \mathcal{H}^{N}\right)  \rightarrow L_{1+}\left(
\mathbb{X}\right)  \subset L_{1}\left(  \mathbb{X}\right)  ,
\end{equation}
where $L_{1+}\left(  \mathbb{X}\right)  $ is the positive cone of
$L_{1}\left(  \mathbb{X}\right)  .$ However its differential
\begin{equation}
D\Xi\left(  a\right)  :\mathcal{B}_{2sa}\left(  \mathcal{H}^{N}\right)
\rightarrow L_{1}\left(  \mathbb{X}\right)
\end{equation}
may be surjective. Its coordinates are
\begin{equation}
\left(  \overset{}{\nabla\Xi\left(  \mathbf{x}\right)  \left(  a\right)
}\right)  _{jk}=D\Xi\left(  a\right)  \left(  \mathbf{x}\right)  \left(
\Lambda_{jk}\right)  =\left(  a\left\vert \widehat{\Phi\left(  \mathbf{x}%
\right)  ^{\dagger}\Phi\left(  \mathbf{x}\right)  }\right.  \Lambda
_{jk}\right)  .
\end{equation}
If the vectors $\left\{  \left\vert \nabla\Xi\left(  \mathbf{x}\right)
\left(  a\right)  \right)  \right\}  $ form a basis (complete or overcomplete)
for $L_{1}\left(  \mathbb{X}\right)  $ then $D\Xi\left(  a\right)  $ is
surjective, i.e. for any for $f\in L_{1}\left(  \mathbb{X}\right)  $ there
exists a $b\in$ $\mathcal{B}_{2sa}\left(  \mathcal{H}^{N}\right)  $ such that
\begin{equation}
f\left(  \mathbf{x}\right)  =\left(  a\left\vert \widehat{\Phi\left(
\mathbf{x}\right)  ^{\dagger}\Phi\left(  \mathbf{x}\right)  }\right.
b\right)  \mathbf{.}%
\end{equation}
$\lceil$An overcomplete continuous set $\left\{  \left\vert y\right\}
\right\}  $ is a proper basis if any element $g$ can be expanded
\begin{equation}
\left\vert g\right\rangle =\int g\left(  y\right)  \left\vert y\right\}  d\mu,
\end{equation}
where
\begin{equation}
\left\vert g\right\rangle \neq\left\vert g^{\prime}\right\rangle
\Leftrightarrow g\left(  y\right)  \neq g^{\prime}\left(  y\right)  \text{
almost everywhere with respect to }\mu.\rfloor
\end{equation}
$\lceil$ let $a$ be pure state, i.e. a projector, that is localized in a
region of space so that
\begin{equation}
\widehat{\Phi\left(  \mathbf{x}\right)  ^{\dagger}\Phi\left(  \mathbf{x}%
\right)  }\left\vert \overset{}{a}\right)  =0\,\forall\mathbf{x\in\Delta}%
\end{equation}
then $\nexists$ $b$ $\backepsilon f\left(  \mathbf{x}\right)  \neq0$ when
$\mathbf{x\in\Delta,}$ hence $D\Xi\left(  a\right)  $ is not
surjective.$\mathbf{\rfloor}$

The grammian matrix for this basis has matrix elements
\begin{align}
\left(  \overset{}{\nabla\Xi\left(  \mathbf{x}\right)  \left(  a\right)
}\left\vert \overset{}{\nabla\Xi\left(  \mathbf{x}^{\prime}\right)  \left(
a\right)  }\right.  \right)   &  =\left(  a\left\vert \widehat{\Phi\left(
\mathbf{x}\right)  ^{\dagger}\Phi\left(  \mathbf{x}\right)  }\widehat
{\Phi\left(  \mathbf{x}^{\prime}\right)  ^{\dagger}\Phi\left(  \mathbf{x}%
^{\prime}\right)  }\right.  a\right) \nonumber\\
&  =\left\{
\begin{array}
[c]{c}%
\left(  a\left\vert \widehat{\Phi\left(  \mathbf{x}\right)  ^{\dagger}%
\Phi\left(  \mathbf{x}\right)  }\right.  a\right)  =\gamma\left(
\mathbf{x}\right)  ,\text{ for }\mathbf{x=x}^{\prime}\\
\left(  a\left\vert \widehat{\Phi\left(  \mathbf{x}\right)  ^{\dagger}%
\Phi\left(  \mathbf{x}^{\prime}\right)  ^{\dagger}\Phi\left(  \mathbf{x}%
^{\prime}\right)  \Phi\left(  \mathbf{x}\right)  }\right.  a\right)
,\text{for }\mathbf{x\neq x}^{\prime}%
\end{array}
.\right.
\end{align}
The vectors $\left\{  \left.  \overset{}{\nabla\Xi\left(  \mathbf{x}\right)
\left(  a\right)  }\right\vert \mathbf{x\in}\mathbb{X}\right\}  $ form a basis
for the normal space $\mathcal{N}_{a}$ on the surface $\Xi\left(  a\right)
=\gamma$ at the point $a.$ If these vectors are linearly dependent then there
exists a $g$ such that
\begin{equation}
\int g\left(  \mathbf{x}\right)  \widehat{\Phi\left(  \mathbf{x}\right)
^{\dagger}\Phi\left(  \mathbf{x}\right)  }\left\vert \overset{}{a}\right)
=\widehat{\Omega\left(  g\right)  }\left\vert \overset{}{a}\right)  =0,
\label{roeig}%
\end{equation}
which in particular implies that
\begin{equation}
\left(  a\left\vert \widehat{\Omega\left(  g\right)  }\right.  a\right)
=\left\langle \left.  g\right\vert \gamma\right\rangle =0.
\end{equation}
(Of course $\widehat{\Omega\left(  g\right)  }\left\vert \overset{}{a}\right)
\neq0$ does not imply that $\left\{  \left.  \overset{}{\nabla\Xi\left(
\mathbf{x}\right)  \left(  a\right)  }\right\vert \mathbf{x\in}\mathbb{X}%
\right\}  $ are linearly independent). The eq. $\left(  \ref{zeroeig}\right)
$ could be satisfied quite often by local potentials $\Omega\left(  g\right)
$ that are localized in a region of space where the state has no density.

The operators $\left\{  \Phi\left(  \mathbf{x}\right)  ^{\dagger}\Phi\left(
\mathbf{x}\right)  \right\}  $ act on the spaces $\left\{  \mathcal{H}%
^{N}\right\}  $ as
\begin{align}
P_{N}\Phi\left(  \mathbf{x}\right)  ^{\dagger}\Phi\left(  \mathbf{x}\right)
P_{N}  &  =\int\left\vert \mathbf{x}_{1}\cdots\mathbf{x}_{N}\right\rangle
\left\langle \mathbf{x}_{1}\cdots\mathbf{x}_{N}\right\vert \Phi\left(
\mathbf{x}\right)  ^{\dagger}\Phi\left(  \mathbf{x}\right)  \left\vert
\mathbf{x}_{1}^{\prime}\cdots\mathbf{x}_{N}^{\prime}\right\rangle \left\langle
\mathbf{x}_{1}^{\prime}\cdots\mathbf{x}_{N}^{\prime}\right\vert d\mathbf{x}%
_{1}\cdots d\mathbf{x}_{N}d\mathbf{x}_{1}^{\prime}\cdots d\mathbf{x}%
_{N}^{\prime}\nonumber\\
&  =\int\left\vert \mathbf{xx}_{2}\cdots\mathbf{x}_{N}\right\rangle
\left\langle \mathbf{xx}_{2}\cdots\mathbf{x}_{N}\right\vert d\mathbf{x}%
_{2}\cdots d\mathbf{x}_{N}%
\end{align}
so
\begin{equation}
\int P_{N}\Phi\left(  \mathbf{x}\right)  ^{\dagger}\Phi\left(  \mathbf{x}%
\right)  P_{N}d\mathbf{x}=NI_{\mathcal{H}^{N}}%
\end{equation}
and thus form a resolution of the identity
\begin{equation}
\frac{1}{N}\int P_{N}\Phi\left(  \mathbf{x}\right)  ^{\dagger}\Phi\left(
\mathbf{x}\right)  P_{N}d\mathbf{x}=I_{\mathcal{H}^{N}}.
\end{equation}
The operators $\left\{  \frac{1}{N}P_{N}\Phi\left(  \mathbf{x}\right)
^{\dagger}\Phi\left(  \mathbf{x}\right)  P_{N}\right\}  $ are projection
operators (not one dimensional unless $N=1$)
\begin{align}
\left[  P_{N}\Phi\left(  \mathbf{x}\right)  ^{\dagger}\Phi\left(
\mathbf{x}\right)  P_{N}\right]  ^{2}  &  =\frac{1}{N^{2}}\int\left\vert
\mathbf{xx}_{2}\cdots\mathbf{x}_{N}\right\rangle \left\langle \mathbf{xx}%
_{2}\cdots\mathbf{x}_{N}\right\vert d\mathbf{x}_{2}\cdots d\mathbf{x}_{N}%
\int\left\vert \mathbf{xx}_{2}\cdots\mathbf{x}_{N}\right\rangle \left\langle
\mathbf{xx}_{2}\cdots\mathbf{x}_{N}\right\vert d\mathbf{x}_{2}\cdots
d\mathbf{x}_{N}\nonumber\\
&  =\frac{1}{N}\int\left\vert \mathbf{xx}_{2}\cdots\mathbf{x}_{N}\right\rangle
\left\langle \mathbf{xx}_{2}\cdots\mathbf{x}_{N}\right\vert d\mathbf{x}%
_{2}\cdots d\mathbf{x}_{N}=\frac{1}{N}P_{N}\Phi\left(  \mathbf{x}\right)
^{\dagger}\Phi\left(  \mathbf{x}\right)  P_{N},
\end{align}
but not in not orthogonal
\begin{align}
&  \frac{1}{N}P_{N}\Phi\left(  \mathbf{x}\right)  ^{\dagger}\Phi\left(
\mathbf{x}\right)  P_{N}P_{N}\frac{1}{N}\Phi\left(  \mathbf{x}^{\prime
}\right)  ^{\dagger}\Phi\left(  \mathbf{x}^{\prime}\right)  P_{N}=\nonumber\\
&  \frac{1}{N^{2}}\int\left\vert \mathbf{xx}_{2}\cdots\mathbf{x}%
_{N}\right\rangle \left\langle \mathbf{xx}_{2}\cdots\mathbf{x}_{N}\right\vert
d\mathbf{x}_{2}\cdots d\mathbf{x}_{N}\int\left\vert \mathbf{x}^{\prime
}\mathbf{x}_{2}\cdots\mathbf{x}_{N}\right\rangle \left\langle \mathbf{x}%
^{\prime}\mathbf{x}_{2}\cdots\mathbf{x}_{N}\right\vert d\mathbf{x}_{2}\cdots
d\mathbf{x}_{N}\nonumber\\
&  =\frac{N-1}{N}\int\left\vert \mathbf{x\mathbf{x}^{\prime}x}_{3}%
\cdots\mathbf{x}_{N}\right\rangle \left\langle \mathbf{x\mathbf{x}^{\prime}%
x}_{3}\cdots\mathbf{x}_{N}\right\vert d\mathbf{x}_{3}\cdots d\mathbf{x}%
_{N}=\frac{1}{N\left(  N-1\right)  }P_{N}\Phi\left(  \mathbf{x}\right)
^{\dagger}\Phi\left(  \mathbf{x}^{\prime}\right)  ^{\dagger}\Phi\left(
\mathbf{x}^{\prime}\right)  \Phi\left(  \mathbf{x}\right)  P_{N}%
\end{align}
Their products are however projection operators as
\begin{align}
&  \Phi\left(  \mathbf{x}\right)  ^{\dagger}\Phi\left(  \mathbf{x}\right)
\Phi\left(  \mathbf{x}^{\prime}\right)  ^{\dagger}\Phi\left(  \mathbf{x}%
^{\prime}\right)  \Phi\left(  \mathbf{x}\right)  ^{\dagger}\Phi\left(
\mathbf{x}\right)  \Phi\left(  \mathbf{x}^{\prime}\right)  ^{\dagger}%
\Phi\left(  \mathbf{x}^{\prime}\right)  =\nonumber\\
&  \Phi\left(  \mathbf{x}\right)  ^{\dagger}\Phi\left(  \mathbf{x}\right)
\Phi\left(  \mathbf{x}^{\prime}\right)  ^{\dagger}\Phi\left(  \mathbf{x}%
^{\prime}\right)  \Phi\left(  \mathbf{x}^{\prime}\right)  ^{\dagger}%
\Phi\left(  \mathbf{x}^{\prime}\right)  \Phi\left(  \mathbf{x}\right)
^{\dagger}\Phi\left(  \mathbf{x}\right)  =\nonumber\\
&  \Phi\left(  \mathbf{x}\right)  ^{\dagger}\Phi\left(  \mathbf{x}\right)
\Phi\left(  \mathbf{x}^{\prime}\right)  ^{\dagger}\Phi\left(  \mathbf{x}%
^{\prime}\right)  \Phi\left(  \mathbf{x}\right)  ^{\dagger}\Phi\left(
\mathbf{x}\right)  =\nonumber\\
&  \Phi\left(  \mathbf{x}\right)  ^{\dagger}\Phi\left(  \mathbf{x}\right)
\Phi\left(  \mathbf{x}\right)  ^{\dagger}\Phi\left(  \mathbf{x}\right)
\Phi\left(  \mathbf{x}^{\prime}\right)  ^{\dagger}\Phi\left(  \mathbf{x}%
^{\prime}\right)  =\Phi\left(  \mathbf{x}\right)  ^{\dagger}\Phi\left(
\mathbf{x}\right)  \Phi\left(  \mathbf{x}^{\prime}\right)  ^{\dagger}%
\Phi\left(  \mathbf{x}^{\prime}\right)  .
\end{align}
They do not form spectral expansions as one can observe
\begin{align}
&  \int f\left(  \mathbf{x}\right)  \Phi\left(  \mathbf{x}\right)  ^{\dagger
}\Phi\left(  \mathbf{x}\right)  d\mathbf{x}\int f\left(  \mathbf{x}\right)
^{-1}\Phi\left(  \mathbf{x}\right)  ^{\dagger}\Phi\left(  \mathbf{x}\right)
d\mathbf{x=}\int f\left(  \mathbf{x}\right)  f\left(  \mathbf{x}^{\prime
}\right)  ^{-1}\Phi\left(  \mathbf{x}\right)  ^{\dagger}\Phi\left(
\mathbf{x}\right)  \Phi\left(  \mathbf{x}^{\prime}\right)  ^{\dagger}%
\Phi\left(  \mathbf{x}^{\prime}\right)  d\mathbf{x}d\mathbf{x}^{\prime
}=\nonumber\\
&  \int\Phi\left(  \mathbf{x}\right)  ^{\dagger}\Phi\left(  \mathbf{x}\right)
d\mathbf{x+}\int f\left(  \mathbf{x}\right)  f\left(  \mathbf{x}^{\prime
}\right)  ^{-1}\Phi\left(  \mathbf{x}\right)  ^{\dagger}\Phi\left(
\mathbf{x}^{\prime}\right)  ^{\dagger}\Phi\left(  \mathbf{x}^{\prime}\right)
\Phi\left(  \mathbf{x}\right)  d\mathbf{x}d\mathbf{x}^{\prime},
\end{align}
so in general
\begin{equation}
\int f\left(  \mathbf{x}\right)  ^{-1}\Phi\left(  \mathbf{x}\right)
^{\dagger}\Phi\left(  \mathbf{x}\right)  d\mathbf{x\neq}\left[  \int f\left(
\mathbf{x}\right)  \Phi\left(  \mathbf{x}\right)  ^{\dagger}\Phi\left(
\mathbf{x}\right)  d\mathbf{x}\right]  ^{-1}.
\end{equation}


The transpose of the map $D\Xi\left(  a\right)  $ is given by
\begin{align}
D\Xi\left(  a\right)  ^{t}  &  :L_{\infty}\left(  \mathbb{X}\right)
\rightarrow\mathcal{B}_{2sa}\left(  \mathcal{H}^{N}\right) \nonumber\\
\left(  \overset{}{D\Xi\left(  a\right)  ^{t}f}\left\vert \overset{}%
{b}\right.  \right)  _{\mathcal{B}_{2sa}\left(  \mathcal{H}^{N}\right)  }  &
=\left(  \overset{}{f}\left\vert \overset{}{D\Xi\left(  a\right)  b}\right.
\right)  _{L_{\infty}\left(  \mathbb{X}\right)  \times L_{1}\left(
\mathbb{X}\right)  }\text{ for }f\in L_{\infty}\left(  \mathbb{X}\right)
\text{ and all }b\in\mathcal{B}_{2sa}\left(  \mathcal{H}^{N}\right)
\end{align}
leading to
\begin{equation}
\ker\left\{  D\Xi\left(  a\right)  \right\}  =\operatorname{Ran}\left\{
D\Xi\left(  a\right)  ^{t}\right\}  ^{\bot}.
\end{equation}
The generalized inverse of the $D\Xi\left(  a\right)  $ is given by
\begin{equation}
D\Xi\left(  a\right)  ^{G}:L_{1}\left(  \mathbb{X}\right)  \rightarrow
\mathcal{B}_{2sa}\left(  \mathcal{H}^{N}\right)  ,
\end{equation}
where
\begin{align}
D\Xi\left(  a\right)  ^{G}  &  =\sum_{1\leq j\leq r^{\prime}}\eta_{j}%
^{-1}\left\vert X_{j}\right)  \left(  \xi_{j}^{d}\right\vert \\
D\Xi\left(  a\right)   &  =\sum_{1\leq j\leq r^{\prime}}\eta_{j}\left\vert
\xi_{j}\right)  \left(  X_{j}\right\vert ,
\end{align}
$\left\{  \left.  \left\vert X_{j}\right)  \right\vert 1\leq j\leq r_{N}%
^{2}\right\}  ,\left\{  \left.  \left(  \xi_{j}^{d}\right\vert \right\vert
1\leq j\leq r\right\}  $ and $\left\{  \left.  \left\vert \xi_{j}\right)
\right\vert 1\leq j\leq r\right\}  $ are dual bases for $\mathcal{B}%
_{2sa}\left(  \mathcal{H}^{N}\right)  ,$ $L_{\infty}\left(  \mathbb{X}\right)
$ and $L_{1}\left(  \mathbb{X}\right)  $ respectively and
\begin{equation}
\ker\left\{  D\Xi\left(  a\right)  \right\}  =\operatorname{linspan}\left\{
\left.  \left\vert \overset{}{X_{j}}\right)  \right\vert r^{\prime}+1\leq
j\leq r_{N}^{2}\right\}  .
\end{equation}


$\lceil$A simple example of a constraining function is
\begin{equation}
F\left(  x,y,z\right)  =2x^{2}+3y^{2}+4z^{2}.
\end{equation}
The level surfaces are
\begin{equation}
F\left(  x,y,z\right)  =c^{2}\equiv G\left(  x,y,z\right)  =0,
\end{equation}
where
\begin{equation}
G\left(  x,y,z\right)  =2x^{2}+3y^{2}+4z^{2}-c^{2}%
\end{equation}
and produce a family of ellipsoids in $\mathbb{R}^{3}.$ The equations of the
ellipsoids are
\begin{equation}
z^{2}=\frac{1}{4}\left(  c^{2}-2x^{2}-3y^{2}\right)
\end{equation}
giving
\begin{equation}
z_{\pm}\left(  x,y,c\right)  =\pm\frac{1}{4}\sqrt{\left(  c^{2}-2x^{2}%
-3y^{2}\right)  },
\end{equation}
showing that for fixed $c,$ $z_{\pm}\left(  x,y,c\right)  $ are
parameterizations of portions of the surface in terms of $\left(  x,y\right)
$. Hence, for instance, the points $\left(  1,1,\pm\frac{1}{4}\sqrt{\left(
c^{2}-5\right)  }\right)  $ are on the surface. As $F\left(  x,y,z\right)  $
and $G\left(  x,y,z\right)  $ are level surfaces
\begin{align}
DF\left(  x,y,z\right)  \left(  \delta x,\delta y,\delta z\right)   &
=DG\left(  x,y,z\right)  \left(  \delta x,\delta y,\delta z\right) \nonumber\\
&  =\frac{\partial F\left(  x,y,z\right)  }{\partial x}\delta x+\frac{\partial
F\left(  x,y,z\right)  }{\partial y}\delta y+\frac{\partial F\left(
x,y,z\right)  }{\partial z}\delta z=0\nonumber\\
&  =\nabla F\left(  x,y,z\right)  ^{t}\cdot\left(
\begin{array}
[c]{c}%
\delta x\\
\delta y\\
\delta z
\end{array}
\right)  =\left(  \left.  \nabla F\left(  x,y,z\right)  \right\vert
\delta\mathbf{r}\right)  =0
\end{align}
A tangent vector $\mathbf{v}$ at the point $\left(  x,y,z\right)  $ on the
surface $F\left(  x,y,z\right)  =c^{2}$ is defined by curves $\mathbf{r}%
\left(  t\right)  $ that lie on the surface for $\left\vert t\right\vert
\leq\delta,$ by considering the derivative of $F\left(  \mathbf{r}\left(
t\right)  \right)  $ at the point $t=0,$ where $\mathbf{r}\left(  0\right)
=\left(  x,y,z\right)  $ and $\mathbf{\dot{r}}\left(  0\right)  =\mathbf{v,}$
i.e.
\begin{align}
&  \left.  \frac{dF\left(  \mathbf{r}\left(  t\right)  \right)  }%
{dt}\right\vert _{t=0}=\frac{\partial F\left(  x,y,z\right)  }{\partial
x}\frac{dx\left(  0\right)  }{dt}+\frac{\partial F\left(  x,y,z\right)
}{\partial y}\frac{dy\left(  0\right)  }{dt}+\frac{\partial F\left(
x,y,z\right)  }{\partial z}\frac{dz\left(  0\right)  }{dt}\nonumber\\
&  =\nabla F\left(  x,y,z\right)  ^{t}\left(
\begin{array}
[c]{c}%
\dot{x}\left(  0\right) \\
\dot{y}\left(  0\right) \\
\dot{z}\left(  0\right)
\end{array}
\right)  =\nabla F\left(  x,y,z\right)  ^{t}\left(
\begin{array}
[c]{c}%
v_{x}\\
v_{y}\\
v_{z}%
\end{array}
\right)  =0=DF\left(  x,y,z\right)  \left(  \mathbf{v}\right)  .
\end{align}
Thus a necessary condition to be a tangent vector is that $\mathbf{v\in}%
\ker\left\{  DF\left(  x,y,z\right)  \right\}  ,$ where
\begin{equation}
\ker\left\{  DF\left(  x,y,z\right)  \right\}  =\operatorname{linspan}\left\{
\left.  \overset{}{\mathbf{h}}\right\vert DF\left(  x,y,z\right)  \left(
\mathbf{h}\right)  =\left(  \left.  \nabla F\left(  x,y,z\right)  \right\vert
\mathbf{h}\right)  =0\right\}  .
\end{equation}
It is not a sufficient condition as there may be elements of $\mathbf{h\in
}\ker\left\{  DF\left(  x,y,z\right)  \right\}  $ such that
\begin{equation}
\mathbf{h\neq\dot{r}}\left(  0\right)  ,
\end{equation}
where $\mathbf{r}\left(  t\right)  $ is curve on the surface. If $F$ is
continuously differentiable at $\mathbf{r}\left(  0\right)  $ then \emph{all
}directional derivatives exist at $\mathbf{r}\left(  0\right)  $ and are given
by
\begin{equation}
\nabla F\left(  \mathbf{r}\left(  0\right)  \right)  ^{t}\cdot\mathbf{v}%
=\nabla_{\mathbf{v}}F\left(  \mathbf{r}\left(  0\right)  \right)  ,
\end{equation}
some of these directions are not along the surface i.e. not in the tangent
space even if $\nabla_{\mathbf{v}}F\left(  \mathbf{r}\left(  0\right)
\right)  =0$.

The tangent space $\mathcal{T}_{\mathbf{r}\left(  0\right)  }$ at any point
$\mathbf{r}\left(  0\right)  $ is given by
\begin{equation}
\mathcal{T}_{\mathbf{r}\left(  0\right)  }=\operatorname{linspan}\left\{
\left.  \mathbf{v}\right\vert \left(  \left.  \nabla F\left(  x,y,z\right)
\right\vert \mathbf{v}\right)  =0;\mathbf{v=\dot{r}}\left(  0\right)
;F\left(  \mathbf{r}\left(  t\right)  \right)  =c^{2}\mathbf{,\quad}\left\vert
t\right\vert \leq\delta\right\}  .
\end{equation}
The normal $\mathcal{N}_{\mathbf{r}\left(  0\right)  }$ space is given by
\begin{equation}
\mathcal{N}_{\mathbf{r}\left(  0\right)  }=\operatorname{linspan}\left\{
\nabla F\left(  x,y,z\right)  \right\}  .
\end{equation}
In this case the normal space is generated by the gradient $\nabla F\left(
x,y,z\right)  $ that is
\begin{equation}
\nabla F\left(  x,y,z\right)  =\nabla G\left(  x,y,z\right)  =\left(
2x,3y,4z\right)  .
\end{equation}
Clearly $\nabla F\left(  x,y,z\right)  \neq\mathbf{0}$ except at the origin,
but then that point does not lie on the level surface unless $c=0$, where the
level surface is just the single point $\left(  0,0,0\right)  .$

We can consider a more general example
\begin{equation}
F\left(  x,y,z\right)  =\left(
\begin{array}
[c]{c}%
F_{1}\left(  x,y,z\right) \\
F_{2}\left(  x,y,z\right)
\end{array}
\right)  =\left(
\begin{array}
[c]{c}%
c_{1}^{2}\\
c_{2}^{2}%
\end{array}
\right)  \equiv G\left(  x,y,z\right)  =\mathbf{0},
\end{equation}
where
\begin{equation}
G\left(  x,y,z\right)  =\left(
\begin{array}
[c]{c}%
2x^{2}+3y^{2}\\
3y^{2}+4z^{2}%
\end{array}
\right)  -\left(
\begin{array}
[c]{c}%
c_{1}^{2}\\
c_{2}^{2}%
\end{array}
\right)  =0.
\end{equation}
This has a parameterization
\begin{align}
x\left(  y\right)   &  =\pm\sqrt{\frac{1}{2}\left(  c_{1}^{2}-3y^{2}\right)
}\nonumber\\
z\left(  y\right)   &  =\pm\sqrt{\frac{1}{3}\left(  c_{2}^{2}-3y^{2}\right)
},
\end{align}
showing that the manifold is one dimensional. As $F\left(  x,y,z\right)  $ and
$G\left(  x,y,z\right)  $ are level surfaces
\begin{align}
DF\left(  x,y,z\right)  \left(  \delta x,\delta y,\delta z\right)   &
=DG\left(  x,y,z\right)  \left(  \delta x,\delta y,\delta z\right) \nonumber\\
&  =\left(
\begin{array}
[c]{ccc}%
\frac{\partial F_{1}\left(  x,y,z\right)  }{\partial x} & \frac{\partial
F_{1}\left(  x,y,z\right)  }{\partial y} & \frac{\partial F_{1}\left(
x,y,z\right)  }{\partial z}\\
\frac{\partial F_{2}\left(  x,y,z\right)  }{\partial x} & \frac{\partial
F_{2}\left(  x,y,z\right)  }{\partial y} & \frac{\partial F_{2}\left(
x,y,z\right)  }{\partial z}%
\end{array}
\right)  \left(
\begin{array}
[c]{c}%
\delta x\\
\delta y\\
\delta z
\end{array}
\right)  =\left(
\begin{array}
[c]{c}%
0\\
0
\end{array}
\right)  .
\end{align}
The tangent space $\mathcal{T}_{\mathbf{r}\left(  0\right)  }$ at any point
$\mathbf{r}\left(  0\right)  $ belonging to the surface is given by
\begin{equation}
\mathcal{T}_{\mathbf{r}\left(  0\right)  }=\operatorname{linspan}\left\{
\left.  \overset{}{\mathbf{v}}\right\vert DF\left(  \mathbf{r}\left(
0\right)  \right)  \left\vert \mathbf{v}\right)  =0;\mathbf{v=\dot{r}}\left(
0\right)  ;F\left(  \mathbf{r}\left(  t\right)  \right)  =c^{2}\mathbf{,\quad
}\left\vert t\right\vert \leq\delta\right\}  \subseteq.\ker\left\{  DF\left(
\mathbf{r}\left(  0\right)  \right)  \right\}  .
\end{equation}
If the dimension of $\,\ker\left\{  DF\left(  x,y,z\right)  \right\}  $ is $1$
then $DF\left(  x,y,z\right)  $ is surjective. \emph{A point where }$DF\left(
x,y,z\right)  $\emph{ is surjective is a called a normal point and all
solutions }$\mathbf{v}$ to $DF\left(  \mathbf{r}\left(  0\right)  \right)
\left\vert \mathbf{v}\right)  =0$ are tangent vectors. \emph{ }If the vectors
$\nabla F_{1}\left(  x,y,z\right)  $ and $\nabla F_{2}\left(  x,y,z\right)  $
are linearly dependent, but not both zero then the dimension is $2.$ \emph{If
all solutions }$\mathbf{v}$ to $DF\left(  \mathbf{r}\left(  0\right)  \right)
\left\vert \mathbf{v}\right)  =0$ are still tangent vectors then such a point
is called a \emph{regular point}.\emph{. }If they are both zero (a critical
point) then the dimension is $3,$then all vectors are solutions and some are
not tangent vectors, (so the perpendicular complement to the normal space is
not the tangent space in this case)$.$ The normal $\mathcal{N}_{\mathbf{r}%
\left(  0\right)  }$ space is
\begin{equation}
\mathcal{N}_{\mathbf{r}\left(  0\right)  }=\operatorname{linspan}\left\{
\nabla F_{1}\left(  \mathbf{r}\left(  0\right)  \right)  ,\nabla F_{2}\left(
\mathbf{r}\left(  0\right)  \right)  \right\}  ,
\end{equation}
where
\begin{equation}
\nabla F_{j}\left(  x,y,z\right)  =\left(
\begin{array}
[c]{ccc}%
\frac{\partial F_{j}\left(  x,y,z\right)  }{\partial x} & \frac{\partial
F_{j}\left(  x,y,z\right)  }{\partial y} & \frac{\partial F_{j}\left(
x,y,z\right)  }{\partial z}%
\end{array}
\right)  .
\end{equation}
Hence
\begin{align}
\dim\left\{  \mathcal{N}_{\mathbf{r}\left(  0\right)  }\right\}   &
=2\Leftrightarrow\dim\left\{  \ker\left\{  DF\left(  \mathbf{r}\left(
0\right)  \right)  \right\}  \right\}  =1\Leftrightarrow\dim\left\{
\mathcal{T}_{\mathbf{r}\left(  0\right)  }\right\}  =1\\
\dim\left\{  \mathcal{N}_{\mathbf{r}\left(  0\right)  }\right\}   &
=1\Leftrightarrow\dim\left\{  \ker\left\{  DF\left(  \mathbf{r}\left(
0\right)  \right)  \right\}  \right\}  =2\\
\dim\left\{  \mathcal{N}_{\mathbf{r}\left(  0\right)  }\right\}   &
=0\Leftrightarrow\dim\left\{  \ker\left\{  DF\left(  \mathbf{r}\left(
0\right)  \right)  \right\}  \right\}  =3
\end{align}
In this case
\begin{equation}
\left(
\begin{array}
[c]{ccc}%
\frac{\partial F_{1}\left(  x,y,z\right)  }{\partial x} & \frac{\partial
F_{1}\left(  x,y,z\right)  }{\partial y} & \frac{\partial F_{1}\left(
x,y,z\right)  }{\partial z}\\
\frac{\partial F_{2}\left(  x,y,z\right)  }{\partial x} & \frac{\partial
F_{2}\left(  x,y,z\right)  }{\partial y} & \frac{\partial F_{2}\left(
x,y,z\right)  }{\partial z}%
\end{array}
\right)  =\left(
\begin{array}
[c]{ccc}%
4x & 6y & 0\\
0 & 6y & 8z
\end{array}
\right)
\end{equation}
and $\nabla F_{1}\left(  x,y,z\right)  $ and $\nabla F_{2}\left(
x,y,z\right)  $ are linearly independent and thus the tangent space is one
dimensional if we choose $\left(  c_{1}^{2},c_{2}^{2}\right)  =\left(
8,27\right)  $ then the point $\left(  2,0,3\right)  $ is on the surface. At
this point the normal space is generated by $\left\{  \left(  8,0,0\right)
,\left(  0,0,24\right)  \right\}  $ and the tangent vector $\mathbf{\dot{r}%
}\left(  0\right)  $ must be perpendicular to both these directions, hence
\begin{equation}
\mathbf{\dot{r}}\left(  0\right)  =\left(  8,0,0\right)  \times\left(
0,0,24\right)  =\left\vert
\begin{array}
[c]{ccc}%
\mathbf{i} & \mathbf{j} & \mathbf{k}\\
8 & 0 & 0\\
0 & 0 & 24
\end{array}
\right\vert =-192\mathbf{j\equiv}\left(  0,192,0\right)  .
\end{equation}


If
\begin{equation}
\left(
\begin{array}
[c]{ccc}%
4x & 6y & 0
\end{array}
\right)  =\,\alpha\left(
\begin{array}
[c]{ccc}%
0 & 6y & 8z
\end{array}
\right)  ,
\end{equation}
then these vectors are linearly dependent and the dimension of the normal
space is $1.$ This occurs at $\left(  x,z\right)  =\left(  0,0\right)  $ and
$\alpha=1\ $at the points on the level curve
\begin{equation}
\left(
\begin{array}
[c]{c}%
2x^{2}+3y^{2}\\
3y^{2}+4z^{2}%
\end{array}
\right)  -\left(
\begin{array}
[c]{c}%
c^{2}\\
c^{2}%
\end{array}
\right)  =0,
\end{equation}
where
\begin{equation}
3y^{2}=c^{2}%
\end{equation}
i.e. at the points $\left(  0,\pm\frac{c}{\sqrt{3}},0\right)  .$ Tangent
vectors at this point must satisfy the equations%

\begin{equation}
\left(
\begin{array}
[c]{ccc}%
4x & 6y & 0\\
0 & 6y & 8z
\end{array}
\right)  \left(
\begin{array}
[c]{c}%
v_{1}\\
v_{2}\\
v_{3}%
\end{array}
\right)  =\left(
\begin{array}
[c]{c}%
0\\
0
\end{array}
\right)
\end{equation}
As $x=z=0$ then there are two vectors that satisfy this equation
\begin{equation}
\left(
\begin{array}
[c]{ccc}%
0 & \frac{\pm c}{\sqrt{3}} & 0\\
0 & \frac{\pm c}{\sqrt{3}} & 0
\end{array}
\right)  \left(
\begin{array}
[c]{c}%
v_{1}\\
v_{2}\\
v_{3}%
\end{array}
\right)  =\left(
\begin{array}
[c]{c}%
0\\
0
\end{array}
\right)  ,
\end{equation}
i.e.
\begin{equation}
\left(
\begin{array}
[c]{c}%
v_{1}\\
0\\
0
\end{array}
\right)  ,\left(
\begin{array}
[c]{c}%
0\\
0\\
v_{3}%
\end{array}
\right)  .
\end{equation}


Curves on the surface
\begin{equation}
G\left(  x,y,z\right)  =\left(
\begin{array}
[c]{c}%
2x^{2}+3y^{2}\\
3y^{2}+4z^{2}%
\end{array}
\right)  -\left(
\begin{array}
[c]{c}%
c_{1}^{2}\\
c_{2}^{2}%
\end{array}
\right)  .
\end{equation}
are given by
\begin{align}
&  g\left(  t\right)  =\left(
\begin{array}
[c]{c}%
2x\left(  t\right)  ^{2}+3y\left(  t\right)  ^{2}\\
3y\left(  t\right)  ^{2}+4z\left(  t\right)  ^{2}%
\end{array}
\right)  -\left(
\begin{array}
[c]{c}%
c_{1}^{2}\\
c_{2}^{2}%
\end{array}
\right)  =\left(
\begin{array}
[c]{c}%
\mathbf{r}\left(  t\right)  ^{t}\left(
\begin{array}
[c]{ccc}%
2 & 0 & 0\\
0 & 3 & 0\\
0 & 0 & 0
\end{array}
\right)  \mathbf{r}\left(  t\right) \\
\mathbf{r}\left(  t\right)  ^{t}\left(
\begin{array}
[c]{ccc}%
0 & 0 & 0\\
0 & 3 & 0\\
0 & 0 & 4
\end{array}
\right)  \mathbf{r}\left(  t\right)
\end{array}
\right)  -\left(
\begin{array}
[c]{c}%
c_{1}^{2}\\
c_{2}^{2}%
\end{array}
\right)  =\\
&  \left(  \mathbf{r}\left(  t\right)  ^{t}\mathbf{r}\left(  t\right)
^{t}\right)  \left(
\begin{array}
[c]{cc}%
\left(
\begin{array}
[c]{ccc}%
2 & 0 & 0\\
0 & 3 & 0\\
0 & 0 & 0
\end{array}
\right)  & \mathbb{O}\\
\mathbb{O} & \left(
\begin{array}
[c]{ccc}%
0 & 0 & 0\\
0 & 3 & 0\\
0 & 0 & 4
\end{array}
\right)
\end{array}
\right)  \left(
\begin{array}
[c]{c}%
\mathbf{r}\left(  t\right) \\
\mathbf{r}\left(  t\right)
\end{array}
\right)  -\left(
\begin{array}
[c]{cc}%
c_{1}^{2} & \\
& c_{2}^{2}%
\end{array}
\right)  =0
\end{align}
where
\begin{equation}
\mathbf{r}\left(  t\right)  =\left(
\begin{array}
[c]{c}%
x\left(  t\right) \\
y\left(  t\right) \\
z\left(  t\right)
\end{array}
\right)  .
\end{equation}
One can observe that if $\left\Vert \mathbf{r}\left(  t\right)  \right\Vert
=1$
\begin{equation}
\left(
\begin{array}
[c]{c}%
\mathbf{r}\left(  t\right)  ^{t}\left(
\begin{array}
[c]{ccc}%
2-c_{1}^{2} & 0 & 0\\
0 & 3-c_{1}^{2} & 0\\
0 & 0 & -c_{1}^{2}%
\end{array}
\right)  \mathbf{r}\left(  t\right) \\
\mathbf{r}\left(  t\right)  ^{t}\left(
\begin{array}
[c]{ccc}%
-c_{2}^{2} & 0 & 0\\
0 & 3-c_{2}^{2} & 0\\
0 & 0 & 4-c_{2}^{2}%
\end{array}
\right)  \mathbf{r}\left(  t\right)
\end{array}
\right)  =\left(
\begin{array}
[c]{c}%
0\\
0
\end{array}
\right)  .
\end{equation}
If the constraint matrices $\mathbb{A}_{1}$ and $\mathbb{A}_{2}$ had been
positive then
\begin{equation}
\mathbb{A}_{1}\mathbf{r}\left(  t\right)  =\mathbb{A}_{2}\mathbf{r}\left(
t\right)  =\mathbf{0.}%
\end{equation}
The Taylor series expansion of the curve around the point $t$ is
\begin{align}
g\left(  t+\delta t\right)   &  =g\left(  t\right)  +2\left(  \mathbf{r}%
\left(  t\right)  ^{t}\mathbf{r}\left(  t\right)  ^{t}\right)  \left(
\begin{array}
[c]{cc}%
\left(
\begin{array}
[c]{ccc}%
2 & 0 & 0\\
0 & 3 & 0\\
0 & 0 & 0
\end{array}
\right)  & \mathbb{O}\\
\mathbb{O} & \left(
\begin{array}
[c]{ccc}%
0 & 0 & 0\\
0 & 3 & 0\\
0 & 0 & 4
\end{array}
\right)
\end{array}
\right)  \left(
\begin{array}
[c]{c}%
\mathbf{\dot{r}}\left(  t\right) \\
\mathbf{\dot{r}}\left(  t\right)
\end{array}
\right)  \delta t\\
&  +\left\{  \left(  \mathbf{r}\left(  t\right)  ^{t}\mathbf{r}\left(
t\right)  ^{t}\right)  \left(
\begin{array}
[c]{cc}%
\left(
\begin{array}
[c]{ccc}%
2 & 0 & 0\\
0 & 3 & 0\\
0 & 0 & 0
\end{array}
\right)  & \mathbb{O}\\
\mathbb{O} & \left(
\begin{array}
[c]{ccc}%
0 & 0 & 0\\
0 & 3 & 0\\
0 & 0 & 4
\end{array}
\right)
\end{array}
\right)  \left(
\begin{array}
[c]{c}%
\mathbf{\ddot{r}}\left(  t\right) \\
\mathbf{\ddot{r}}\left(  t\right)
\end{array}
\right)  \right. \\
&  +\left.  \left(  \mathbf{\dot{r}}\left(  t\right)  ^{t}\mathbf{\dot{r}%
}\left(  t\right)  ^{t}\right)  \left(
\begin{array}
[c]{cc}%
\left(
\begin{array}
[c]{ccc}%
2 & 0 & 0\\
0 & 3 & 0\\
0 & 0 & 0
\end{array}
\right)  & \mathbb{O}\\
\mathbb{O} & \left(
\begin{array}
[c]{ccc}%
0 & 0 & 0\\
0 & 3 & 0\\
0 & 0 & 4
\end{array}
\right)
\end{array}
\right)  \left(
\begin{array}
[c]{c}%
\mathbf{\dot{r}}\left(  t\right) \\
\mathbf{\dot{r}}\left(  t\right)
\end{array}
\right)  \right\}  \delta t^{2}\\
&  +\text{higher order terms}%
\end{align}
If
\begin{equation}
\mathbf{r}\left(  0\right)  ^{t}\left(
\begin{array}
[c]{ccc}%
2 & 0 & 0\\
0 & 3 & 0\\
0 & 0 & 0
\end{array}
\right)  =\mathbf{0}%
\end{equation}
then we can find solutions $\mathbf{\dot{r}}\left(  0\right)  $
\begin{equation}
\left(
\begin{array}
[c]{ccc}%
2 & 0 & 0\\
0 & 3 & 0\\
0 & 0 & 0
\end{array}
\right)  \mathbf{\dot{r}}\left(  0\right)  =0
\end{equation}
that make the second order term in that direction zero and complement the
normal space. If
\begin{equation}
\mathbf{r}\left(  t\right)  ^{t}\left(
\begin{array}
[c]{ccc}%
2 & 0 & 0\\
0 & 3 & 0\\
0 & 0 & 0
\end{array}
\right)  =\mathbf{0}%
\end{equation}
or%
\begin{equation}
\mathbf{r}\left(  t\right)  ^{t}\left(
\begin{array}
[c]{ccc}%
0 & 0 & 0\\
0 & 3 & 0\\
0 & 0 & 4
\end{array}
\right)  =\mathbf{0}%
\end{equation}
then the leading term of the Taylor series for that component gives%
\begin{equation}
\mathbf{\dot{r}}\left(  t\right)  ^{t}\left(
\begin{array}
[c]{ccc}%
2 & 0 & 0\\
0 & 3 & 0\\
0 & 0 & 0
\end{array}
\right)  \mathbf{\dot{r}}\left(  t\right)  =0
\end{equation}
or%
\begin{equation}
\mathbf{\dot{r}}\left(  t\right)  ^{t}\left(
\begin{array}
[c]{ccc}%
0 & 0 & 0\\
0 & 3 & 0\\
0 & 0 & 4
\end{array}
\right)  \mathbf{\dot{r}}\left(  t\right)  =0.
\end{equation}
$\rfloor$

We can introduce the approximation
\begin{equation}
\Omega\left(  \lambda\right)  =\int\lambda\left(  \mathbf{x}\right)
\Phi\left(  \mathbf{x}\right)  ^{\dagger}\Phi\left(  \mathbf{x}\right)
d\mathbf{x\simeq}\sum_{1\leq j\leq r}\lambda_{j}\Phi_{j}^{\dagger}\Phi_{j}%
\end{equation}
allowing us to write
\begin{equation}
\left\langle \lambda\left\vert \left\{  \Xi\left(  a\right)  -\gamma\right\}
\right.  \right\rangle =\left(  a\left\vert \widehat{\Omega\left(
\lambda\right)  }\right.  a\right)  -\left\langle \lambda\left\vert \overset
{}{\gamma}\right.  \right\rangle .
\end{equation}
As we can expand $a$ as
\begin{equation}
a=\sum_{1\leq j,k\leq r_{N}}a_{jk}\Lambda_{jk}%
\end{equation}
the Lagrangian is a function of the variables $\left\{  \left.  a_{jk}%
\right\vert 1\leq j,k\leq r_{N}\right\}  ,\left\{  \lambda\in L_{\infty
}\left(  \mathbb{X}\right)  \right\}  $ and $\left\{  \gamma\in L_{1}\left(
\mathbb{X}\right)  \right\}  .$ The super operator action on $\mathcal{B}%
_{2sa}\left(  \mathcal{H}^{N}\right)  $ is defined by the semi bounded
\emph{real }symmetric quadratic forms
\begin{align}
Q_{1}\left(  a,b\right)   &  =\left(  a\left\vert \left\{  \left(
K+V_{2}\right)  b\right\}  \right.  \right)  =\operatorname{Tr}\left\{
a\left(  K+V_{2}\right)  b\right\} \nonumber\\
Q_{2}\left(  a,b\right)   &  =\left(  a\left\vert \left\{  \Omega\left(
\lambda\right)  b\right\}  \right.  \right)  =\operatorname{Tr}\left\{
a\Omega\left(  \lambda\right)  b\right\}  ,
\end{align}
which have domains
\begin{equation}
\mathcal{D}\left(  Q\right)  =\left\{  a\left\vert Q\left(  a\right)  \equiv
Q\left(  a,q\right)  \geq\mu\left\Vert a\right\Vert ^{2}\right.  \quad
-\infty<\mu_{Q}<\infty\right\}
\end{equation}
that are dense subsets of $\mathcal{B}_{2sa}\left(  \mathcal{H}^{N}\right)  .$
These define super operators $\widehat{X}$ with domains
\begin{equation}
\mathcal{D}\left(  \widehat{X}\right)  =\left\{  a\left\vert \overset{}{a\in
Q\left(  a\right)  \,\backepsilon\exists b\in\mathcal{B}_{2sa}\left(
\mathcal{H}^{N}\right)  \,\forall c\in\mathcal{D}\left(  Q\right)  \quad
Q\left(  a,b)=\left(  \left.  b\right\vert c\right)  \right)  }\right.
\right\}
\end{equation}
by
\begin{align}
\left(  a\left\vert \left(  \widehat{K+V_{2}}\right)  \right.  b\right)   &
=Q_{1}\left(  a,b\right) \nonumber\\
\left(  a\left\vert \widehat{\Omega\left(  \lambda\right)  }\right.  b\right)
&  =Q_{2}\left(  a,b\right)  ,
\end{align}
which are self adjoint operators $:\mathcal{B}_{2sa}\left(  \mathcal{H}%
^{N}\right)  \rightarrow\mathcal{B}_{2sa}\left(  \mathcal{H}^{N}\right)  .$
Note wrt the inner product $\left(  |\right)  $
\begin{align}
\left(  \widehat{K+V_{2}}\right)   &  =\left(  \widehat{K+V_{2}}\right)
^{t}\\
\widehat{\Omega\left(  \lambda\right)  }  &  =\widehat{\Omega\left(
\lambda\right)  }^{t}.
\end{align}


$\lceil$We always have the spectral decomposition
\begin{equation}
X=\int_{\sigma_{c}\left(  X\right)  }xP_{X}\left(  x\right)  dx+\sum_{j}%
x_{j}P_{X_{j}}%
\end{equation}
that gives in particular
\begin{equation}
X^{2}=\int_{\sigma_{c}\left(  X\right)  }x^{2}P_{X}\left(  x\right)
dx+\sum_{j}x_{j}^{2}P_{X_{j}}.
\end{equation}
Consider
\begin{equation}
a=\int_{\sigma_{c}\left(  X\right)  }\rho\left(  x\right)  P_{X}\left(
E\right)  dx+\sum_{j}\rho_{j}P_{X_{j}},\quad
\end{equation}
then
\begin{equation}
a^{2}=\int_{\sigma_{c}\left(  X\right)  }\rho\left(  x\right)  ^{2}%
P_{X}\left(  x\right)  dx+\sum_{j}\rho_{j}^{2}P_{X_{j}},
\end{equation}
where%
\begin{equation}
\int_{\sigma_{c}\left(  X\right)  }\rho\left(  x\right)  ^{2}dx+\,\sum_{j}%
\rho_{j}^{2}=1;\quad\left\vert \rho\left(  x\right)  \right\vert
\leq1;\left\vert x_{j}\right\vert \leq1.
\end{equation}
and $a^{2}$ is in general a mixed state. Thus
\begin{equation}
\left\vert \overset{}{a}\right)  =\left\vert \int_{\sigma_{c}\left(  X\right)
}\rho\left(  x\right)  P_{X}\left(  x\right)  dx+\sum_{j}\rho_{j}P_{X_{j}%
}\right)  =\int_{\sigma_{c}\left(  X\right)  }\rho\left(  x\right)  \left\vert
\overset{}{P_{X}\left(  x\right)  dx}\right)  dx+\sum_{j}\rho_{j}\left\vert
\overset{}{P_{X_{j}}}\right)  ,
\end{equation}
where $\left\vert a\right)  \in\mathcal{B}_{2sa}\left(  \mathcal{H}%
^{N}\right)  .$ If
\begin{equation}
X=\int_{\sigma_{c}\left(  X\right)  }xP_{X}\left(  x\right)  dx+\sum_{j}%
x_{j}P_{X_{j}}%
\end{equation}
then
\begin{align}
\widehat{X}\left\vert \overset{}{a}\right)   &  =\left\vert \overset{}%
{Xa}\right)  =\left\vert \left\{  \int_{\sigma_{c}\left(  X\right)  }%
xP_{X}\left(  x\right)  dx+\sum_{j}x_{j}P_{X_{j}}\right\}  \left\{
\int_{\sigma_{c}\left(  X\right)  }\rho\left(  x\right)  P_{X}\left(
x\right)  dx+\sum_{j}\rho_{j}P_{X_{j}}\right\}  \right) \nonumber\\
&  =\left\vert \int_{\sigma_{c}\left(  X\right)  }x\rho\left(  x\right)
P_{X}\left(  x\right)  dx+\sum_{j}x_{j}\rho_{j}P_{X_{j}}\right)  =\int
_{\sigma_{c}\left(  X\right)  }x\rho\left(  x\right)  \left\vert P_{X}\left(
x\right)  \right)  dx+\sum_{j}x_{j}\rho_{j}\left\vert P_{X_{j}}\right)  .
\end{align}
Thus the eigen vectors of $\widehat{X}$ are pure and mixed states of the form
\begin{equation}
\left\vert \overset{}{a}\right)  =\sum_{1\leq j\leq m}\rho_{j}\left\vert
P_{X_{j}}\right)  ,
\end{equation}
where $x_{j}=c$ for $1\leq j\leq m$.$\rfloor$

The normal Lagrangian \emph{for regular} constraining points can be expressed
as
\begin{subequations}
\begin{align}
\mathcal{L}\left(  a,\lambda,\gamma\right)   &  =\left(  a\left\vert \left(
\widehat{K+V_{2}}\right)  \right.  a\right)  -\left\langle \lambda\left\vert
\left\{  \left(  a\left\vert \widehat{\Phi^{\dagger}\Phi}\right.  a\right)
-\gamma\right\}  \right.  \right\rangle \nonumber\\
&  \equiv\left(  a\left\vert \left\{  \widehat{K+V_{2}}-\left\langle
\lambda\left\vert \widehat{\Phi^{\dagger}\Phi}\right.  \right\rangle \right\}
\right.  a\right)  +\left\langle \lambda\left\vert \overset{}{\gamma}\right.
\right\rangle \nonumber\\
&  \equiv\left(  a\left\vert \left\{  \widehat{K+V_{2}}-\widehat{\Omega\left(
\lambda\right)  }\right\}  \right.  a\right)  +\left\langle \lambda\left\vert
\overset{}{\gamma}\right.  \right\rangle .
\end{align}
It will be useful to express this Lagrangian as
\end{subequations}
\begin{equation}
\mathcal{L}\left(  a,\lambda,\gamma\right)  =\mathfrak{L}\left(
a,\lambda\right)  +\left\langle \lambda\left\vert \overset{}{\gamma}\right.
\right\rangle ,
\end{equation}
where
\begin{equation}
\mathfrak{L}\left(  a,\lambda\right)  =\left(  a\left\vert \left\{
\widehat{K+V_{2}}-\widehat{\Omega\left(  \lambda\right)  }\right\}  \right.
a\right)  .
\end{equation}
The Lagrangian can be expanded about a point $\left(  a,\lambda\right)  $ as
\begin{equation}
\mathcal{L}\left(  a+\delta a,\lambda+\delta\lambda,\gamma\right)  =\left(
a+\delta a\left\vert \left\{  \widehat{K+V_{2}}-\widehat{\Omega\left(
\lambda+\delta\lambda\right)  }\right\}  \right.  a+\delta a\right)
+\left\langle \lambda+\delta\lambda\left\vert \overset{}{\gamma}\right.
\right\rangle
\end{equation}
The first order term in the expansion of the Lagrangian in terms of the
coordinates $\left\{  \left.  a_{jk}\right\vert 1\leq j,k\leq r_{N}\right\}  $
and $\left\{  \left.  \lambda_{l}\right\vert 1\leq l\leq r\right\}  ,$ where
we approximate $\lambda$ and $\gamma$ by
\begin{align}
\lambda &  =\sum_{1\leq j\leq r}\lambda_{j}F_{j}\in\mathbb{R}^{r}\subset
L_{\infty}\left(  \mathbb{X}\right) \nonumber\\
\gamma &  =\sum_{1\leq j\leq r}\gamma_{j}G_{j}\in\mathbb{R}^{r}\subset
L_{1}\left(  \mathbb{X}\right)  ,
\end{align}
is
\begin{equation}
D\mathcal{L}\left(  a,\lambda,\gamma\right)  \left(  \delta a,\delta
\lambda\right)  =\left\langle \delta\lambda\left\vert \overset{}{\gamma
}\right.  \right\rangle -\left(  a\left\vert \widehat{\Omega\left(
\delta\lambda\right)  }\right.  a\right)  +2\left(  a\left\vert \left\{
\widehat{K+V_{2}}-\widehat{\Omega\left(  \lambda\right)  }\right\}  \right.
\delta a\right)  .
\end{equation}
The second order term is
\begin{equation}
D^{2}\mathcal{L}\left(  a,\lambda,\gamma\right)  \left(  \delta a,\delta
\lambda\right)  =2\left\{  \left(  \delta a\left\vert \left\{  \widehat
{K+V_{2}}-\widehat{\Omega\left(  \lambda\right)  }\right\}  \right.  \delta
a\right)  -\left(  a\left\vert \widehat{\Omega\left(  \delta\lambda\right)
}\right.  \delta a\right)  \right\}  .
\end{equation}
Using the bases $\left\{  \Lambda_{jk}\right\}  $ of $\mathcal{B}_{2sa}\left(
\mathcal{H}^{N}\right)  $ the first derivatives can be expressed as
\begin{equation}
D\mathcal{L}\left(  a,\lambda,\gamma\right)  =\left(  D_{a}\mathcal{L}\left(
a,\lambda,\gamma\right)  ,D_{\lambda}\mathcal{L}\left(  a,\lambda
,\gamma\right)  \right)  ,
\end{equation}
where
\begin{align}
D_{a}\mathcal{L}\left(  a,\lambda,\gamma\right)  _{jk}  &  =2\left(
a\left\vert \left\{  \widehat{K+V_{2}}-\widehat{\Omega\left(  \lambda\right)
}\right\}  \right.  \Lambda_{jk}\right) \nonumber\\
D_{\lambda}\mathcal{L}\left(  a,\lambda,\gamma\right)  _{j}  &  =\gamma
_{j}-\left(  a\left\vert \widehat{\Phi_{j}^{\dagger}\Phi_{j}}\right.
a\right)  =\Xi_{j}\left(  a\right)  .
\end{align}
The second order derivatives can be expressed as%
\begin{equation}
D^{2}\mathcal{L}\left(  a,\lambda,\gamma\right)  =\left(
\begin{array}
[c]{cc}%
D_{aa}^{2}\mathcal{L}\left(  a,\lambda,\gamma\right)  & D_{a\lambda}%
^{2}\mathcal{L}\left(  a,\lambda,\gamma\right) \\
D_{\lambda a}^{2}\mathcal{L}\left(  a,\lambda,\gamma\right)  & D_{\lambda}%
^{2}\mathcal{L}\left(  a,\lambda,\gamma\right)
\end{array}
\right)  =\left(
\begin{array}
[c]{cc}%
D_{aa}^{2}\mathcal{L}\left(  a,\lambda,\gamma\right)  & D_{a\lambda}%
^{2}\mathcal{L}\left(  a,\lambda,\gamma\right) \\
D_{\lambda a}^{2}\mathcal{L}\left(  a,\lambda,\gamma\right)  & \mathbb{O}%
\end{array}
\right)  ,
\end{equation}
where
\begin{align}
D_{aa}^{2}\mathcal{L}\left(  a,\lambda,\gamma\right)  _{lmjk}  &  =2\left(
\Lambda_{lm}\left\vert \left\{  \widehat{K+V_{2}}-\widehat{\Omega\left(
\lambda\right)  }\right\}  \right.  \Lambda_{jk}\right) \nonumber\\
D_{a\lambda}^{2}\mathcal{L}\left(  a,\lambda,\gamma\right)  _{jkl}  &
=-2\left(  a\left\vert \widehat{\Phi_{j}^{\dagger}\Phi_{j}}\right.
\Lambda_{jk}\right)  .
\end{align}
The normal space, $\mathcal{N}_{a_{0}},$ to the surface
\begin{equation}
\Xi\left(  a\right)  =\left(  a\left\vert \widehat{\Phi^{\dagger}\Phi}\right.
a\right)  =\left(
\begin{array}
[c]{c}%
\Xi_{1}\left(  a\right) \\
\vdots\\
\Xi_{r}\left(  a\right)
\end{array}
\right)  =\gamma
\end{equation}
at the point $a_{0}$ is spanned by the vectors
\begin{equation}
\nabla\Xi_{j}\left(  a_{0}\right)  =2\left(  \overset{}{a_{0}}\right\vert
\widehat{\Phi_{j}^{\dagger}\Phi_{j}},\quad1\leq j\leq r
\end{equation}
with coordinates
\begin{equation}
2\left(  a_{0}\left\vert \widehat{\Phi_{j}^{\dagger}\Phi_{j}}\right.
\Lambda_{kl}\right)  .
\end{equation}
One should observe that if $\gamma_{j}=0$ then $\left(  \overset{}{a_{0}%
}\right\vert \widehat{\Phi_{j}^{\dagger}\Phi_{j}}=0$ and vice versa as
$\widehat{\Phi_{j}^{\dagger}\Phi_{j}}$ is a positive operator and the
constraint function becomes
\begin{equation}
\Xi:\mathcal{B}_{2sa}\left(  \mathcal{H}^{N}\right)  \rightarrow
\mathbb{R}^{r^{\prime}};\quad r^{\prime}<r.
\end{equation}


The differential of the map $\Xi$ at the point $a_{0}$ is
\begin{align}
D_{a}\Xi\left(  a_{0}\right)   &  =\left(
\begin{array}
[c]{c}%
\nabla\Xi_{1}\left(  a_{0}\right) \\
\vdots\\
\nabla\Xi_{r}\left(  a_{0}\right)
\end{array}
\right)  =\left(
\begin{array}
[c]{ccc}%
\frac{\partial\Xi_{1}\left(  a_{0}\right)  }{\partial a_{11}} & \cdots &
\frac{\partial\Xi_{1}\left(  a_{0}\right)  }{\partial a_{1r_{N}}}\\
\vdots & \vdots & \vdots\\
\frac{\partial\Xi_{r}\left(  a_{0}\right)  }{\partial a_{r_{N}1}} & \cdots &
\frac{\partial\Xi_{r}\left(  a_{0}\right)  }{\partial a_{r_{N}r_{N}}}%
\end{array}
\right) \nonumber\\
&  =2\left(
\begin{array}
[c]{ccc}%
\left(  a_{0}\left\vert \widehat{\Phi_{1}^{\dagger}\Phi_{1}}\right.
\Lambda_{11}\right)  & \cdots & \left(  a_{0}\left\vert \widehat{\Phi
_{1}^{\dagger}\Phi_{1}}\right.  \Lambda_{r_{N}r_{N}}\right) \\
\vdots & \vdots & \vdots\\
\left(  a_{0}\left\vert \widehat{\Phi_{r}^{\dagger}\Phi_{r}}\right.
\Lambda_{11}\right)  & \cdots & \left(  a_{0}\left\vert \widehat{\Phi
_{r}^{\dagger}\Phi_{r}}\right.  \Lambda_{r_{N}r_{N}}\right)
\end{array}
\right)  .
\end{align}
Hence
\begin{equation}
D_{a}\Xi\left(  a_{0}\right)  \left(  \delta a\right)  =\delta\gamma,
\end{equation}
where we can write the matrix representation of this mapping in the following
various ways
\begin{align}
\left(
\begin{array}
[c]{c}%
\delta\gamma_{1}\\
\vdots\\
\delta\gamma_{r}%
\end{array}
\right)   &  =2\left(
\begin{array}
[c]{ccc}%
\left(  a_{0}\left\vert \widehat{\Phi_{1}^{\dagger}\Phi_{1}}\right.
\Lambda_{11}\right)  & \cdots & \left(  a_{0}\left\vert \widehat{\Phi
_{1}^{\dagger}\Phi_{1}}\right.  \Lambda_{r_{N}r_{N}}\right) \\
\vdots & \vdots & \vdots\\
\left(  a_{0}\left\vert \widehat{\Phi_{r}^{\dagger}\Phi_{r}}\right.
\Lambda_{11}\right)  & \cdots & \left(  a_{0}\left\vert \widehat{\Phi
_{r}^{\dagger}\Phi_{r}}\right.  \Lambda_{r_{N}r_{N}}\right)
\end{array}
\right)  \left(
\begin{array}
[c]{c}%
\delta a_{11}\\
\vdots\\
\delta a_{r_{N}r_{N}}%
\end{array}
\right) \nonumber\\
&  =2\left(
\begin{array}
[c]{c}%
\mathbf{a}_{0}^{t}\widehat{\mathbf{\Phi}_{1}^{\dagger}\mathbf{\Phi}_{1}}\\
\vdots\\
\mathbf{a}_{0}^{t}\widehat{\mathbf{\Phi}_{r}^{\dagger}\mathbf{\Phi}_{r}}%
\end{array}
\right)  \delta\mathbf{a}=2\left(
\begin{array}
[c]{ccc}%
\mathbf{a}_{0}^{t} & \cdots & \mathbf{a}_{0}^{t}%
\end{array}
\right)  \left(
\begin{array}
[c]{c}%
\widehat{\mathbf{\Phi}_{1}^{\dagger}\mathbf{\Phi}_{1}}\\
\vdots\\
\widehat{\mathbf{\Phi}_{r}^{\dagger}\mathbf{\Phi}_{r}}%
\end{array}
\right)  \delta\mathbf{a}\nonumber\\
&  =2\left(
\begin{array}
[c]{ccc}%
\mathbf{a}_{0}^{t} & \cdots & \mathbf{a}_{0}^{t}%
\end{array}
\right)  \left(
\begin{array}
[c]{c}%
\widehat{\mathbf{\Phi}_{1}^{\dagger}\mathbf{\Phi}_{1}}\delta\mathbf{a}\\
\vdots\\
\widehat{\mathbf{\Phi}_{r}^{\dagger}\mathbf{\Phi}_{r}}\delta\mathbf{a}%
\end{array}
\right)  =2\left(
\begin{array}
[c]{c}%
\mathbf{a}_{0}^{t}\widehat{\mathbf{\Phi}_{1}^{\dagger}\mathbf{\Phi}_{1}}%
\delta\mathbf{a}\\
\vdots\\
\mathbf{a}_{0}^{t}\widehat{\mathbf{\Phi}_{r}^{\dagger}\mathbf{\Phi}_{r}}%
\delta\mathbf{a}%
\end{array}
\right) \nonumber\\
&  =2\left(
\begin{array}
[c]{ccc}%
\mathbf{a}_{0}^{t} & \cdots & \mathbf{a}_{0}^{t}%
\end{array}
\right)  \left(
\begin{array}
[c]{ccc}%
\widehat{\mathbf{\Phi}_{1}^{\dagger}\mathbf{\Phi}_{1}} & \mathbb{O} &
\mathbb{O}\\
\mathbb{O} & \ddots & \mathbb{O}\\
\mathbb{O} & \mathbb{O} & \widehat{\mathbf{\Phi}_{r}^{\dagger}\mathbf{\Phi
}_{r}}%
\end{array}
\right)  \left(
\begin{array}
[c]{c}%
\delta\mathbf{a}\\
\vdots\\
\delta\mathbf{a}%
\end{array}
\right) \nonumber\\
&  =\left(
\begin{array}
[c]{ccc}%
\delta\gamma_{1} & \mathbb{O} & \mathbb{O}\\
\mathbb{O} & \ddots & \mathbb{O}\\
\mathbb{O} & \mathbb{O} & \delta\gamma_{r}%
\end{array}
\right)  ,
\end{align}
where $\widehat{\mathbf{\Phi}_{j}^{\dagger}\mathbf{\Phi}_{j}}$ is an
$r_{N}^{2}\times r_{N}^{2}$ matrix given by
\begin{equation}
\left(  \widehat{\mathbf{\Phi}_{j}^{\dagger}\mathbf{\Phi}_{j}}\right)
_{lmjk}=\operatorname{Tr}\left\{  \Lambda_{lm}\Phi_{j}^{\dagger}\Phi
_{j}\Lambda_{jk}\right\}  .
\end{equation}
If $\dim\left\{  \operatorname{Ran}\left\{  D_{a}\Xi\left(  a_{0}\right)
\right\}  \right\}  =r,$ i.e. the map $D_{a}\Xi\left(  a_{0}\right)  $ is
surjective, \emph{and }$\Xi\left(  a_{0}\right)  =\gamma$ then this point is
regular \cite{Hestenes}. The tangent space $\mathcal{T}_{a_{0}}$ to this level
surface at this point is the space of vectors
\begin{equation}
\mathcal{T}_{a_{0}}=\operatorname{linspan}\left\{  h\left\vert \left(
a_{0}\left\vert \widehat{\Phi_{j}^{\dagger}\Phi_{j}}\right.  h\right)
=0\right.  ;1\leq j\leq r;\text{ }\mathbf{v=\dot{r}}\left(  0\right)
;\Xi\left(  \mathbf{r}\left(  t\right)  \right)  =\gamma,\left\vert
t\right\vert \leq\delta;\mathbf{r}\left(  0\right)  =a_{0}\right\}  .
\end{equation}
Clearly
\begin{equation}
\mathcal{T}_{a_{0}}\subseteq\ker\left\{  D\Xi\left(  a_{0}\right)  \right\}
\subset\mathcal{B}_{2sa}\left(  \mathcal{H}^{N}\right)  .
\end{equation}
The normal space $\mathcal{N}_{a_{0}}$ at the point $a_{0}$ is defined as
\begin{equation}
\operatorname{linspan}\left\{  \left.  \left(  \overset{}{\nabla\Xi_{j}\left(
a_{0}\right)  }\right\vert =\left(  \overset{}{a_{0}}\right\vert \Phi
_{j}^{\dagger}\Phi_{j}\quad\right\vert 1\leq j\leq r\right\}  .
\end{equation}
For regular points $\ker\left\{  D\Xi\left(  a_{0}\right)  \right\}
=\mathcal{T}_{a_{0}}$, thus $\dim\left\{  \mathcal{T}_{a_{0}}\right\}
=\left(  r_{N}^{2}-r\right)  $ and
\begin{equation}
\mathcal{B}_{2sa}\left(  \mathcal{H}^{N}\right)  =\mathcal{T}_{a_{0}}%
\oplus\mathcal{N}_{a_{0}}.
\end{equation}
For non regular points $\mathcal{T}_{a_{0}}\oplus\mathcal{N}_{a_{0}}%
\subset\mathcal{B}_{2sa}\left(  \mathcal{H}^{N}\right)  .$

We can consider curves in the manifold generated by
\begin{align}
&  \left(
\begin{array}
[c]{ccc}%
\mathbf{a}_{0}^{t}\mathbf{V}\left(  t\right)  ^{t} & \cdots & \mathbf{a}%
_{0}^{t}\mathbf{V}\left(  t\right)  ^{t}%
\end{array}
\right)  \left(
\begin{array}
[c]{ccc}%
\widehat{\mathbf{\Phi}_{1}^{\dagger}\mathbf{\Phi}_{1}} & \mathbb{O} &
\mathbb{O}\\
\mathbb{O} & \ddots & \mathbb{O}\\
\mathbb{O} & \mathbb{O} & \widehat{\mathbf{\Phi}_{r}^{\dagger}\mathbf{\Phi
}_{r}}%
\end{array}
\right)  \left(
\begin{array}
[c]{c}%
\mathbf{V}\left(  t\right)  \mathbf{a}_{0}\\
\vdots\\
\mathbf{V}\left(  t\right)  \mathbf{a}_{0}%
\end{array}
\right)  =\left(
\begin{array}
[c]{ccc}%
\gamma_{1} & \mathbb{O} & \mathbb{O}\\
\mathbb{O} & \ddots & \mathbb{O}\\
\mathbb{O} & \mathbb{O} & \gamma_{r}%
\end{array}
\right) \\
=  &  \left(
\begin{array}
[c]{ccc}%
\mathbf{a}_{0}^{t} & \cdots & \mathbf{a}_{0}^{t}%
\end{array}
\right)  \left(
\begin{array}
[c]{ccc}%
\mathbf{V}\left(  t\right)  ^{t} & \mathbb{O} & \mathbb{O}\\
\mathbb{O} & \ddots & \mathbb{O}\\
\mathbb{O} & \mathbb{O} & \mathbf{V}\left(  t\right)  ^{t}%
\end{array}
\right)  \left(
\begin{array}
[c]{ccc}%
\widehat{\mathbf{\Phi}_{1}^{\dagger}\mathbf{\Phi}_{1}} & \mathbb{O} &
\mathbb{O}\\
\mathbb{O} & \ddots & \mathbb{O}\\
\mathbb{O} & \mathbb{O} & \widehat{\mathbf{\Phi}_{r}^{\dagger}\mathbf{\Phi
}_{r}}%
\end{array}
\right)  \times\\
&  \qquad\qquad\qquad\qquad\qquad\qquad\qquad\qquad\qquad\qquad\times\left(
\begin{array}
[c]{ccc}%
\mathbf{V}\left(  t\right)  & \mathbb{O} & \mathbb{O}\\
\mathbb{O} & \ddots & \mathbb{O}\\
\mathbb{O} & \mathbb{O} & \mathbf{V}\left(  t\right)
\end{array}
\right)  \left(
\begin{array}
[c]{c}%
\mathbf{a}_{0}\\
\vdots\\
\mathbf{a}_{0}%
\end{array}
\right) \\
&  =\left(
\begin{array}
[c]{ccc}%
\mathbf{a}_{0}^{t} & \cdots & \mathbf{a}_{0}^{t}%
\end{array}
\right)  \left(
\begin{array}
[c]{ccc}%
\mathbf{V}\left(  t\right)  ^{t}\widehat{\mathbf{\Phi}_{1}^{\dagger
}\mathbf{\Phi}_{1}}\mathbf{V}\left(  t\right)  & \mathbb{O} & \mathbb{O}\\
\mathbb{O} & \ddots & \mathbb{O}\\
\mathbb{O} & \mathbb{O} & \mathbf{V}\left(  t\right)  ^{t}\widehat
{\mathbf{\Phi}_{r}^{\dagger}\mathbf{\Phi}_{r}}\mathbf{V}\left(  t\right)
\end{array}
\right)  \left(
\begin{array}
[c]{c}%
\mathbf{a}_{0}\\
\vdots\\
\mathbf{a}_{0}%
\end{array}
\right)
\end{align}
so that
\begin{align}
&  \left(  V\left(  t\right)  a_{0}\left\vert \widehat{\mathbf{\Phi}%
_{j}^{\dagger}\mathbf{\Phi}_{j}}V\left(  t\right)  a_{0}\right.  \right)
=\gamma_{j}\\
&  \equiv\left(  \left[  \widehat{\mathbf{\Phi}_{j}^{\dagger}\mathbf{\Phi}%
_{j}}\right]  ^{\frac{1}{2}}V\left(  t\right)  \left[  \widehat{\mathbf{\Phi
}_{j}^{\dagger}\mathbf{\Phi}_{j}}\right]  ^{-\frac{1}{2}}\left[
\widehat{\mathbf{\Phi}_{j}^{\dagger}\mathbf{\Phi}_{j}}\right]  ^{\frac{1}{2}%
}a_{0}\left\vert \left[  \widehat{\mathbf{\Phi}_{j}^{\dagger}\mathbf{\Phi}%
_{j}}\right]  ^{\frac{1}{2}}V\left(  t\right)  \left[  \widehat{\mathbf{\Phi
}_{j}^{\dagger}\mathbf{\Phi}_{j}}\right]  ^{-\frac{1}{2}}\left[
\widehat{\mathbf{\Phi}_{j}^{\dagger}\mathbf{\Phi}_{j}}\right]  ^{\frac{1}{2}%
}a_{0}\right.  \right)
\end{align}
and clearly giving that
\begin{equation}
\left[  \widehat{\mathbf{\Phi}_{j}^{\dagger}\mathbf{\Phi}_{j}}\right]
^{\frac{1}{2}}V\left(  t\right)  \left[  \widehat{\mathbf{\Phi}_{j}^{\dagger
}\mathbf{\Phi}_{j}}\right]  ^{-\frac{1}{2}}\in O\left(  \mathbb{R}^{r_{N}^{2}%
}\right)  .
\end{equation}
Let
\begin{equation}
V\left(  t\right)  =e^{\mathsf{v}t}%
\end{equation}
then
\begin{equation}
\frac{d\mathbf{V}\left(  0\right)  ^{t}\widehat{\mathbf{\Phi}_{1}^{\dagger
}\mathbf{\Phi}_{1}}\mathbf{V}\left(  0\right)  }{dt}=\mathsf{v}^{t}%
\widehat{\mathbf{\Phi}_{1}^{\dagger}\mathbf{\Phi}_{1}}+\widehat{\mathbf{\Phi
}_{1}^{\dagger}\mathbf{\Phi}_{1}}\mathsf{v,}%
\end{equation}
giving
\begin{equation}
\frac{d\left(  a_{0}\left\vert V\left(  0\right)  ^{t}\widehat{\mathbf{\Phi
}_{j}^{\dagger}\mathbf{\Phi}_{j}}V\left(  0\right)  a_{0}\right.  \right)
}{dt}=\left(  a_{0}\left\vert \left\{  \mathsf{v}^{t}\mathbf{\Phi}%
_{j}^{\dagger}\mathbf{\Phi}_{j}+\widehat{\mathbf{\Phi}_{j}^{\dagger
}\mathbf{\Phi}_{j}}\mathsf{v}\right\}  a_{0}\right.  \right)  =0,
\end{equation}
hence
\begin{align}
&  \left(  a_{0}\left\vert \left\{  \mathsf{v}^{t}\widehat{\mathbf{\Phi}%
_{j}^{\dagger}\mathbf{\Phi}_{j}}\right\}  a_{0}\right.  \right)  =-\left(
a_{0}\left\vert \left\{  \widehat{\mathbf{\Phi}_{j}^{\dagger}\mathbf{\Phi}%
_{j}}\mathsf{v}\right\}  a_{0}\right.  \right)  =-\left(  \left.  \left\{
\widehat{\mathbf{\Phi}_{j}^{\dagger}\mathbf{\Phi}_{j}}\mathsf{v}\right\}
a_{0}\right\vert a_{0}\right) \\
&  =-\left(  a_{0}\left\vert \left\{  \mathsf{v}^{t}\widehat{\mathbf{\Phi}%
_{j}^{\dagger}\mathbf{\Phi}_{j}}\right\}  a_{0}\right.  \right)
\Rightarrow\left(  a_{0}\left\vert \left\{  \mathsf{v}^{t}\widehat
{\mathbf{\Phi}_{j}^{\dagger}\mathbf{\Phi}_{j}}\right\}  a_{0}\right.  \right)
=0.
\end{align}


The Hessian of the constraint map is
\begin{equation}
D_{a}^{2}\Xi\left(  a_{0}\right)  =2\widehat{\mathbf{\Phi}^{\dagger
}\mathbf{\Phi}}\equiv2\left(
\begin{array}
[c]{ccc}%
\widehat{\mathbf{\Phi}_{1}^{\dagger}\mathbf{\Phi}_{1}} & \mathbb{O} &
\mathbb{O}\\
\mathbb{O} & \ddots & \mathbb{O}\\
\mathbb{O} & \mathbb{O} & \widehat{\mathbf{\Phi}_{r}^{\dagger}\mathbf{\Phi
}_{r}}%
\end{array}
\right)  .
\end{equation}
The constraint map can be expanded about the point $a_{0}$ as
\begin{align}
\Xi\left(  a_{0}+h\right)   &  =\Xi\left(  a_{0}\right)  +D_{a}\Xi\left(
a_{0}\right)  \mathbf{h}+\frac{1}{2}\mathbf{h}^{t}D_{a}^{2}\Xi\left(
a_{0}\right)  \mathbf{h}\nonumber\\
&  =\Xi\left(  a_{0}\right)  +2\left(
\begin{array}
[c]{ccc}%
\mathbf{a}_{0}^{t} & \cdots & \mathbf{a}_{0}^{t}%
\end{array}
\right)  \left(
\begin{array}
[c]{ccc}%
\widehat{\mathbf{\Phi}_{1}^{\dagger}\mathbf{\Phi}_{1}} & \mathbb{O} &
\mathbb{O}\\
\mathbb{O} & \ddots & \mathbb{O}\\
\mathbb{O} & \mathbb{O} & \widehat{\mathbf{\Phi}_{r}^{\dagger}\mathbf{\Phi
}_{r}}%
\end{array}
\right)  \left(
\begin{array}
[c]{c}%
\mathbf{h}\\
\vdots\\
\mathbf{h}%
\end{array}
\right) \nonumber\\
&  +\left(
\begin{array}
[c]{ccc}%
\mathbf{h}^{t} & \cdots & \mathbf{h}^{t}%
\end{array}
\right)  \left(
\begin{array}
[c]{ccc}%
\widehat{\mathbf{\Phi}_{1}^{\dagger}\mathbf{\Phi}_{1}} & \mathbb{O} &
\mathbb{O}\\
\mathbb{O} & \ddots & \mathbb{O}\\
\mathbb{O} & \mathbb{O} & \widehat{\mathbf{\Phi}_{r}^{\dagger}\mathbf{\Phi
}_{r}}%
\end{array}
\right)  \left(
\begin{array}
[c]{c}%
\mathbf{h}\\
\vdots\\
\mathbf{h}%
\end{array}
\right)  \mathbf{.}%
\end{align}
$\mathbf{h\in}\mathcal{T}_{a_{0}}$ if for small enough $\varkappa$
\begin{align}
&  2\left(
\begin{array}
[c]{ccc}%
\mathbf{a}_{0}^{t} & \cdots & \mathbf{a}_{0}^{t}%
\end{array}
\right)  \left(
\begin{array}
[c]{ccc}%
\widehat{\mathbf{\Phi}_{1}^{\dagger}\mathbf{\Phi}_{1}} & \mathbb{O} &
\mathbb{O}\\
\mathbb{O} & \ddots & \mathbb{O}\\
\mathbb{O} & \mathbb{O} & \widehat{\mathbf{\Phi}_{r}^{\dagger}\mathbf{\Phi
}_{r}}%
\end{array}
\right)  \left(
\begin{array}
[c]{c}%
\mathbf{h}\\
\vdots\\
\mathbf{h}%
\end{array}
\right)  \varkappa\nonumber\\
+  &  \left(
\begin{array}
[c]{ccc}%
\mathbf{h}^{t} & \cdots & \mathbf{h}^{t}%
\end{array}
\right)  \left(
\begin{array}
[c]{ccc}%
\widehat{\mathbf{\Phi}_{1}^{\dagger}\mathbf{\Phi}_{1}} & \mathbb{O} &
\mathbb{O}\\
\mathbb{O} & \ddots & \mathbb{O}\\
\mathbb{O} & \mathbb{O} & \widehat{\mathbf{\Phi}_{r}^{\dagger}\mathbf{\Phi
}_{r}}%
\end{array}
\right)  \left(
\begin{array}
[c]{c}%
\mathbf{h}\\
\vdots\\
\mathbf{h}%
\end{array}
\right)  \varkappa^{2}=\left(
\begin{array}
[c]{c}%
0\\
\vdots\\
0
\end{array}
\right)  .
\end{align}
This gives%
\begin{align}
&  \left(
\begin{array}
[c]{ccc}%
\left(  a_{0}\left\vert \widehat{\mathbf{\Phi}_{1}^{\dagger}\mathbf{\Phi}_{1}%
}\right.  h\right)  & \mathbb{O} & \mathbb{O}\\
\mathbb{O} & \ddots & \mathbb{O}\\
\mathbb{O} & \mathbb{O} & \left(  a_{0}\left\vert \widehat{\mathbf{\Phi}%
_{r}^{\dagger}\mathbf{\Phi}_{r}}\right.  h\right)
\end{array}
\right)  \varkappa\nonumber\\
&  \qquad\qquad\qquad\qquad\qquad\left.  +\right.  \left(
\begin{array}
[c]{ccc}%
\left(  h\left\vert \widehat{\mathbf{\Phi}_{1}^{\dagger}\mathbf{\Phi}_{1}%
}\right.  h\right)  & \mathbb{O} & \mathbb{O}\\
\mathbb{O} & \ddots & \mathbb{O}\\
\mathbb{O} & \mathbb{O} & \left(  h\left\vert \widehat{\mathbf{\Phi}%
_{r}^{\dagger}\mathbf{\Phi}_{r}}\right.  h\right)
\end{array}
\right)  \varkappa^{2}=\left(
\begin{array}
[c]{c}%
0\\
\vdots\\
0
\end{array}
\right)  .
\end{align}
Hence $\mathbf{h\in}\mathcal{T}_{a_{0}}$ if (sufficient condition)
\begin{equation}
h\mathbf{\bot}\left(  a_{0}\right\vert \widehat{\mathbf{\Phi}_{j}^{\dagger
}\mathbf{\Phi}_{j}};\quad1\leq j\leq r\Leftrightarrow h\in\ker\left\{
D_{a}\Xi\left(  a_{0}\right)  \right\}
\end{equation}
and
\begin{equation}
h\in\bigcap\limits_{1\leq j\leq r}\ker\left\{  \widehat{\mathbf{\Phi}%
_{j}^{\dagger}\mathbf{\Phi}_{j}}\right\}  .
\end{equation}
Let $a_{0}$ be an IPS\ wlog we can consider the occupied space $\left\{
\left.  \widehat{\mathbf{\Phi}_{j}^{\dagger}\mathbf{\Phi}_{j}}\right\vert
1\leq j\leq N\right\}  $, then $h\in\ker\left\{  D_{a}\Xi\left(  a_{0}\right)
\right\}  $ if $a_{0}\mathbf{\bot}h$ e.g.%
\begin{align}
a_{0}  &  =\left\vert \psi_{1}\cdots\psi_{N}\right\rangle \left\langle
\psi_{1}\cdots\psi_{N}\right\vert \nonumber\\
h  &  =\left\vert \psi_{1}\cdots\psi_{N}\right\rangle \left\langle \psi
_{2}\cdots\psi_{N+1}\right\vert +\left\vert \psi_{2}\cdots\psi_{N+1}%
\right\rangle \left\langle \psi_{1}\cdots\psi_{N}\right\vert .
\end{align}
The manifold is then given by
\begin{equation}
\left(  a_{0}\left\vert \widehat{\mathbf{\Phi}_{j}^{\dagger}\mathbf{\Phi}_{j}%
}a_{0}\right.  \right)  =1;1\leq j\leq N
\end{equation}
i.e. is an intersection of hyperspheres in $\mathcal{B}_{2sa}\left(
\mathcal{H}^{N}\right)  .$

If $h$ is a tangent vector then (necessary condition)
\begin{equation}
\frac{d}{dt}\left.  \left(
\begin{array}
[c]{ccc}%
\left(  a_{0}+ht\left\vert \widehat{\mathbf{\Phi}_{1}^{\dagger}\mathbf{\Phi
}_{1}}\right.  a_{0}+ht\right)  & \mathbb{O} & \mathbb{O}\\
\mathbb{O} & \ddots & \mathbb{O}\\
\mathbb{O} & \mathbb{O} & \left(  a_{0}+ht\left\vert \widehat{\mathbf{\Phi
}_{r}^{\dagger}\mathbf{\Phi}_{r}}\right.  a_{0}+ht\right)
\end{array}
\right)  \right\vert _{t=0}=D_{a}\Xi\left(  a_{0}\right)  \left(  h\right)
=\mathbf{0.}%
\end{equation}
The state defined by $a_{0}+ht$ is $\left(  a_{0}+ht\right)  ^{2}%
=a_{0}+\left\{  a_{0}h+ha_{0}\right\}  t+h^{2}t^{2}$. If $a_{0}$ is an IPS
given by $\left\vert \psi_{1}\cdots\psi_{N}\right\rangle \left\langle \psi
_{1}\cdots\psi_{N}\right\vert $ we can consider the basis
\begin{align}
\Lambda_{JK}  &  =\frac{1}{\sqrt{2}}\left\{  \left\vert \psi_{j_{1}}\cdots
\psi_{j_{N}}\right\rangle \left\langle \psi_{k_{1}}\cdots\psi_{k_{N}%
}\right\vert +\left\vert \psi_{k_{1}}\cdots\psi_{k_{N}}\right\rangle
\left\langle \psi_{j_{1}}\cdots\psi_{j_{N}}\right\vert \right\}  ;\quad1\leq
J<K\leq r_{N}\nonumber\\
\Lambda_{JJ}  &  =\left\vert \psi_{j_{1}}\cdots\psi_{j_{N}}\right\rangle
\left\langle \psi_{j_{1}}\cdots\psi_{j_{N}}\right\vert ;\quad1\leq J\leq
r_{N}\nonumber\\
\Lambda_{JK}  &  =\frac{i}{\sqrt{2}}\left\{  \left\vert \psi_{j_{1}}\cdots
\psi_{j_{N}}\right\rangle \left\langle \psi_{k_{1}}\cdots\psi_{k_{N}%
}\right\vert -\left\vert \psi_{k_{1}}\cdots\psi_{k_{N}}\right\rangle
\left\langle \psi_{j_{1}}\cdots\psi_{j_{N}}\right\vert \right\}  ;\quad1\leq
J>K\leq r_{N}. \label{basis}%
\end{align}
The necessary tangential condition for occupied orbitals gives
\begin{equation}
\left(  a_{0}\left\vert \widehat{\mathbf{\Phi}_{j}^{\dagger}\mathbf{\Phi}_{j}%
}\right.  h\right)  =0\text{ for }\left(  a_{0}\left\vert \widehat
{\mathbf{\Phi}_{j}^{\dagger}\mathbf{\Phi}_{j}}\right.  a_{0}\right)
=1;\quad1\leq j\leq N
\end{equation}
gives
\begin{equation}
\left(  a_{0}\left\vert h\right.  \right)  =0.
\end{equation}
The states produced along the direction $h$ are
\begin{equation}
\left(  a_{0}+ht\right)  ^{2}=a_{0}+\left\{  ha_{0}+a_{0}h\right\}
t+h^{2}t^{2}.
\end{equation}
For basis elements $h$
\begin{align}
h^{2}  &  =\left\{  \frac{1}{\sqrt{2}}\left\{  \left\vert \psi_{j_{1}}%
\cdots\psi_{j_{N}}\right\rangle \left\langle \psi_{k_{1}}\cdots\psi_{k_{N}%
}\right\vert \pm\left\vert \psi_{k_{1}}\cdots\psi_{k_{N}}\right\rangle
\left\langle \psi_{j_{1}}\cdots\psi_{j_{N}}\right\vert \right\}  \right\}
^{2}=\nonumber\\
&  \pm\frac{1}{2}\left\{  \left\vert \psi_{j_{1}}\cdots\psi_{j_{N}%
}\right\rangle \left\langle \psi_{j_{1}}\cdots\psi_{j_{N}}\right\vert
+\left\vert \psi_{k_{1}}\cdots\psi_{k_{N}}\right\rangle \left\langle
\psi_{k_{1}}\cdots\psi_{k_{N}}\right\vert \right\}  .
\end{align}
and $ha_{0}$ and $a_{0}h$ are only non zero when
\begin{equation}
h=\left\vert \psi_{1}\cdots\psi_{N}\right\rangle \left\langle \psi_{k_{1}%
}\cdots\psi_{k_{N}}\right\vert \pm\left\vert \psi_{k_{1}}\cdots\psi_{k_{N}%
}\right\rangle \left\langle \psi_{1}\cdots\psi_{N}\right\vert ;k_{1}%
,\cdots,k_{N}\neq1,\cdots,N \label{excit}%
\end{equation}
giving
\begin{align}
\left\{  \left\vert \psi_{1}\cdots\psi_{N}\right\rangle \left\langle
\psi_{k_{1}}\cdots\psi_{k_{N}}\right\vert \pm\left\vert \psi_{k_{1}}\cdots
\psi_{k_{N}}\right\rangle \left\langle \psi_{1}\cdots\psi_{N}\right\vert
\right\}  \left\vert \psi_{1}\cdots\psi_{N}\right\rangle \left\langle \psi
_{1}\cdots\psi_{N}\right\vert  &  =\pm\left\vert \psi_{k_{1}}\cdots\psi
_{k_{N}}\right\rangle \left\langle \psi_{1}\cdots\psi_{N}\right\vert
\nonumber\\
\left\vert \psi_{1}\cdots\psi_{N}\right\rangle \left\langle \psi_{1}\cdots
\psi_{N}\right\vert \left\{  \left\vert \psi_{1}\cdots\psi_{N}\right\rangle
\left\langle \psi_{k_{1}}\cdots\psi_{k_{N}}\right\vert \pm\left\vert
\psi_{k_{1}}\cdots\psi_{k_{N}}\right\rangle \left\langle \psi_{1}\cdots
\psi_{N}\right\vert \right\}   &  =\left\vert \psi_{1}\cdots\psi
_{N}\right\rangle \left\langle \psi_{k_{1}}\cdots\psi_{k_{N}}\right\vert
\end{align}
and hence
\begin{equation}
ha_{0}+a_{0}h=\frac{1}{\sqrt{2}}\left\{  \left\vert \psi_{1}\cdots\psi
_{N}\right\rangle \left\langle \psi_{k_{1}}\cdots\psi_{k_{N}}\right\vert
\pm\left\vert \psi_{k_{1}}\cdots\psi_{k_{N}}\right\rangle \left\langle
\psi_{1}\cdots\psi_{N}\right\vert \right\}  .
\end{equation}
The space generated by basis vectors of the form of eq. $\left(
\ref{excit}\right)  $ is
\begin{equation}
\mathcal{M}_{a_{0}}=\mathbb{R}-\operatorname{linspan}\left\{  \left.
\overset{}{\left\vert \psi_{1}\cdots\psi_{N}\right\rangle \left\langle
\psi_{k_{1}}\cdots\psi_{k_{N}}\right\vert \pm\left\vert \psi_{k_{1}}\cdots
\psi_{k_{N}}\right\rangle \left\langle \psi_{1}\cdots\psi_{N}\right\vert
}\right\vert k_{1},\cdots,k_{N}\neq1,\cdots,N\right\}  .
\end{equation}
The curves generated by the basis vectors are eq. $\left(  \ref{basis}\right)
$
\begin{equation}
\mathfrak{f}_{a_{0},h}\left(  t\right)  =\left(  a_{0}+ht\left\vert
\widehat{\Phi^{\dagger}\Phi}\right.  a_{0}+ht\right)  =\operatorname{Tr}%
\left\{  \overset{}{\Phi^{\dagger}\Phi\left\{  \overset{}{a_{0}+\left\{
ha_{0}+a_{0}h\right\}  t+h^{2}t^{2}}\right\}  }\right\}
\end{equation}
If $h\in$ basis eq. $\left(  \ref{excit}\right)  $ then
\begin{align}
\mathfrak{f}_{a_{0},h}\left(  t\right)   &  =\operatorname{Tr}\left\{
\overset{}{\Phi^{\dagger}\Phi\left\{  \overset{}{a_{0}+\left\{  ha_{0}%
+a_{0}h\right\}  t+h^{2}t^{2}}\right\}  }\right\} \nonumber\\
&  =\operatorname{Tr}\left\{  \overset{}{\Phi^{\dagger}\Phi\left\{  \overset
{}{a_{0}+\left\{  ha_{0}+a_{0}h\right\}  t+\frac{a_{0}}{2}t^{2}}\right\}
}\right\}  \neq\gamma,
\end{align}
while curves generated by a basis vector not in this basis are given by
\begin{align}
\mathfrak{g}_{a_{0},h}\left(  t\right)   &  =\operatorname{Tr}\left\{
\overset{}{\Phi^{\dagger}\Phi\left\{  \overset{}{a_{0}+h^{2}t^{2}}\right\}
}\right\} \nonumber\\
&  =\operatorname{Tr}\left\{  \overset{}{\Phi^{\dagger}\Phi\left\{  \overset
{}{a_{0}+h^{2}t^{2}}\right\}  }\right\}  =\gamma.
\end{align}
The curves $\mathfrak{g}_{a_{0},h}\left(  t\right)  $ lie in the constraint
manifold, while the curves $\mathfrak{f}_{a_{0},h}\left(  t\right)  $ do not.
The derivatives of these curves are
\begin{align}
\frac{d\mathfrak{f}_{a_{0},h}\left(  t\right)  }{dt}  &  =\operatorname{Tr}%
\left\{  \overset{}{\Phi^{\dagger}\Phi\left\{  \overset{}{\left\{
ha_{0}+a_{0}h\right\}  +a_{0}t}\right\}  }\right\}  \neq\gamma\\
\frac{d\mathfrak{g}_{a_{0},h}\left(  t\right)  }{dt}  &  =\operatorname{Tr}%
\left\{  \overset{}{\Phi^{\dagger}\Phi\overset{}{2h^{2}t}}\right\}  =\gamma,
\end{align}
which shows that the latter lines lie in the surface and cannot thus lie in
the tangent plane unless the surface is planar. The preceding equations show
that the surface is not planar, thus $h\notin\mathcal{M}_{a_{0}}$ cannot
produce tangent lines. Hence $\mathcal{M}_{a_{0}}$ is the tangent space. Note
in both cases
\begin{equation}
\frac{d\mathfrak{f}_{a_{0},h}\left(  0\right)  }{dt}=\frac{d\mathfrak{g}%
_{a_{0},h}\left(  0\right)  }{dt}=\mathbf{0}.
\end{equation}


The differential can always be expanded as

$\lceil$%
\begin{equation}
D_{a}\Xi\left(  a_{0}\right)  =\sum\limits_{1\leq j\leq r^{\prime}}%
\varepsilon_{j}\left\vert w_{j}\right)  \left(  v_{j}^{d}\right\vert ,
\end{equation}
where $\left(  v_{j}\right\vert \in\mathcal{B}_{2sa}\left(  \mathcal{H}%
^{N}\right)  $ and $\left\vert w_{j}\right)  \in L_{1}\left(  \mathbb{X}%
\right)  .$ The generalized inverse is
\begin{equation}
D_{a}\Xi\left(  a_{0}\right)  ^{G}=\sum\limits_{1\leq j\leq r^{\prime}}\left(
\varepsilon_{j}\right)  ^{-1}\left\vert v_{j}\right)  \left(  w_{j}%
^{d}\right\vert ,
\end{equation}
where $\left\vert v_{j}^{d}\right)  \in\mathcal{B}_{2sa}\left(  \mathcal{H}%
^{N}\right)  $ and $\left(  w_{j}^{d}\right\vert \in L_{\infty}\left(
\mathbb{X}\right)  .$Hence
\begin{align}
D_{a}\Xi\left(  a_{0}\right)  D_{a}\Xi\left(  a_{0}\right)  ^{G}  &
=\sum\limits_{1\leq j\leq r^{\prime}}\left\vert w_{j}\right)  \left(
w_{j}^{d}\right\vert =\mathcal{P}_{\operatorname{Ran}\left\{  D_{a}\Xi\left(
a_{0}\right)  \right\}  }\\
D_{a}\Xi\left(  a_{0}\right)  ^{G}D_{a}\Xi\left(  a_{0}\right)   &
=\sum\limits_{1\leq j\leq r^{\prime}}\left\vert v_{j}^{d}\right)  \left(
v_{j}\right\vert =\mathcal{P}_{\ker\left\{  D_{a}\Xi\left(  a_{0}\right)
\right\}  ^{\bot}}%
\end{align}
and $\dim\left\{  \operatorname{Ran}\left\{  D_{a}\Xi\left(  a_{0}\right)
\right\}  \right\}  =\dim\left\{  \ker\left\{  D_{a}\Xi\left(  a_{0}\right)
\right\}  ^{\bot}\right\}  =r^{\prime}.\rfloor$

The the first order KKT stationarity conditions for regular points are
satisfied by solutions $\left\{  a_{k}\right\}  $ to the equation
\begin{equation}
\left(  \overset{}{a_{k}}\right\vert \left\{  \left(  \widehat{K+V_{2}%
}\right)  -\widehat{\Omega\left(  \lambda\right)  }\right\}  \Leftrightarrow
\left\{  \left(  \widehat{K+V_{2}}\right)  -\widehat{\Omega\left(
\lambda\right)  }\right\}  \left\vert \overset{}{a_{k}}\right)  =0
\label{EnStation}%
\end{equation}
i.e.
\begin{equation}
\left\{  \left(  \widehat{K+V_{2}}\right)  -\widehat{\Omega\left(
\lambda\right)  }\right\}  \left\vert \overset{}{a_{k}}\right)  =0.
\end{equation}
These correspond to solutions $\left\{  a_{k}\right\}  $ of the equations
\begin{equation}
\nabla_{a}E_{U}\left(  a_{k}\right)  \in\mathcal{N}_{a_{k}}\Leftrightarrow
\nabla_{a}E_{U}\left(  a_{k}\right)  \left(  h\right)  =0,\,\,h\in
\mathcal{T}_{a_{k}}%
\end{equation}
and
\begin{equation}
\left(  a_{k}\left\vert \widehat{\Phi^{\dagger}\Phi}\right.  a_{k}\right)
-\gamma=0. \label{surface}%
\end{equation}
The big question is whether we can always find a $\lambda$ for a given
$\gamma$ such that eq. $\left(  \ref{EnStation}\right)  $ has at least one
solution $\left\vert \overset{}{a}\right)  .$ In finite dimensions it always
has $r$ solutions, but in $\infty$ dimensions it may have none.

The super operator $\widehat{\Omega\left(  \lambda\right)  }$ is a local
potential and eq. $\left(  \ref{EnStation}\right)  $ is an eigen equation for
a ground state in this potential that has a Hamiltonian that is shifted to
make it positive

$\lceil$%
\begin{align}
H\left\vert \Psi\right\rangle  &  =\left(  H_{0}+V\right)  \left\vert
\Psi\right\rangle =E\left\vert \Psi\right\rangle \nonumber\\
&  \equiv\left(  H_{0}+V-EI\right)  \left\vert \Psi\right\rangle =0
\end{align}
$\rfloor$. One can find the solution to this equation by finding generalized
eigenvectors and eigenvalues $\left\{  \left\vert a_{k}\left(  \lambda\right)
\right)  ,\eta_{k}\left(  \lambda\right)  \right\}  ,$ of $\widehat{K+V_{2}}$
with respect to $\widehat{\Omega\left(  \lambda\right)  },$ for a fixed
arbitrary $\lambda$ corresponding to a nonsingular super operator
$\widehat{\Omega\left(  \lambda\right)  ,}$ that produces at least one
eigenvalue, then defining
\begin{equation}
\tilde{\lambda}_{k}=\eta_{k}\left(  \lambda\right)  \lambda.
\end{equation}
$\lceil$Is $\widehat{\Omega\left(  \lambda\right)  }$ nonsingular?
\begin{equation}
\widehat{\Omega\left(  \lambda\right)  }=\sum_{1\leq j\leq r}\lambda
_{j}\widehat{\mathbf{\Phi}_{j}^{\dagger}\mathbf{\Phi}_{j}}%
\end{equation}
If it is singular it has a zero eigenvalue and
\begin{equation}
\sum_{1\leq j\leq r}\lambda_{j}\widehat{\mathbf{\Phi}_{j}^{\dagger
}\mathbf{\Phi}_{j}}\left\vert \overset{}{a\left(  \lambda\right)  }\right)  =0
\end{equation}
and the vectors $\left\{  \left.  \widehat{\mathbf{\Phi}_{j}^{\dagger
}\mathbf{\Phi}_{j}}\left\vert \overset{}{a\left(  \lambda\right)  }\right)
\right\vert 1\leq j\leq r\right\}  $ are linearly dependent and conversely.
Moreover in the singular case if
\begin{equation}
\left\{  \left(  \widehat{K+V_{2}}\right)  -\widehat{\Omega\left(
\lambda\right)  }\right\}  \left\vert \overset{}{a\left(  \lambda\right)
}\right)  =0
\end{equation}
then
\begin{equation}
\left(  \widehat{K+V_{2}}\right)  \left\vert \overset{}{a\left(
\lambda\right)  }\right)  =0.
\end{equation}
In the continuous case $\widehat{\Omega\left(  \lambda\right)  }$ if singular
then there exists at least one state $\left\vert a\left(  \lambda\right)
\right)  $ such that
\begin{equation}
\int\lambda\left(  \mathbf{x}\right)  \widehat{\mathbf{\Phi}\left(
\mathbf{x}\right)  ^{\dagger}\mathbf{\Phi}\left(  \mathbf{x}\right)
}\left\vert \overset{}{a\left(  \lambda\right)  }\right)  d\mathbf{x}=0
\end{equation}
and the vectors $\left\{  \widehat{\mathbf{\Phi}\left(  \mathbf{x}\right)
^{\dagger}\mathbf{\Phi}\left(  \mathbf{x}\right)  }\left\vert \overset
{}{a\left(  \lambda\right)  }\right)  \right\}  $ do not form a basis. Is
$D_{a}\Xi\left(  a_{0}\right)  $ surjective? $\lceil W$ is a proper subspace
of $L_{1}\left(  \Delta\right)  $ if
\begin{equation}
W\subset L_{1}\left(  \Delta\right)  \Rightarrow\exists\lambda\in L_{\infty
}\left(  \Delta\right)  \backepsilon\left\langle \left.  \lambda\right\vert
w\right\rangle =0\forall w\in W\rfloor
\end{equation}
If it is not then there exists a non zero $\lambda$ such that
\begin{equation}
\int\lambda\left(  \mathbf{x}\right)  \left(  a_{0}\left\vert \widehat
{\mathbf{\Phi}\left(  \mathbf{x}\right)  ^{\dagger}\mathbf{\Phi}\left(
\mathbf{x}\right)  }\right.  \delta a\right)  d\mathbf{x}=0\quad\forall\delta
a\,
\end{equation}
which of course implies that
\begin{equation}
\int\lambda\left(  \mathbf{x}\right)  \widehat{\mathbf{\Phi}\left(
\mathbf{x}\right)  ^{\dagger}\mathbf{\Phi}\left(  \mathbf{x}\right)
}\left\vert \overset{}{a_{0}}\right)  d\mathbf{x}=\widehat{\Omega\left(
\lambda\right)  }\left\vert \overset{}{a_{0}}\right)  =0.
\end{equation}
We can always limit the integration to
\begin{equation}
\int_{\Delta}\lambda\left(  \mathbf{x}\right)  \widehat{\mathbf{\Phi}\left(
\mathbf{x}\right)  ^{\dagger}\mathbf{\Phi}\left(  \mathbf{x}\right)
}\left\vert \overset{}{a_{0}}\right)  d\mathbf{x}=\widehat{\Omega\left(
\lambda\right)  }\left\vert \overset{}{a_{0}}\right)  =0\,,
\end{equation}
where
\begin{equation}
\Delta=\left\{  \mathbf{x}\left\vert \widehat{\mathbf{\Phi}\left(
\mathbf{x}\right)  ^{\dagger}\mathbf{\Phi}\left(  \mathbf{x}\right)
}\left\vert \overset{}{a_{0}}\right)  \neq0\right.  \right\}  .
\end{equation}
Hence
\begin{equation}
\operatorname{Tr}\left\{  b\Omega\left(  \lambda\right)  a_{0}\right\}
=0\forall b
\end{equation}
in particular if $a_{0}=$ $\left\vert \Psi_{1}\right\rangle \left\langle
\Psi_{1}\right\vert $ then
\begin{equation}
\operatorname{Tr}\left\{  \left\{  \left\vert \Psi_{k}\right\rangle
\left\langle \Psi_{1}\right\vert \pm\left\vert \Psi_{k}\right\rangle
\left\langle \Psi_{1}\right\vert \right\}  \Omega\left(  \lambda\right)
\left\vert \Psi_{1}\right\rangle \left\langle \Psi_{1}\right\vert \right\}
=0,
\end{equation}
and
\begin{equation}
\left\langle \Psi_{k}\right\vert \Omega\left(  \lambda\right)  \left\vert
\Psi_{1}\right\rangle =0\forall k
\end{equation}
so
\begin{equation}
\Omega\left(  \lambda\right)  \left\vert \Psi_{1}\right\rangle =0.
\end{equation}
This implies that
\begin{equation}
e^{\Omega\left(  \lambda\right)  }\left\vert \Psi_{1}\right\rangle =\left\vert
\Psi_{1}\right\rangle ,
\end{equation}
which gives
\begin{equation}
\Omega\left(  \lambda\right)  \left\vert \psi_{1}\cdots\psi_{N}\right\rangle
=\left\vert \Omega\left(  \lambda\right)  \psi_{1}\cdots\psi_{N}\right\rangle
+\cdots+\left\vert \psi_{1}\cdots\Omega\left(  \lambda\right)  \psi
_{N}\right\rangle .
\end{equation}
The following theorem is very important%
\begin{align}
\Omega\left(  \lambda\right)  \left\vert \psi_{1}\cdots\psi_{N}\right\rangle
& =\left\vert \Omega\left(  \lambda\right)  \psi_{1}\cdots\psi_{N}%
\right\rangle +\cdots+\left\vert \psi_{1}\cdots\Omega\left(  \lambda\right)
\psi_{N}\right\rangle =0\nonumber\\
& \Leftrightarrow\Omega\left(  \lambda\right)  \left\vert \psi_{j}%
\right\rangle =0\Leftrightarrow\lambda=0.
\end{align}
One can see this by considering $N=2$
\begin{equation}
\left\langle \left.  \mathbf{xx}^{\prime}\right\vert \psi_{1}\psi
_{2}\right\rangle =\frac{1}{\sqrt{2}}\left\{  \psi_{1}\left(  \mathbf{x}%
\right)  \psi_{2}\left(  \mathbf{x}^{\prime}\right)  -\psi_{2}\left(
\mathbf{x}\right)  \psi_{1}\left(  \mathbf{x}^{\prime}\right)  \right\}
\end{equation}
then%
\begin{align}
& \left\langle \left.  \mathbf{xx}^{\prime}\right\vert \Omega\left(
\lambda\right)  \psi_{1}\psi_{2}\right\rangle +\left\langle \left.
\mathbf{x}\right\vert \psi_{1}\Omega\left(  \lambda\right)  \psi
_{2}\right\rangle \nonumber\\
& =\frac{1}{\sqrt{2}}\left\{  \lambda\left(  \mathbf{x}\right)  \psi
_{1}\left(  \mathbf{x}\right)  \psi_{2}\left(  \mathbf{x}^{\prime}\right)
-\psi_{2}\left(  \mathbf{x}\right)  \lambda\left(  \mathbf{x}^{\prime}\right)
\psi_{1}\left(  \mathbf{x}^{\prime}\right)  +\psi_{1}\left(  \mathbf{x}%
\right)  \lambda\left(  \mathbf{x}^{\prime}\right)  \psi_{2}\left(
\mathbf{x}^{\prime}\right)  -\lambda\left(  \mathbf{x}\right)  \psi_{2}\left(
\mathbf{x}\right)  \psi_{1}\left(  \mathbf{x}^{\prime}\right)  \right\}
\nonumber\\
& =\frac{1}{\sqrt{2}}\left\{  \lambda\left(  \mathbf{x}\right)  \left\{
\psi_{1}\left(  \mathbf{x}\right)  \psi_{2}\left(  \mathbf{x}^{\prime}\right)
-\psi_{2}\left(  \mathbf{x}\right)  \psi_{1}\left(  \mathbf{x}^{\prime
}\right)  \right\}  +\lambda\left(  \mathbf{x}^{\prime}\right)  \left\{
\psi_{1}\left(  \mathbf{x}\right)  \psi_{2}\left(  \mathbf{x}^{\prime}\right)
-\psi_{2}\left(  \mathbf{x}\right)  \psi_{1}\left(  \mathbf{x}^{\prime
}\right)  \right\}  \right\}  \nonumber\\
& =\frac{1}{\sqrt{2}}\left\{  \lambda\left(  \mathbf{x}\right)  +\lambda
\left(  \mathbf{x}^{\prime}\right)  \right\}  \left\{  \psi_{1}\left(
\mathbf{x}\right)  \psi_{2}\left(  \mathbf{x}^{\prime}\right)  -\psi
_{2}\left(  \mathbf{x}\right)  \psi_{1}\left(  \mathbf{x}^{\prime}\right)
\right\}  .
\end{align}
Hence if
\begin{equation}
\left\langle \left.  \mathbf{xx}^{\prime}\right\vert \Omega\left(
\lambda\right)  \psi_{1}\psi_{2}\right\rangle +\left\langle \left.
\mathbf{xx}^{\prime}\right\vert \psi_{1}\Omega\left(  \lambda\right)  \psi
_{2}\right\rangle \overset{ae}{=}0
\end{equation}
then either
\begin{equation}
\left\langle \left.  \mathbf{xx}^{\prime}\right\vert \psi_{1}\psi
_{2}\right\rangle +\left\langle \left.  \mathbf{xx}^{\prime}\right\vert
\psi_{1}\psi_{2}\right\rangle \overset{ae}{=}0
\end{equation}
or
\begin{equation}
\lambda\left(  \mathbf{x}\right)  +\lambda\left(  \mathbf{x}^{\prime}\right)
\overset{ae}{=}0.
\end{equation}
As we are considering $\left\langle \left.  \mathbf{xx}^{\prime}\right\vert
\psi_{1}\psi_{2}\right\rangle +\left\langle \left.  \mathbf{xx}^{\prime
}\right\vert \psi_{1}\psi_{2}\right\rangle \overset{ae}{=}0$ for
$\mathbf{x,x}^{\prime}\in\Delta$ this gives
\begin{align}
& \lambda\left(  \mathbf{x}\right)  \overset{ae}{=}-\lambda\left(
\mathbf{x}^{\prime}\right)  \overset{ae}{=}-\lambda\left(  \mathbf{x}%
^{\prime\prime}\right)  \nonumber\\
\lambda\left(  \mathbf{x}^{\prime}\right)  \overset{ae}{=}-\lambda\left(
\mathbf{x}^{\prime\prime}\right)    & \Rightarrow\lambda\left(  \mathbf{x}%
^{\prime}\right)  \overset{ae}{=}-\lambda\left(  \mathbf{x}^{\prime}\right)
\end{align}
and hence%
\begin{equation}
\lambda\left(  \mathbf{x}\right)  \overset{ae}{=}0.
\end{equation}


One can observe that if $\left\vert a\right)  $ is normalized then
\begin{equation}
\left(  a_{k}\left(  \lambda\right)  \left\vert \left(  \widehat{K+V_{2}%
}\right)  \right.  a_{k}\left(  \lambda\right)  \right)  =\left(  a_{k}\left(
\lambda\right)  \left\vert \widehat{\Omega\left(  \eta_{k}\left(
\lambda\right)  \lambda\right)  }\right.  a_{k}\left(  \lambda\right)
\right)  =\eta_{k}\left(  \lambda\right)  \left(  a_{k}\left(  \lambda\right)
\left\vert \widehat{\Omega\left(  \lambda\right)  }\right.  a_{k}\left(
\lambda\right)  \right)
\end{equation}
is the universal energy (for the fixed density $\gamma$), and is related to
the value of the eigenvalue by
\begin{equation}
E_{U}\left(  a_{k}\left(  \lambda\right)  \right)  =\eta_{k}\left(
\lambda\right)  =\left(  a_{k}\left(  \lambda\right)  \left\vert \left(
\widehat{K+V_{2}}\right)  \right.  a_{k}\left(  \lambda\right)  \right)  ,
\end{equation}
as one also has
\begin{equation}
\left(  a_{k}\left(  \lambda\right)  \left\vert \widehat{\Omega\left(
\lambda\right)  }\right.  a_{k}\left(  \lambda\right)  \right)  =1.
\end{equation}
If $\left\{  \left\vert \overset{}{a_{k}}\right)  \right\}  $ are normalized
then we have
\begin{subequations}
\begin{align}
\left(  a_{k}\left(  \lambda\right)  \left\vert \widehat{\Omega\left(
\lambda\right)  }\right.  a_{l}\left(  \lambda\right)  \right)   &
=\delta_{kl}\\
\left(  a_{k}\left(  \lambda\right)  \left\vert \overset{}{a_{k}\left(
\lambda\right)  }\right.  \right)   &  =1
\end{align}
but
\end{subequations}
\begin{equation}
\left(  a_{k}\left(  \lambda\right)  \left\vert \overset{}{a_{l}\left(
\lambda\right)  }\right.  \right)  \neq0.
\end{equation}
Thus one would select the states with the lowest $\eta_{k}\left(
\lambda\right)  $ for the study of ground states.

If $\left\vert \overset{}{a}\right)  $ are solutions of $\left\{  \left(
\widehat{K+V_{2}}\right)  -\widehat{\Omega\left(  \lambda\right)  }\right\}
\left\vert \overset{}{a}\right)  =0$ for a given $\lambda,$ then $\left\vert
\overset{}{a}\right)  $ is a generalized eigenvector wrt $\widehat
{\Omega\left(  \lambda\right)  }$ with eigenvalue $1.$ If $\left\{  \left(
\widehat{K+V_{2}}\right)  -\eta\widehat{\Omega\left(  \lambda\right)
}\right\}  \left\vert \overset{}{a}\right)  =0$ then $\left\{  \left(
\widehat{K+V_{2}}\right)  -\widehat{\Omega\left(  \eta\lambda\right)
}\right\}  \left\vert \overset{}{a}\right)  =0,$ and
\begin{equation}
\left(  a\left\vert \widehat{\Omega\left(  \eta\lambda\right)  }\right.
a\right)  =\eta.
\end{equation}
If $\left\vert \overset{}{a}\right)  $ is normalized then
\begin{subequations}
\begin{align}
\left(  a\left\vert a\right.  \right)   &  =1\label{prop1}\\
\left(  a\left\vert \widehat{\Omega\left(  \lambda\right)  }\right.  a\right)
&  =1\label{prop2}\\
\left(  a\left\vert \left(  \widehat{K+V_{2}}\right)  \right.  a\right)   &
=\eta. \label{prop3}%
\end{align}
If we have degenerate states associated with an eigenvalue $\eta_{k},$ i.e.
\end{subequations}
\begin{equation}
\left(  a_{kj}\left\vert \left(  \widehat{K+V_{2}}\right)  \right.
a_{kj^{\prime}}\right)  =\eta_{k}\left(  a_{kj}\left\vert \widehat
{\Omega\left(  \lambda\right)  }\right.  a_{kj^{\prime}}\right)  =\eta
_{k}\delta_{jj^{\prime}}%
\end{equation}
then all real linear combinations $\left\vert a_{k}\left(  \mathbf{\beta
}\right)  \right)  $of $\left\{  \left.  \left\vert a_{kj}\right)  \right\vert
1\leq j\leq d_{k}\right\}  ,$ where
\begin{equation}
\left\vert a_{k\ast}\left(  \mathbf{\beta}\right)  \right)  =\sum_{1\leq j\leq
d_{k}}\beta_{j}\left\vert a_{kj}\right)  ,
\end{equation}
are eigenstates and satisfy the stationarity equation with the same $\eta
_{k}\lambda.$Thus one can have many $\left\vert \overset{}{a}\right)
^{\prime}s$ with \emph{the same} Lagrange multipliers (corresponding to a
singular Hessian $D_{aa}^{2}\mathcal{L}$ ), but if the constraint is
surjective we can only have one Lagrange multiplier vector for any $\left\vert
\overset{}{a}\right)  $ (thus the Hessian $D_{a\lambda}^{2}\mathcal{L}$ has
linearly independent rows)$.$ Thus we have
\begin{equation}
\left(  \sum_{1\leq j\leq d_{k}}\beta_{j}\left\vert a_{kj}\right)  \left\vert
\widehat{\Omega\left(  \lambda\right)  }\sum_{1\leq j\leq d_{k}}\beta
_{j}\left\vert a_{kj}\right)  \right.  \right)  =\sum_{1\leq j\leq d_{k}}%
\beta_{j}^{2}=1
\end{equation}
and for normalized states
\begin{equation}
\left(  \sum_{1\leq j\leq d_{k}}\beta_{j}\left\vert a_{kj}\right)  \left\vert
\sum_{1\leq j\leq d_{k}}\beta_{j}\left\vert a_{kj}\right)  \right.  \right)
=\sum_{1\leq j\leq d_{k}}\beta_{j}^{2}+2\sum_{1\leq j<j^{\prime}\leq d_{k}%
}\beta_{j}\beta_{j^{\prime}}\left(  \left.  \overset{}{a_{kj}}\right\vert
a_{kj^{\prime}}\right)  =1.
\end{equation}
This gives
\begin{equation}
\sum_{1\leq j<j^{\prime}\leq d_{k}}\beta_{j}\beta_{j^{\prime}}\left(  \left.
\overset{}{a_{kj}}\right\vert a_{kj^{\prime}}\right)  =0\quad\forall
\mathbf{\beta\backepsilon}\left\Vert \mathbf{\beta}\right\Vert =1,
\end{equation}
which implies that
\begin{equation}
\left(  \left.  \overset{}{a_{kj}}\right\vert a_{kj^{\prime}}\right)  =0\quad
j\neq j^{\prime}%
\end{equation}
and hence
\begin{equation}
\left(  \left.  \overset{}{a_{kj}}\right\vert a_{kj^{\prime}}\right)
=\delta_{jj^{\prime}}=\left(  a_{kj}\left\vert \widehat{\Omega\left(
\lambda\right)  }\right.  a_{kj^{\prime}}\right)  .
\end{equation}
$\lceil$ In the case $d_{k}=2$ this gives
\begin{align}
\left(  \left.  \overset{}{a_{k1}}\right\vert a_{k1}\right)   &  =\left(
\left.  \overset{}{a_{k2}}\right\vert a_{k2}\right)  =\left(  a_{k1}\left\vert
\widehat{\Omega\left(  \lambda\right)  }\right.  a_{k1}\right)  =\left(
a_{k2}\left\vert \widehat{\Omega\left(  \lambda\right)  }\right.
a_{k2}\right)  =1\nonumber\\
\left(  \left.  \overset{}{a_{k1}}\right\vert a_{k2}\right)   &  =\left(
a_{k1}\left\vert \widehat{\Omega\left(  \lambda\right)  }\right.
a_{k2}\right)  =0\rfloor.
\end{align}


Let $\left\vert \overset{}{a}\right)  $ be a constrained minimum i.e.
\begin{align}
\Xi\left(  a\right)   &  =\gamma,\nonumber\\
DE\left(  a\right)  \left(  h\right)   &  =\mathbf{0}\quad\forall
h\in\mathcal{T}_{a},\nonumber\\
D^{2}E\left(  a\right)  \left(  h\right)   &  \geq0\quad\forall h\in
\mathcal{T}_{a}.
\end{align}
Without loss of generality we can always replace $\Xi:\mathcal{B}_{2sa}\left(
\mathcal{H}^{N}\right)  \rightarrow\mathbb{R}^{r}$ by $\tilde{\Xi}%
:\mathcal{B}_{2sa}\left(  \mathcal{H}^{N}\right)  \rightarrow\mathbb{R}%
^{r^{\prime}}$ where we have excluded $r-r^{\prime}$ zero constraints $\Xi
_{j}\left(  a\right)  =0.$ However the solution can still be a non regular
point i.e. the vectors $\left\{  \left(  \overset{}{a}\right\vert
\widehat{\Phi_{j}^{\dagger}\Phi_{j}}\right\}  $ are linearly dependent or
$\left\{  \left(  \overset{}{a}\right\vert \mathbf{\Phi}\left(  \mathbf{x}%
\right)  ^{\dagger}\mathbf{\Phi}\left(  \mathbf{x}\right)  \right\}  $is not
basis. An example of this is when $\left(  a\right\vert $ is an IPS\ and
$\left\{  \left.  \widehat{\Phi_{j}^{\dagger}\Phi_{j}}\right\vert 1\leq j\leq
N\right\}  $refer to GSO's of occupied states.

If the solution is a mixed state $a^{2}$ then
\begin{align}
\left(  a\left\vert \left\{  \widehat{K+V_{2}}\right\}  \right.  a\right)   &
=\sum\alpha_{k}\left(  a_{k}\left\vert \left\{  \widehat{K+V_{2}}\right\}
\right.  a_{k}\right) \nonumber\\
\left(  a\left\vert \widehat{\Phi^{\dagger}\Phi}\right.  a\right)   &
=\sum\alpha_{k}\left(  a_{k}\left\vert \widehat{\Phi^{\dagger}\Phi}\right.
a_{k}\right)  =\gamma
\end{align}
so the pure states $\left\{  a_{k}\right\}  $ might not satisfy the
constraints individually and the mixed state cannot be replaced by any of
these pure states.

The tangent space has the following property
\begin{equation}
\mathcal{N}_{a}^{\bot}=\ker D\Xi\left(  a\right)  \supseteq\mathcal{T}_{a}.
\end{equation}
Hence
\begin{align}
DE\left(  a\right)   &  =n\left(  a,\lambda\right)  +\xi\left(  a,\eta
_{j}\right) \nonumber\\
&  =\sum_{1\leq j\leq r}\lambda_{1j}\nabla\Xi_{j}\left(  a\right)
+\sum_{1\leq j\leq r}\lambda_{2j}\nabla\Xi_{j}\left(  b\right)  ,\nonumber\\
&  =\left(  a\right\vert \Omega\left(  \lambda_{1}\right)  +\left(
b\right\vert \Omega\left(  \lambda_{2}\right)
\end{align}
where
\begin{align}
n\left(  a,\lambda_{1}\right)   &  \in\mathcal{N}_{a}\nonumber\\
\xi\left(  a,\lambda_{1}\right)   &  \in\mathcal{X}_{a}^{\bot}=\mathcal{N}%
_{a}^{\bot}\ominus\mathcal{T}_{a}\nonumber\\
\mathcal{T}_{a}  &  =\left(  \mathcal{N}_{a}\oplus\mathcal{X}_{a}\right)
^{\bot}\nonumber\\
\mathcal{T}_{a}^{\bot}  &  =\mathcal{N}_{a}\oplus\mathcal{X}_{a}.
\end{align}
The solution satisfies
\begin{equation}
\left\{  \left(  \widehat{K+V_{2}}\right)  -\widehat{\Omega\left(
\lambda\right)  }\right\}  \left\vert \overset{}{a}\right)  =\widehat
{\Omega\left(  \lambda\right)  }\left\vert \overset{}{b}\right)  .
\end{equation}
$\lceil$Choose a $\lambda$ such that $\widehat{\Omega\left(  \lambda\right)
}\left\vert \overset{}{a}\right)  \neq\mathbf{0}$ , $\lambda\neq\mathbf{0}$
and $\widehat{\Omega\left(  \lambda\right)  }$ is non singular$\mathbf{,}$
(this is always possible when $\widehat{\mathbf{\Phi}^{\dagger}\left(
\mathbf{x}\right)  \mathbf{\Phi}\left(  \mathbf{x}\right)  }\left\vert
\overset{}{a}\right)  \neq\mathbf{0),}$ then choose $\left\vert \overset{}%
{b}\right)  $ such that $\widehat{\mathbf{\Phi}^{\dagger}\left(
\mathbf{x}\right)  \mathbf{\Phi}\left(  \mathbf{x}\right)  }\left\vert
\overset{}{b}\right)  $ span $\mathcal{X}_{a}.$ One needs to prove a theorem
showing that a unique $\lambda$ exists for a given $\left\vert \overset{}%
{b}\right)  ,$ one also must prove that a $\left\vert \overset{}{b}\right)  $
exists$.\rfloor$%
\begin{equation}
\left(  \widehat{K+V_{2}}\right)  \left\vert \overset{}{a}\right)
=\widehat{\Omega\left(  \lambda\right)  }\left\{  \left\vert \overset{}%
{a}\right)  +\left\vert \overset{}{b}\right)  \right\}
\end{equation}


One can find a solutions
\begin{equation}
\left\{  \left(  \widehat{K+V_{2}}\right)  -\widehat{\Omega\left(
\lambda\right)  }\right\}  ^{G}\left\vert \overset{}{c\left(  \lambda\right)
}\right)  +\left\vert \overset{}{\tilde{a}}\right)  ,
\end{equation}
where
\begin{equation}
\left\{  \left(  \widehat{K+V_{2}}\right)  -\widehat{\Omega\left(
\lambda\right)  }\right\}  \left\vert \overset{}{\tilde{a}}\right)
=\mathbf{0.}%
\end{equation}


\subsection{Functional Definition}

One can define a density functional
\begin{equation}
\mathcal{F}_{U}\left(  \gamma\right)  =E_{U}\left(  a_{1\ast}\left(
\mathbf{\beta,}\gamma\right)  \right)  =\eta_{1}\left(  \gamma\right)
\end{equation}
such that
\begin{equation}
\frac{\delta\mathcal{F}_{U}\left(  \gamma\right)  }{\delta\gamma}=\frac
{\delta\mathcal{E}_{U}\left(  a_{1\ast}\left(  \mathbf{\beta,}\gamma\right)
\right)  }{\delta\gamma}=\frac{\delta\eta_{1}\left(  \gamma\right)  }%
{\delta\gamma}=\lambda\left(  \gamma\right)  .
\end{equation}
The functional $\mathcal{F}_{U}\left(  \gamma\right)  $ nor its derivatives
depend on $\mathbf{\beta}$ so one can just choose a fixed arbitrary
$\mathbf{\beta}$ to work with.

As $\left[  \widehat{\Phi_{j}^{\dagger}\Phi_{j}}\widehat{,\Phi_{k}^{\dagger
}\Phi_{k}}\right]  =0$ $\forall j,k$ we can find a conb $\left\{  \left.
\left\vert \sigma_{j}\right)  \right\vert 1\leq j\leq r_{N}^{2}\right\}  $ of
$\mathcal{B}_{2sa}\left(  \mathcal{H}^{N}\right)  $ such that
\begin{equation}
\widehat{\Phi_{j}^{\dagger}\Phi_{j}}=\sum_{1\leq k\leq r_{N}^{2}}\rho
_{jk}\left\vert \sigma_{k}\right)  \left(  \sigma_{k}\right\vert ;\quad1\leq
j\leq r
\end{equation}
giving
\begin{align}
\widehat{\Omega\left(  \mathbf{\lambda}\right)  }  &  =\sum_{\substack{1\leq
k\leq r_{N}^{2}\\1\leq j\leq r}}\lambda_{j}\rho_{jk}\left\vert \sigma
_{k}\right)  \left(  \sigma_{k}\right\vert =\sum_{1\leq k\leq r_{N}^{2}%
}\left(  \lambda\mathbf{\rho}\right)  _{k}\left\vert \sigma_{k}\right)
\left(  \sigma_{k}\right\vert \nonumber\\
&  =\sum_{1\leq k\leq r_{N}^{2}}\mathbf{\rho}^{t}\left(  \lambda\right)
_{k}\left\vert \sigma_{k}\right)  \left(  \sigma_{k}\right\vert =\sum_{1\leq
k\leq r_{N}^{2}}\xi_{k}\left\vert \sigma_{k}\right)  \left(  \sigma
_{k}\right\vert .
\end{align}
Clearly
\begin{align}
&  \mathbf{\rho}^{t}:\mathbb{R}^{r}\rightarrow\mathbb{R}^{r_{N}^{2}%
}\nonumber\\
&  \mathbf{\rho}^{t}\left(  \lambda\right)  =\xi.
\end{align}
If the linear map $\mathbf{\rho}^{t}$ is $1-1$ then $\mathbf{\rho}$ is onto.
Clearly $\mathbf{\rho}$ is the matrix representation of $\widehat{\Omega
}:L_{\infty}\left(  \mathbb{X}\right)  \rightarrow\mathcal{B}_{2sa\Phi}\left(
\mathcal{H}^{N}\right)  $, where this is the space of functions of $\left\{
\widehat{\Phi_{j}^{\dagger}\Phi_{j}}\right\}  .$

The constraint equations
\begin{equation}
\left(  a\left\vert \widehat{\Phi_{j}^{\dagger}\Phi_{j}}\right.  a\right)
=\gamma_{j}%
\end{equation}
can also be expressed as
\begin{align}
&  \left(  a\left\vert \sum_{1\leq k\leq r_{N}^{2}}\rho_{jk}\left\vert
\sigma_{k}\right)  \left(  \sigma_{k}\right\vert \right.  a\right)
=\sum_{1\leq k\leq r_{N}^{2}}\rho_{jk}\left(  \left.  \overset{}{a}\right\vert
\sigma_{k}\right)  \left(  \left.  \overset{}{\sigma}_{k}\right\vert a\right)
\nonumber\\
&  =\sum_{1\leq k\leq r_{N}^{2}}\rho_{jk}\left(  \left.  \overset{}%
{a}\right\vert \sigma_{k}\right)  ^{2}=\sum_{1\leq k\leq r_{N}^{2}}\rho
_{jk}a_{k}^{2}=\mathbf{\rho}\left(  a^{\left[  2\right]  }\right)  _{j}%
=\gamma_{j},
\end{align}
where
\begin{align}
\left\vert a^{\left[  2\right]  }\right)   &  =\sum_{1\leq k\leq r_{N}^{2}%
}a_{k}^{2}\left\vert \sigma_{k}\right) \nonumber\\
a_{k}  &  =\left(  \left.  \overset{}{\sigma}_{k}\right\vert a\right)  .
\end{align}
Showing that
\begin{align}
&  \mathbf{\rho}:\mathbb{R}^{r_{N}^{2}}\rightarrow\mathbb{R}^{r}\nonumber\\
&  \mathbf{\rho}\left(  a^{\left[  2\right]  }\right)  =\gamma.
\end{align}
The kernel of the linear map $\mathbf{\rho}$ has dimension $r_{N}^{2}$ $-r$ if
$\mathbf{\rho}$ is onto.

$\lceil$ The notation comes from
\begin{align}
\left(  |\right)  _{\mathcal{B}_{2sa}\left(  \mathcal{H}^{N}\right)  }  &
:\mathcal{B}_{2sa}\left(  \mathcal{H}^{N}\right)  \times\mathcal{B}%
_{2sa}\left(  \mathcal{H}^{N}\right)  \rightarrow\mathbb{R}\nonumber\\
\left(  |\right)  _{1}  &  :L_{\infty}\times L_{1}\rightarrow\mathbb{R}%
\end{align}
and
\begin{equation}
\left(  \lambda|\mathbf{\rho}\left(  a\right)  \right)  _{1}=\left(
\mathbf{\rho}^{t}\left(  \lambda\right)  |a\right)  _{\mathcal{B}_{2sa}\left(
\mathcal{H}^{N}\right)  }\rfloor
\end{equation}


\subsection{Newton Raphson Solution}

The Lagrangian can expanded up to second order as
\begin{align}
\mathcal{L}\left(  a+\delta a,\lambda+\delta\lambda,\gamma\right)   &
=\mathcal{L}\left(  a,\lambda,\gamma\right)  +\left(
\begin{array}
[c]{cc}%
D_{a}\mathcal{L}\left(  a,\lambda,\gamma\right)  & D_{\lambda}\mathcal{L}%
\left(  a,\lambda,\gamma\right)
\end{array}
\right)  \left(
\begin{array}
[c]{c}%
\delta a\\
\delta\lambda
\end{array}
\right) \nonumber\\
&  +\frac{1}{2}\left(
\begin{array}
[c]{cc}%
\delta a & \delta\lambda
\end{array}
\right)  \left(
\begin{array}
[c]{cc}%
D_{aa}^{2}\mathcal{L}\left(  a,\lambda,\gamma\right)  & D_{a\lambda}%
^{2}\mathcal{L}\left(  a,\lambda,\gamma\right) \\
D_{\lambda a}^{2}\mathcal{L}\left(  a,\lambda,\gamma\right)  & \mathbb{O}%
\end{array}
\right)  \left(
\begin{array}
[c]{c}%
\delta a\\
\delta\lambda
\end{array}
\right)  .
\end{align}
In matrix notation this becomes
\begin{align}
&  \mathcal{L}\left(  a,\lambda,\gamma\right)  +\left(
\begin{array}
[c]{cc}%
\cdots,2\left(  a\left\vert \left\{  \widehat{K+V_{2}}-\widehat{\Omega\left(
\lambda\right)  }\right\}  \right.  \Lambda_{jk}\right)  ,\cdots &
\cdots,\gamma_{j}-\left(  a\left\vert \widehat{\Phi_{j}^{\dagger}\Phi_{j}%
}\right.  a\right)  ,\cdots
\end{array}
\right)  \left(
\begin{array}
[c]{c}%
\delta\mathbf{a}\\
\delta\mathbf{\lambda}%
\end{array}
\right) \nonumber\\
&  +\frac{1}{2}\left(
\begin{array}
[c]{cc}%
\delta\mathbf{a} & \delta\mathbf{\lambda}%
\end{array}
\right)  \left(
\begin{array}
[c]{cc}%
2\left(  \Lambda_{lm}\left\vert \left\{  \widehat{K+V_{2}}-\widehat
{\Omega\left(  \lambda\right)  }\right\}  \right.  \Lambda_{jk}\right)  &
-2\left(  a\left\vert \widehat{\Phi_{j}^{\dagger}\Phi_{j}}\right.
\Lambda_{jk}\right) \\
-2\left(  \Lambda_{jk}\left\vert \widehat{\Phi_{j}^{\dagger}\Phi_{j}}\right.
a\right)  & \mathbb{O}%
\end{array}
\right)  \left(
\begin{array}
[c]{c}%
\delta\mathbf{a}\\
\delta\mathbf{\lambda}%
\end{array}
\right)  .
\end{align}
This can also be expressed as
\begin{align}
&  \mathcal{L}\left(  a+\delta a,\lambda+\delta\lambda,\gamma\right)
=\mathcal{L}\left(  a,\lambda,\gamma\right) \nonumber\\
&  +\left(
\begin{array}
[c]{cc}%
2\left(  a\left\vert \left\{  \widehat{K+V_{2}}-\widehat{\Omega\left(
\lambda\right)  }\right\}  \right.  \cdot\right)  & -\left(  a\left\vert
\widehat{\Omega\left(  \cdot\right)  }\right.  a\right)  +\left\langle
\cdot\left\vert \gamma\right.  \right\rangle
\end{array}
\right)  \left(
\begin{array}
[c]{c}%
\delta a\\
\delta\lambda
\end{array}
\right) \nonumber\\
&  +\left(
\begin{array}
[c]{cc}%
\delta a & \delta\lambda
\end{array}
\right)  \left(
\begin{array}
[c]{cc}%
\left\{  \widehat{K+V_{2}}-\widehat{\Omega\left(  \lambda\right)  }\right\}  &
\left(  a\left\vert \left\{  \widehat{K+V_{2}}-\widehat{\Omega\left(
\cdot\right)  }\right\}  \right.  \cdot\right) \\
\left(  a\left\vert \left\{  \widehat{K+V_{2}}-\widehat{\Omega\left(
\cdot\right)  }\right\}  \right.  \cdot\right)  & \mathbb{O}%
\end{array}
\right)  \left(
\begin{array}
[c]{c}%
\delta a\\
\delta\lambda
\end{array}
\right)  .
\end{align}
The stationary points of the second order approximation are then given by the
solutions of
\begin{align}
&  \left(
\begin{array}
[c]{cc}%
\cdots,2\left(  a\left\vert \left\{  \widehat{K+V_{2}}-\widehat{\Omega\left(
\lambda\right)  }\right\}  \right.  \Lambda_{jk}\right)  ,\cdots &
\cdots,\gamma_{j}-\left(  a\left\vert \widehat{\Phi_{j}^{\dagger}\Phi_{j}%
}\right.  a\right)  ,\cdots
\end{array}
\right)  +\nonumber\\
2  &  \left(
\begin{array}
[c]{cc}%
\delta\mathbf{a} & \delta\mathbf{\lambda}%
\end{array}
\right)  \left(
\begin{array}
[c]{cc}%
\begin{array}
[c]{c}%
\vdots\\
\cdots\left(  \Lambda_{lm}\left\vert \left\{  \widehat{K+V_{2}}-\widehat
{\Omega\left(  \lambda\right)  }\right\}  \right.  \Lambda_{jk}\right)
\cdots\\
\vdots
\end{array}
&
\begin{array}
[c]{c}%
\vdots\\
\cdots-\left(  a\left\vert \widehat{\Phi_{j}^{\dagger}\Phi_{j}}\right.
\Lambda_{jk}\right)  \cdots\\
\vdots
\end{array}
\\%
\begin{array}
[c]{c}%
\vdots\\
\cdots-\left(  \Lambda_{jk}\left\vert \widehat{\Phi_{j}^{\dagger}\Phi_{j}%
}\right.  a\right)  \cdots\\
\vdots
\end{array}
& \mathbb{O}%
\end{array}
\right)  =\left(
\begin{array}
[c]{cc}%
\mathbf{0} & \mathbf{0}%
\end{array}
\right)  \mathbf{,}%
\end{align}
which gives%
\begin{align}
&  \delta\mathbf{a}^{t}\left(
\begin{array}
[c]{c}%
\vdots\\
\cdots\left(  \Lambda_{lm}\left\vert \left\{  \widehat{K+V_{2}}-\widehat
{\Omega\left(  \lambda\right)  }\right\}  \right.  \Lambda_{jk}\right)
\cdots\\
\vdots
\end{array}
\right)  +\delta\mathbf{\lambda}^{t}\left(
\begin{array}
[c]{c}%
\vdots\\
-\cdots\left(  \Lambda_{jk}\left\vert \widehat{\Phi_{j}^{\dagger}\Phi_{j}%
}\right.  a\right)  \cdots\\
\vdots
\end{array}
\right) \nonumber\\
&  \qquad\qquad\qquad\qquad\qquad\qquad\qquad\qquad\qquad\left.  =\right.
\left(  \cdots,\left(  a\left\vert \left\{  \widehat{K+V_{2}}-\widehat
{\Omega\left(  \lambda\right)  }\right\}  \right.  \Lambda_{jk}\right)
,\cdots\right) \nonumber\\
&  2\delta\mathbf{a}^{t}\left(
\begin{array}
[c]{c}%
\vdots\\
\cdots-\left(  a\left\vert \widehat{\Phi_{j}^{\dagger}\Phi_{j}}\right.
\Lambda_{jk}\right)  \cdots\\
\vdots
\end{array}
\right)  =\mathbf{\gamma-}\left(  \cdots,\left(  a\left\vert \widehat{\Phi
_{j}^{\dagger}\Phi_{j}}\right.  a\right)  ,\cdots\right)  \mathbf{.}%
\end{align}
Hence
\begin{align}
\delta\mathbf{\lambda}^{t}\left(
\begin{array}
[c]{c}%
\vdots\\
-\cdots\left(  \Lambda_{jk}\left\vert \widehat{\Phi_{j}^{\dagger}\Phi_{j}%
}\right.  a\right)  \cdots\\
\vdots
\end{array}
\right)   &  =\left(  \cdots,\left(  a\left\vert \left\{  \widehat{K+V_{2}%
}-\widehat{\Omega\left(  \lambda\right)  }\right\}  \right.  \Lambda
_{jk}\right)  ,\cdots\right) \nonumber\\
&  -\delta\mathbf{a}^{t}\left(
\begin{array}
[c]{c}%
\vdots\\
\cdots\left(  \Lambda_{lm}\left\vert \left\{  \widehat{K+V_{2}}-\widehat
{\Omega\left(  \lambda\right)  }\right\}  \right.  \Lambda_{jk}\right)
\cdots\\
\vdots
\end{array}
\right)
\end{align}
giving
\begin{align}
\delta\mathbf{\lambda}^{t}  &  =\left(
\begin{array}
[c]{c}%
\vdots\\
-\cdots\left(  \Lambda_{jk}\left\vert \widehat{\Phi_{j}^{\dagger}\Phi_{j}%
}\right.  a\right)  \cdots\\
\vdots
\end{array}
\right)  ^{G}\left\{  \left(  \cdots,\left(  a\left\vert \left\{
\widehat{K+V_{2}}-\widehat{\Omega\left(  \lambda\right)  }\right\}  \right.
\Lambda_{jk}\right)  ,\cdots\right)  \right. \nonumber\\
&  \qquad\qquad\qquad\qquad\left.  -\,\delta\mathbf{a}^{t}\left(
\begin{array}
[c]{c}%
\vdots\\
\cdots\left(  \Lambda_{lm}\left\vert \left\{  \widehat{K+V_{2}}-\widehat
{\Omega\left(  \lambda\right)  }\right\}  \right.  \Lambda_{jk}\right)
\cdots\\
\vdots
\end{array}
\right)  \right\}  .
\end{align}
This leads to
\begin{align}
&  \left\{  \delta\mathbf{\lambda}^{t}\left(
\begin{array}
[c]{c}%
\vdots\\
-\cdots\left(  \Lambda_{jk}\left\vert \widehat{\Phi_{j}^{\dagger}\Phi_{j}%
}\right.  a\right)  \cdots\\
\vdots
\end{array}
\right)  -\left(  \cdots,\left(  a\left\vert \left\{  \widehat{K+V_{2}%
}-\widehat{\Omega\left(  \lambda\right)  }\right\}  \right.  \Lambda
_{jk}\right)  ,\cdots\right)  \right\}  \times\nonumber\\
&  \left(
\begin{array}
[c]{c}%
\vdots\\
\cdots\left(  \Lambda_{lm}\left\vert \left\{  \widehat{K+V_{2}}-\widehat
{\Omega\left(  \lambda\right)  }\right\}  \right.  \Lambda_{jk}\right)
\cdots\\
\vdots
\end{array}
\right)  ^{G}=\delta\mathbf{a}^{t}.
\end{align}
Thus
\begin{align}
&  2\delta\mathbf{a}^{t}\left\{  \delta\mathbf{\lambda}^{t}\left(
\begin{array}
[c]{c}%
\vdots\\
-\cdots\left(  \Lambda_{jk}\left\vert \widehat{\Phi_{j}^{\dagger}\Phi_{j}%
}\right.  a\right)  \cdots\\
\vdots
\end{array}
\right)  -\left(  \cdots,\left(  a\left\vert \left\{  \widehat{K+V_{2}%
}-\widehat{\Omega\left(  \lambda\right)  }\right\}  \right.  \Lambda
_{jk}\right)  ,\cdots\right)  \right\}  \times\nonumber\\
&  \left(
\begin{array}
[c]{c}%
\vdots\\
\cdots\left(  \Lambda_{lm}\left\vert \left\{  \widehat{K+V_{2}}-\widehat
{\Omega\left(  \lambda\right)  }\right\}  \right.  \Lambda_{jk}\right)
\cdots\\
\vdots
\end{array}
\right)  ^{G}\left(
\begin{array}
[c]{c}%
\vdots\\
\cdots-\left(  a\left\vert \widehat{\Phi_{j}^{\dagger}\Phi_{j}}\right.
\Lambda_{jk}\right)  \cdots\\
\vdots
\end{array}
\right) \nonumber\\
&  =\mathbf{\gamma-}\left(  \cdots,\left(  a\left\vert \widehat{\Phi
_{j}^{\dagger}\Phi_{j}}\right.  a\right)  ,\cdots\right)
\end{align}
and
\begin{align}
&  \delta\mathbf{\lambda}^{t}\nonumber\\
&  =\left\{  \mathbf{\gamma-}\left(  \cdots,\left(  a\left\vert \widehat
{\Phi_{j}^{\dagger}\Phi_{j}}\right.  a\right)  ,\cdots\right)  \right\}
\times\nonumber\\
&  \left\{  \frac{1}{2}\left\{  \left(
\begin{array}
[c]{c}%
\vdots\\
\cdots\left(  \Lambda_{lm}\left\vert \left\{  \widehat{K+V_{2}}-\widehat
{\Omega\left(  \lambda\right)  }\right\}  \right.  \Lambda_{jk}\right)
\cdots\\
\vdots
\end{array}
\right)  ^{G}\left(
\begin{array}
[c]{c}%
\vdots\\
\cdots-\left(  a\left\vert \widehat{\Phi_{j}^{\dagger}\Phi_{j}}\right.
\Lambda_{jk}\right)  \cdots\\
\vdots
\end{array}
\right)  \right\}  \right.  ^{G}\nonumber\\
&  \left.  \left(  a\left\vert \left\{  \widehat{K+V_{2}}-\widehat
{\Omega\left(  \lambda\right)  }\right\}  \right.  \delta\mathbf{a}\right)
\right\}  \left(
\begin{array}
[c]{c}%
\vdots\\
-\cdots\left(  \delta\mathbf{a}\left\vert \widehat{\Phi_{j}^{\dagger}\Phi_{j}%
}\right.  a\right)  \cdots\\
\vdots
\end{array}
\right)  ^{G}%
\end{align}
i.e.
\begin{align}
\left(
\begin{array}
[c]{cc}%
\delta a\left(  \kappa_{a}\right)  & \delta\lambda\left(  \kappa_{\lambda
}\right)
\end{array}
\right)   &  =-\left(
\begin{array}
[c]{cc}%
D_{a}\mathcal{L}\left(  a,\lambda,\gamma\right)  & D_{\lambda}\mathcal{L}%
\left(  a,\lambda,\gamma\right)
\end{array}
\right)  \left(
\begin{array}
[c]{cc}%
D_{aa}^{2}\mathcal{L}\left(  a,\lambda,\gamma\right)  & D_{a\lambda}%
^{2}\mathcal{L}\left(  a,\lambda,\gamma\right) \\
D_{\lambda a}^{2}\mathcal{L}\left(  a,\lambda,\gamma\right)  & \mathbb{O}%
\end{array}
\right)  ^{G}\nonumber\\
&  \qquad\qquad\qquad\qquad\qquad\qquad\qquad\qquad\qquad\qquad\qquad
\qquad+\left(
\begin{array}
[c]{cc}%
\kappa_{a} & \kappa_{\lambda}%
\end{array}
\right)  ,
\end{align}
where
\begin{equation}
\left(
\begin{array}
[c]{cc}%
\kappa_{a} & \kappa_{\lambda}%
\end{array}
\right)  \in\ker D^{2}\mathcal{L}\left(  a,\lambda,\gamma\right)  .
\end{equation}
Recall
\begin{subequations}
\begin{align}
D_{a}\mathcal{L}\left(  a,\lambda,\gamma\right)  \left(  \delta a\right)   &
=2\left(  a\left\vert \left\{  \widehat{K+V_{2}}-\widehat{\Omega\left(
\lambda\right)  }\right\}  \right.  \delta a\right) \\
D_{\lambda}\mathcal{L}\left(  a,\lambda,\gamma\right)  \left(  \delta
\lambda\right)   &  =\left\langle \delta\lambda\left\vert \gamma\right.
\right\rangle -\left(  a\left\vert \widehat{\Omega\left(  \delta
\lambda\right)  }\right.  a\right) \\
D_{aa}^{2}\mathcal{L}\left(  a,\lambda,\gamma\right)  \left(  \delta a,\delta
a\right)   &  =\left(  \delta a\left\vert \left\{  \widehat{K+V_{2}}%
-\widehat{\Omega\left(  \lambda\right)  }\right\}  \right.  \delta a\right) \\
D_{a\lambda}^{2}\mathcal{L}\left(  a,\lambda,\gamma\right)  \left(  \delta
a,\delta\lambda\right)   &  =2\left(  a\left\vert \widehat{\Omega\left(
\delta\lambda\right)  }\right.  \delta a\right)
\end{align}
We can applying partitioning to the quadratic approximation equations for for
stationary points to obtain
\end{subequations}
\begin{equation}
\left(
\begin{array}
[c]{cc}%
D_{aa}^{2}\mathcal{L}\left(  a,\lambda,\gamma\right)  & D_{a\lambda}%
^{2}\mathcal{L}\left(  a,\lambda,\gamma\right) \\
D_{\lambda a}^{2}\mathcal{L}\left(  a,\lambda,\gamma\right)  & \mathbb{O}%
\end{array}
\right)  \left(
\begin{array}
[c]{c}%
\delta a\\
\delta\lambda
\end{array}
\right)  =-\left(
\begin{array}
[c]{c}%
D_{a}\mathcal{L}\left(  a,\lambda,\gamma\right) \\
D_{\lambda}\mathcal{L}\left(  a,\lambda,\gamma\right)
\end{array}
\right)  \label{ME1}%
\end{equation}
that is equivalent to the coupled equations
\begin{align}
D_{a\lambda}^{2}\mathcal{L}\left(  a,\lambda,\gamma\right)  \delta\lambda &
=-\left\{  \overset{}{D_{a}\mathcal{L}\left(  a,\lambda,\gamma\right)
+D_{aa}^{2}\mathcal{L}\left(  a,\lambda,\gamma\right)  \delta a}\right\}
\label{x1eq}\\
D_{\lambda a}^{2}\mathcal{L}\left(  a,\lambda,\gamma\right)  \delta a  &
=-D_{\lambda}\mathcal{L}\left(  a,\lambda,\gamma\right)  .
\end{align}
Hence
\begin{equation}
\delta\lambda=-\left[  D_{a\lambda}^{2}\mathcal{L}\left(  a,\lambda
,\gamma\right)  \right]  ^{G}\left\{  \overset{}{D_{a}\mathcal{L}\left(
a,\lambda,\gamma\right)  -D_{aa}^{2}\mathcal{L}\left(  a,\lambda
,\gamma\right)  \left[  D_{\lambda a}^{2}\mathcal{L}\left(  a,\lambda
,\gamma\right)  \right]  ^{G}D_{\lambda}\mathcal{L}\left(  a,\lambda
,\gamma\right)  }\right\}  .
\end{equation}
If $D_{a}\mathcal{L}\left(  a,\lambda,\gamma\right)  =0$ then this becomes
\begin{equation}
\delta\lambda=\,\overset{}{\left[  D_{a\lambda}^{2}\mathcal{L}\left(
a,\lambda,\gamma\right)  \right]  ^{G}D_{aa}^{2}\mathcal{L}\left(
a,\lambda,\gamma\right)  \left[  D_{\lambda a}^{2}\mathcal{L}\left(
a,\lambda,\gamma\right)  \right]  ^{G}D_{\lambda}\mathcal{L}\left(
a,\lambda,\gamma\right)  }.
\end{equation}
One should observe that
\begin{align}
D_{a\lambda}^{2}\mathcal{L}\left(  a,\lambda,\gamma\right)   &  =\left[
D_{\lambda a}^{2}\mathcal{L}\left(  a,\lambda,\gamma\right)  \right]  ^{t}\\
\left[  D_{a\lambda}^{2}\mathcal{L}\left(  a,\lambda,\gamma\right)  \right]
^{G}\left[  D_{a\lambda}^{2}\mathcal{L}\left(  a,\lambda,\gamma\right)
\right]   &  =P_{\ker D_{a\lambda}^{2}\mathcal{L}\left(  a,\lambda
,\gamma\right)  ^{\bot}}\\
\left[  D_{a\lambda}^{2}\mathcal{L}\left(  a,\lambda,\gamma\right)  \right]
\left[  D_{a\lambda}^{2}\mathcal{L}\left(  a,\lambda,\gamma\right)  \right]
^{G}  &  =P_{\operatorname{Ran}D_{a\lambda}^{2}\mathcal{L}\left(
a,\lambda,\gamma\right)  }%
\end{align}
Thus the procedure is to first $a$ for a fixed $\lambda$ such that
$D_{a}\mathcal{L}\left(  a,\lambda,\gamma\right)  =0$ then find $\delta
\lambda,$ change $\lambda\rightarrow$ $\lambda+\delta\lambda$ and find a new
$a$ and so on.

\subsection{Examples}

Let $r=1$ then eq. $\left(  \ref{EnStation}\right)  $
\begin{equation}
\left\{  \left(  \widehat{K+V_{2}}\right)  -\lambda\widehat{\Phi_{1}^{\dagger
}\Phi_{1}}\right\}  \left\vert \overset{}{a}\right)  =0
\end{equation}
leading to eigenvalue equation
\begin{equation}
\left(  \widehat{K+V_{2}}\right)  \left\vert \overset{}{a}\right)
=\lambda\widehat{\Phi_{1}^{\dagger}\Phi_{1}}\left\vert \overset{}{a}\right)  .
\label{e1}%
\end{equation}
We can choose normalized eigenvectors $\left\vert b_{p}\right)  $ of
$\widehat{\Phi_{1}^{\dagger}\Phi_{1}}^{-\frac{1}{2}}\left(  \widehat{K+V_{2}%
}\right)  \widehat{\Phi_{1}^{\dagger}\Phi_{1}}^{-\frac{1}{2}}$ such that
\begin{subequations}
\begin{equation}
\left(  b_{p}\left\vert \widehat{\Phi_{1}^{\dagger}\Phi_{1}}^{-\frac{1}{2}%
}\left(  \widehat{K+V_{2}}\right)  \widehat{\Phi_{1}^{\dagger}\Phi_{1}%
}^{-\frac{1}{2}}b_{q}\right.  \right)  =\lambda_{p}\delta_{pq}%
\end{equation}
and
\end{subequations}
\begin{equation}
\left(  b_{p}\left\vert b_{q}\right.  \right)  =\left(  \tilde{a}%
_{p}\left\vert \widehat{\Phi_{1}^{\dagger}\Phi_{1}}\tilde{a}_{q}\right.
\right)  =\delta_{pq},
\end{equation}
where
\begin{equation}
\left\vert \overset{}{\tilde{a}_{p}}\right)  =\widehat{\Phi_{1}^{\dagger}%
\Phi_{1}}^{-\frac{1}{2}}\left\vert \overset{}{b_{p}}\right)  .
\end{equation}
This gives
\begin{equation}
\left(  \tilde{a}_{p}\left\vert \left(  \widehat{K+V_{2}}\right)  \tilde
{a}_{q}\right.  \right)  =\lambda_{p}\left(  \tilde{a}_{p}\left\vert
\widehat{\Phi_{1}^{\dagger}\Phi_{1}}\tilde{a}_{q}\right.  \right)
=\lambda_{p}\delta_{pq}%
\end{equation}
and
\begin{equation}
\left(  \widehat{K+V_{2}}\right)  \left\vert \overset{}{\tilde{a}_{p}}\right)
=\lambda_{p}\widehat{\Phi_{1}^{\dagger}\Phi_{1}}\left\vert \overset{}%
{\tilde{a}_{p}}\right)  .
\end{equation}
We can then define
\begin{equation}
\left\vert \overset{}{a_{p}}\right)  =\gamma^{\frac{1}{2}}\left\vert
\overset{}{\tilde{a}_{p}}\right)  ,
\end{equation}
so that
\begin{equation}
\left(  \widehat{K+V_{2}}\right)  \left\vert \overset{}{a_{p}}\right)
=\lambda_{p}\widehat{\Phi_{1}^{\dagger}\Phi_{1}}\left\vert \overset{}{a_{p}%
}\right)
\end{equation}
and
\begin{equation}
\left(  a_{p}\left\vert \widehat{\Phi_{1}^{\dagger}\Phi_{1}}a_{p}\right.
\right)  =\gamma.
\end{equation}
Hence
\begin{equation}
\left(  \overset{}{a_{p}}\left\vert \left(  \widehat{K+V_{2}}\right)  \right.
\overset{}{a_{p}}\right)  =\lambda_{p}\left(  \overset{}{a_{p}}\left\vert
\widehat{\Phi_{1}^{\dagger}\Phi_{1}}\right.  \overset{}{a_{p}}\right)
=\lambda_{p}\gamma
\end{equation}
and
\begin{equation}
\left(  \overset{}{a_{p}}\left\vert \overset{}{a_{p}}\right.  \right)
=\gamma\left(  \overset{}{b_{p}}\left\vert \widehat{\Phi_{1}^{\dagger}\Phi
_{1}}^{-1}\overset{}{b_{p}}\right.  \right)  ,
\end{equation}
thus the energy, $E_{U}\left(  a_{p}\right)  ,$ of the state $\left\vert
\overset{}{a_{p}}\right)  $ is
\begin{equation}
E_{U}\left(  a_{p}\right)  =\lambda_{p}\left(  \overset{}{b_{p}}\left\vert
\widehat{\Phi_{1}^{\dagger}\Phi_{1}}^{-1}\overset{}{b_{p}}\right.  \right)
^{-1},
\end{equation}
where $\left\vert b_{p}\right)  $ is \emph{independent }of $\gamma$.

The eigenvalues of eq. $\left(  \ref{e1}\right)  $ are solutions to
\begin{equation}
\det\left\{  \left(  \widehat{K+V_{2}}\right)  -\lambda_{p}\widehat{\Phi
_{1}^{\dagger}\Phi_{1}}\right\}  =0
\end{equation}
and the eigenvectors associated with the eigenvalue $\lambda_{p}$ are the
vectors $\left\{  \left\vert \overset{}{a_{p1}}\right)  ,\ldots,\left\vert
\overset{}{a_{pm_{p}}}\right)  \right\}  $ that span the $\ker\left\{  \left(
\widehat{K+V_{2}}\right)  -\lambda_{p}\widehat{\Phi_{1}^{\dagger}\Phi_{1}%
}\right\}  $. These (in general non orthonormal) vectors have the property
\begin{equation}
\left(  \overset{}{a_{pq}}\left\vert \left(  \widehat{K+V_{2}}\right)
\right.  \left\vert \overset{}{a_{pq}}\right)  \right)  =\lambda_{p1}\left(
\overset{}{a_{pq}}\left\vert \widehat{\Phi_{1}^{\dagger}\Phi_{1}}\right.
\left\vert \overset{}{a_{pq}}\right)  \right)  =\lambda_{p}\gamma,
\end{equation}
showing that if $\left(  \left.  \overset{}{a_{1p}}\right\vert \overset
{}{a_{1p}}\right)  $ $=\gamma=1$ then $\lambda_{1}$ is the energy.

In the case $m_{p}>1$ we obtain a possible space of solutions
\begin{equation}
\sum_{1\leq q\leq m_{p}}\alpha_{q}\left\vert \overset{}{a_{pq}}\right)  .
\end{equation}
In order that these are constrained solutions they must also satisfy
\begin{equation}
\sum_{1\leq q,q^{\prime}\leq m}\alpha_{q}\alpha_{q^{\prime}}\left(  \overset
{}{a_{pq}}\left\vert \widehat{\Phi_{1}^{\dagger}\Phi_{1}}\right.  \overset
{}{a_{pq^{\prime}}}\right)  =\gamma. \label{cond1}%
\end{equation}
The vectors $\left\{  \left\vert \overset{}{a_{p1}}\right)  ,\ldots,\left\vert
\overset{}{a_{pm_{p}}}\right)  \right\}  $ have the property that
\begin{equation}
\left(  \overset{}{a_{pq}}\left\vert \widehat{\Phi_{1}^{\dagger}\Phi_{1}%
}\right.  \overset{}{a_{1q^{\prime}}}\right)  =\gamma\delta_{qq^{\prime}}%
\end{equation}
and
\begin{equation}
\left(  \widehat{K+V_{2}}\right)  \left\vert \overset{}{a_{pq}}\right)
=\lambda_{p}\left\vert \overset{}{a_{pq}}\right)  ,
\end{equation}
so that eq. $\left(  \ref{cond1}\right)  $ becomes
\begin{equation}
\sum_{1\leq q\leq m_{p}}\alpha_{q}^{2}=\gamma.
\end{equation}
Thus one has the solutions $\left\{  \left.  \left\vert a_{p}\left(
\mathbf{\alpha}\left(  \gamma\right)  \right)  \right)  ;\lambda_{p}\left(
\gamma\right)  ;\right\vert \,\left\vert a_{1}\left(  \mathbf{\alpha}%
,\gamma\right)  \right)  \in\mathbb{R}^{r_{N}^{2}};\left\Vert \mathbf{\alpha
}\right\Vert =\gamma,\mathbf{\alpha\in}\mathbb{R}^{m_{p}},\gamma\mathbf{\in
}\mathbb{R}^{1}\right\}  .$ One can introduce a foliation of $\mathbb{R}^{m}$
by the subspaces $\left\{  \mathbb{R}^{1},\cdots,\mathbb{R}^{m}\right\}  $ and
express changes in $\gamma$ by the scaling of the $\left\{  \alpha
_{p}\right\}  $ i.e.
\begin{align}
\mathbf{\alpha}  &  \rightarrow\eta\mathbf{\alpha}\nonumber\\
\gamma &  \rightarrow\eta\gamma
\end{align}
and then follow a paths of solutions $\left(  \left\vert a_{p}\left(
\mathbf{\alpha}\left(  \eta\gamma\right)  \right)  \right)  ,\lambda
_{p}\left(  \eta\gamma\right)  \right)  $ for on each foliation a fixed
initial $\gamma$. In this case with $r=1,$ $\left(  \left\vert a_{p}\left(
\mathbf{\alpha}\left(  \eta\gamma\right)  \right)  \right)  ,\lambda
_{p}\left(  \eta\gamma\right)  \right)  \equiv\left(  \left\vert \eta
^{\frac{1}{2}}a_{p}\left(  \mathbf{\alpha}\left(  \gamma\right)  \right)
\right)  ,\lambda_{p}\left(  \gamma\right)  \right)  .$

One repeats this for $\lambda_{2},\cdots,\lambda_{q}$ and thus finds all the
solutions. If $\widehat{\Phi_{1}^{\dagger}\Phi_{1}}$ = $I$ and $\gamma=1$then
this equation gives the stationary $N$-electron states.

\subsubsection{$\lceil$Generalized Eigenvalue Problem}

We want to find solutions to
\begin{equation}
H\left\vert v\right\rangle =\lambda S\left\vert v\right\rangle ,
\end{equation}
where $H=H^{\dagger}$ and $S>0.$ This equation can be rewritten as
\begin{align}
HS^{-\frac{1}{2}}S^{\frac{1}{2}}\left\vert v\right\rangle  &  =\lambda
SS^{-\frac{1}{2}}S^{\frac{1}{2}}\left\vert v\right\rangle \Rightarrow
\nonumber\\
S^{-\frac{1}{2}}HS^{-\frac{1}{2}}S^{\frac{1}{2}}\left\vert v\right\rangle  &
=\lambda S^{\frac{1}{2}}\left\vert v\right\rangle ,
\end{align}
where $S^{-\frac{1}{2}}HS^{-\frac{1}{2}}=\left(  S^{-\frac{1}{2}}HS^{-\frac
{1}{2}}\right)  ^{\dagger}$. Thus we can choose normalized eigenvectors
$\left\vert \varphi_{j}\right\rangle $ of $S^{-\frac{1}{2}}HS^{-\frac{1}{2}}$
such that
\begin{subequations}
\begin{align}
&  \left\langle \varphi_{j}\left\vert S^{-\frac{1}{2}}HS^{-\frac{1}{2}}%
\varphi_{k}\right.  \right\rangle =\lambda_{j}\delta_{jk},\nonumber\\
&  \left\langle \varphi_{j}\left\vert \varphi_{k}\right.  \right\rangle
=\delta_{jk},\nonumber\\
&  \left\vert \nu_{j}\right\rangle =S^{-\frac{1}{2}}\left\vert \varphi
_{j}\right\rangle \Leftrightarrow S^{\frac{1}{2}}\left\vert \nu_{j}%
\right\rangle =\left\vert \varphi_{j}\right\rangle \nonumber\\
&  \Rightarrow\left\langle \nu_{j}\left\vert S\nu_{k}\right.  \right\rangle
=\delta_{jk},
\end{align}
so that
\end{subequations}
\begin{equation}
\left\langle \nu_{j}\left\vert H\nu_{k}\right.  \right\rangle =\lambda
_{j}\left\langle \nu_{j}\left\vert S\nu_{k}\right.  \right\rangle =\lambda
_{j}\delta_{jk}.
\end{equation}
If $S$ is singular the eigenvalue equation can only have solutions that span
the space if $\ker S$ = $\ker H.$ If $S=S^{\dagger}$ but is not positive then
\begin{align}
S  &  =\sum_{j}\beta_{j}\left\vert \omega_{j}\right\rangle \left\langle
\omega_{j}\right\vert \\
S\left(  \mathbf{\theta}\right)  ^{\frac{1}{2}}  &  =\sum_{j}\beta_{j}%
^{\frac{1}{2}}e^{i\theta_{j}}\left\vert \omega_{j}\right\rangle \left\langle
\omega_{j}\right\vert
\end{align}
so that $S^{\frac{1}{2}}\neq\left(  S\left(  \mathbf{\theta}\right)
^{\frac{1}{2}}\right)  ^{\dagger}$ but
\begin{align}
&  S\left(  \mathbf{\theta}\right)  ^{\frac{1}{2}}\left(  S\left(
\mathbf{\theta}\right)  ^{\frac{1}{2}}\right)  ^{\dagger}=\left(  S\left(
\mathbf{\theta}\right)  ^{\frac{1}{2}}\right)  ^{\dagger}S\left(
\mathbf{\theta}\right)  ^{\frac{1}{2}}=S\\
&  SS\left(  \mathbf{\theta}\right)  ^{-\frac{1}{2}}=\left(  S\left(
\mathbf{\theta}\right)  ^{\frac{1}{2}}\right)  ^{\dagger}.
\end{align}
Hence
\begin{align}
&  HS\left(  \mathbf{\theta}\right)  ^{-\frac{1}{2}}S\left(  \mathbf{\theta
}\right)  ^{\frac{1}{2}}\left\vert v\right\rangle =\lambda SS\left(
\mathbf{\theta}\right)  ^{-\frac{1}{2}}S\left(  \mathbf{\theta}\right)
^{\frac{1}{2}}\left\vert v\right\rangle \Rightarrow\nonumber\\
&  \left(  S\left(  \mathbf{\theta}\right)  ^{-\frac{1}{2}}\right)  ^{\dagger
}HS\left(  \mathbf{\theta}\right)  ^{-\frac{1}{2}}S\left(  \mathbf{\theta
}\right)  ^{\frac{1}{2}}\left\vert v\right\rangle =\lambda S\left(
\mathbf{\theta}\right)  ^{\frac{1}{2}}\left\vert v\right\rangle ,
\end{align}
so $\left(  S\left(  \mathbf{\theta}\right)  ^{-\frac{1}{2}}\right)
^{\dagger}HS\left(  \mathbf{\theta}\right)  ^{-\frac{1}{2}}$ is self adjoint
and has normalized eigenvectors
\begin{equation}
\left\vert \psi_{j}\right\rangle =S\left(  \mathbf{\theta}\right)  ^{\frac
{1}{2}}\left\vert v_{j}\right\rangle ,
\end{equation}
such that
\begin{equation}
\left\langle \psi_{j}\left\vert \psi_{k}\right.  \right\rangle =\left\langle
\nu_{j}\left\vert \left(  S\left(  \mathbf{\theta}\right)  ^{\frac{1}{2}%
}\right)  ^{\dagger}S\left(  \mathbf{\theta}\right)  ^{\frac{1}{2}}\nu
_{k}\right.  \right\rangle =\left\langle \nu_{j}\left\vert S\nu_{k}\right.
\right\rangle =\delta_{jk}%
\end{equation}
and
\begin{equation}
\left\langle \psi_{j}\left\vert \left(  S\left(  \mathbf{\theta}\right)
^{-\frac{1}{2}}\right)  ^{\dagger}HS\left(  \mathbf{\theta}\right)
^{-\frac{1}{2}}\psi_{k}\right.  \right\rangle =\left\langle \nu_{j}\left\vert
H\nu_{k}\right.  \right\rangle =\lambda_{j}\delta_{jk}.
\end{equation}
We can generalize this in various ways and obtain the following theorems

\begin{theorem}
Let A and B be normal operators with spectral expansions
\begin{equation}
A=\int_{\sigma\left(  A\right)  }\alpha P_{A}\left(  \alpha\right)
d\alpha\quad\text{and \quad}B=\int_{\sigma\left(  B\right)  }\beta
P_{B}\left(  \beta\right)  d\beta
\end{equation}

\end{theorem}

$\rfloor$

Let $r=2$ and consider operators $\left\{  \widehat{\Phi_{1}^{\dagger}\Phi
_{1}},\widehat{\Phi_{2}^{\dagger}\Phi_{2}}\right\}  $ that commute i.e.
\begin{equation}
\left[  \widehat{\Phi_{1}^{\dagger}\Phi_{1}},\widehat{\Phi_{2}^{\dagger}%
\Phi_{2}}\right]  =0.
\end{equation}
Eq. $\left(  \ref{EnStation}\right)  $ becomes
\begin{equation}
\left\{  \left(  \widehat{K+V_{2}}\right)  -\lambda_{1}\widehat{\Phi
_{1}^{\dagger}\Phi_{1}}-\lambda_{2}\widehat{\Phi_{2}^{\dagger}\Phi_{2}%
}\right\}  \left\vert \overset{}{a}\right)  =0
\end{equation}
and again one must find solutions $\left(  \lambda_{11},\lambda_{21}\right)  $
such that
\begin{equation}
\det\left\{  \left(  \widehat{K+V_{2}}\right)  -\lambda_{p1}\widehat{\Phi
_{1}^{\dagger}\Phi_{1}}-\lambda_{p2}\widehat{\Phi_{2}^{\dagger}\Phi_{2}%
}\right\}  =0
\end{equation}
and vectors $\left\{  \left\vert \overset{}{a_{p1}}\right)  ,\ldots,\left\vert
\overset{}{a_{pm_{p}}}\right)  \right\}  $that span the $\ker\left\{  \left(
\widehat{K+V_{2}}\right)  -\lambda_{p1}\widehat{\Phi_{1}^{\dagger}\Phi_{1}%
}-\lambda_{p2}\widehat{\Phi_{2}^{\dagger}\Phi_{2}}\right\}  $ and satisfy
\begin{equation}
\left(  a_{pq}\left\vert \left(
\begin{array}
[c]{c}%
\widehat{\Phi_{1}^{\dagger}\Phi_{1}}\\
\widehat{\Phi_{2}^{\dagger}\Phi_{2}}%
\end{array}
\right)  \right.  a_{pq}\right)  =\left(
\begin{array}
[c]{c}%
\gamma_{1}\\
\gamma_{2}%
\end{array}
\right)  .
\end{equation}
This constraint implies
\begin{equation}
\left(  a_{pq}\left\vert \widehat{\Phi_{1}^{\dagger}\Phi_{1}}\right.
a_{pq}\right)  +\left(  a_{pq}\left\vert \widehat{\Phi_{2}^{\dagger}\Phi_{2}%
}\right.  a_{pq}\right)  =\gamma_{1}+\gamma_{2}.
\end{equation}
The operator
\begin{equation}
\widehat{\Omega\left(  \mathbf{\lambda}\right)  }=\lambda_{1}\widehat{\Phi
_{1}^{\dagger}\Phi_{1}}+\lambda_{2}\widehat{\Phi_{2}^{\dagger}\Phi_{2}}%
\end{equation}
is self adjoint and we can find generalized eigen solutions \emph{for any}
$\mathbf{\lambda}$ that produce a non singular operator such that
\begin{equation}
\left(  \widehat{K+V_{2}}\right)  \left\vert \overset{}{a_{\eta}}\left(
\mathbf{\lambda}\right)  \right)  =\eta\left(  \mathbf{\lambda}\right)
\left\{  \lambda_{1}\widehat{\Phi_{1}^{\dagger}\Phi_{1}}+\lambda_{2}%
\widehat{\Phi_{2}^{\dagger}\Phi_{2}}\right\}  \left\vert \overset{}{a_{\eta}%
}\left(  \mathbf{\lambda}\right)  \right)  =\eta\left(  \mathbf{\lambda
}\right)  \widehat{\Omega\left(  \mathbf{\lambda}\right)  }\left\vert
\overset{}{a_{\eta}}\left(  \mathbf{\lambda}\right)  \right)  .
\end{equation}
The solutions are
\begin{equation}
\left\vert \overset{}{\tilde{a}_{_{pq}}}\left(  \mathbf{\lambda}\right)
\right)  =\widehat{\Omega\left(  \mathbf{\lambda}\right)  }^{-\frac{1}{2}%
}\left\vert \overset{}{b_{pq}}\left(  \mathbf{\lambda}\right)  \right)  ,
\end{equation}
where $\left\{  \left\vert \overset{}{b_{pq}}\right)  \right\}  $ are
orthonormal,
\begin{align}
\left(  \tilde{a}_{pq}\left(  \mathbf{\lambda}\right)  \left\vert \left(
\widehat{K+V_{2}}\right)  \tilde{a}_{pq}\left(  \mathbf{\lambda}\right)
\right.  \right)   &  =\eta_{p}\left(  \mathbf{\lambda}\right)  \left(
\tilde{a}_{pq}\left(  \mathbf{\lambda}\right)  \left\vert \widehat
{\Omega\left(  \mathbf{\lambda}\right)  }\tilde{a}_{p^{\prime}q^{\prime}%
}\left(  \mathbf{\lambda}\right)  \right.  \right)  =\eta_{p}\left(
\mathbf{\lambda}\right)  \delta_{pp^{\prime}}\delta_{qq^{\prime}},\\
\left(  \tilde{a}_{pq}\left(  \mathbf{\lambda}\right)  \left\vert
\widehat{\Omega\left(  \mathbf{\lambda}\right)  }\tilde{a}_{p^{\prime
}q^{\prime}}\left(  \mathbf{\lambda}\right)  \right.  \right)   &
=\delta_{pp^{\prime}}\delta_{qq^{\prime}}%
\end{align}
and
\begin{equation}
\left(  \widehat{K+V_{2}}\right)  \left\vert \overset{}{\tilde{a}_{pq}}\left(
\mathbf{\lambda}\right)  \right)  =\eta_{p}\left(  \mathbf{\lambda}\right)
\widehat{\Omega\left(  \mathbf{\lambda}\right)  }\left\vert \overset{}%
{\tilde{a}_{pq}}\left(  \mathbf{\lambda}\right)  \right)  .
\end{equation}
Thus by construction
\begin{align}
\left(  \tilde{a}_{pq}\left(  \mathbf{\lambda}\right)  \left\vert \left\{
\lambda_{1}\widehat{\Phi_{1}^{\dagger}\Phi_{1}}+\lambda_{2}\widehat{\Phi
_{2}^{\dagger}\Phi_{2}}\right\}  \tilde{a}_{p^{\prime}q^{\prime}}\left(
\mathbf{\lambda}\right)  \right.  \right)   &  =\\
\lambda_{1}\left(  \tilde{a}_{pq}\left(  \mathbf{\lambda}\right)  \left\vert
\widehat{\Phi_{1}^{\dagger}\Phi_{1}}\tilde{a}_{p^{\prime}q^{\prime}}\left(
\mathbf{\lambda}\right)  \right.  \right)  +\lambda_{2}\left(  \tilde{a}%
_{pq}\left(  \mathbf{\lambda}\right)  \left\vert \widehat{\Phi_{2}^{\dagger
}\Phi_{2}}\tilde{a}_{p^{\prime}q^{\prime}}\left(  \mathbf{\lambda}\right)
\right.  \right)   &  =\delta_{pqp^{\prime}q^{\prime}}%
\end{align}
i.e.
\begin{equation}
\left(
\begin{array}
[c]{cc}%
\lambda_{1} & \lambda_{2}%
\end{array}
\right)  \left(
\begin{array}
[c]{c}%
\left(  \tilde{a}_{pq}\left(  \mathbf{\lambda}\right)  \left\vert
\widehat{\Phi_{1}^{\dagger}\Phi_{1}}\tilde{a}_{p^{\prime}q^{\prime}}\left(
\mathbf{\lambda}\right)  \right.  \right) \\
\left(  \tilde{a}_{pq}\left(  \mathbf{\lambda}\right)  \left\vert
\widehat{\Phi_{2}^{\dagger}\Phi_{2}}\tilde{a}_{p^{\prime}q^{\prime}}\left(
\mathbf{\lambda}\right)  \right.  \right)
\end{array}
\right)  =\delta_{pqp^{\prime}q^{\prime}}\quad\forall\mathbf{\lambda}%
\end{equation}
Say $m_{1}=1$ then
\begin{equation}
\left(
\begin{array}
[c]{cc}%
\lambda_{1} & \lambda_{2}%
\end{array}
\right)  \left(
\begin{array}
[c]{c}%
\left(  \tilde{a}_{1}\left(  \mathbf{\lambda}\right)  \left\vert \widehat
{\Phi_{1}^{\dagger}\Phi_{1}}\tilde{a}_{1}\left(  \mathbf{\lambda}\right)
\right.  \right) \\
\left(  \tilde{a}_{1}\left(  \mathbf{\lambda}\right)  \left\vert \widehat
{\Phi_{2}^{\dagger}\Phi_{2}}\tilde{a}_{1}\left(  \mathbf{\lambda}\right)
\right.  \right)
\end{array}
\right)  =1\quad\forall\mathbf{\lambda}%
\end{equation}
and
\begin{equation}
\left(  \tilde{a}_{1}\left(  \mathbf{\lambda}\right)  \left\vert \left(
\begin{array}
[c]{c}%
\widehat{\Phi_{1}^{\dagger}\Phi_{1}}\\
\widehat{\Phi_{2}^{\dagger}\Phi_{2}}%
\end{array}
\right)  \right.  \tilde{a}_{1}\left(  \mathbf{\lambda}\right)  \right)
=\left(
\begin{array}
[c]{c}%
\gamma_{1}\\
\gamma_{2}%
\end{array}
\right)  .
\end{equation}
This constraint implies
\begin{equation}
\left(  a_{pq}\left\vert \widehat{\Phi_{1}^{\dagger}\Phi_{1}}\right.
a_{pq}\right)  +\left(  a_{pq}\left\vert \widehat{\Phi_{2}^{\dagger}\Phi_{2}%
}\right.  a_{pq}\right)  =\gamma_{1}+\gamma_{2}.
\end{equation}
such that
\begin{align}
\left(  a_{\eta}\left\vert \left\{  \lambda_{1}\widehat{\Phi_{1}^{\dagger}%
\Phi_{1}}+\lambda_{2}\widehat{\Phi_{2}^{\dagger}\Phi_{2}}\right\}
a_{\eta^{\prime}}\right.  \right)   &  =\left(  \lambda_{1}\gamma_{1}%
+\lambda_{2}\gamma_{2}\right)  \delta_{\eta\eta^{\prime}}\\
\left(  a_{\eta}\left\vert \left\{  \lambda_{1}\widehat{\Phi_{1}^{\dagger}%
\Phi_{1}}\right\}  a_{\eta}\right.  \right)  +\left(  a_{\eta}\left\vert
\left\{  \lambda_{2}\widehat{\Phi_{2}^{\dagger}\Phi_{2}}\right\}  a_{\eta
}\right.  \right)   &  =1.
\end{align}
This gives
\begin{equation}
\lambda_{1}\gamma_{1}+\lambda_{2}\gamma_{2}=1.
\end{equation}
As the family $\left\{  \widehat{\Phi_{j}^{\dagger}\Phi_{j}}\right\}  $
commute one can find a basis that simultaneously diagonalizes these operators.

For solutions corresponding to a fixed $\lambda\in\mathbb{R}^{2}$ the
constraints are
\begin{equation}
\sum_{1\leq p,q\leq m}\alpha_{p}\alpha_{q}\left(  a_{1p}\left\vert \left(
\begin{array}
[c]{c}%
\widehat{\Phi_{1}^{\dagger}\Phi_{1}}\\
\widehat{\Phi_{2}^{\dagger}\Phi_{2}}%
\end{array}
\right)  \right.  a_{1q}\right)  =\left(
\begin{array}
[c]{c}%
\gamma_{1}\\
\gamma_{2}%
\end{array}
\right)  .
\end{equation}


Eq. $\left(  \ref{surface}\right)  $ is a level surface of the function
\begin{equation}
\mathcal{\Xi}\left(  a\right)  \in\mathbb{R}^{r}%
\end{equation}
where
\begin{equation}
\mathcal{\Xi}\left(  a\right)  =\left(  \left(  a\left\vert \widehat
{\Phi^{\dagger}\Phi}_{1}\right.  a\right)  ,\cdots,\left(  a\left\vert
\widehat{\Phi^{\dagger}\Phi}_{r}\right.  a\right)  \right)  .
\end{equation}
The surface $\mathcal{S}\left(  \gamma\right)  $ $\equiv$ $\left\{  \left.
a\right\vert \mathcal{\Xi}\left(  a\right)  =\gamma\right\}  $ can be locally
parameterized by $r$ parameters. $\lceil$%
\begin{equation}
\gamma_{j}=\left(  a\left\vert \widehat{\Phi_{j}^{\dagger}\Phi_{j}}\right.
a\right)  =\sum a_{k}\left(  \widehat{\Phi_{j}^{\dagger}\Phi_{j}}\right)
_{km}a_{m}=\sum_{\substack{1\leq j\leq r\\1\leq k,m\leq r_{N}^{2}}}\left(
\widehat{\Phi_{j}^{\dagger}\Phi_{j}}\right)  _{km}a_{k}a_{m}%
\end{equation}
i.e.
\begin{align}
\gamma &  =\left[  \widehat{\Phi^{\dagger}\Phi}\right]  \left(  a\otimes
a\right) \\
\left[  \widehat{\Phi^{\dagger}\Phi}\right]   &  :\mathbb{R}^{r_{N}^{4}%
}\rightarrow\mathbb{R}^{r}%
\end{align}
$\left[  \widehat{\Phi^{\dagger}\Phi}\right]  $ is a positive linear map
\begin{equation}
\left[  \widehat{\Phi^{\dagger}\Phi}\right]  :\mathcal{B}_{2sa}\left(
\mathcal{H}^{N}\right)  \otimes\mathcal{B}_{2sa}\left(  \mathcal{H}%
^{N}\right)  \rightarrow L_{1}\left(  \mathbb{X}\right)
\end{equation}
If $\left[  \widehat{\Phi^{\dagger}\Phi}\right]  $ is surjective
\begin{equation}
\left[  \widehat{\Phi^{\dagger}\Phi}\right]  =\sum_{1\leq j\leq r}\alpha
_{j}\left\vert \mathbf{f}_{j}\right)  \left(  \mathbf{e}_{j}\right\vert
\end{equation}
and
\begin{equation}
\left[  \widehat{\Phi^{\dagger}\Phi}\right]  ^{G}=\sum_{1\leq j\leq r}%
\alpha_{j}^{-1}\left\vert \mathbf{e}_{j}\right)  \left(  \mathbf{f}%
_{j}\right\vert ,
\end{equation}
where $\left\{  \left\vert \mathbf{f}_{j}\right)  ;1\leq j\leq r\right\}  $ is
a basis of $L_{1}\left(  \mathbb{X}\right)  $ and $\left\{  \left\vert
\mathbf{e}_{j}\right)  ;1\leq j\leq r\right\}  $ is a basis of $\ker\left[
\widehat{\Phi^{\dagger}\Phi}\right]  ^{\bot},$ which gives that
\begin{equation}
\left[  \widehat{\Phi^{\dagger}\Phi}\right]  ^{-1}\left(  \gamma\right)
=\left[  \widehat{\Phi^{\dagger}\Phi}\right]  ^{G}\left(  \gamma\right)
\oplus\ker\left[  \widehat{\Phi^{\dagger}\Phi}\right]  .
\end{equation}
Thus we can consider
\begin{equation}
\left[  a\otimes a\right]  \left(  \kappa,\gamma\right)  =\left[
\widehat{\Phi^{\dagger}\Phi}\right]  ^{G}\left(  \gamma\right)  +\kappa_{\mu
};\quad\kappa_{\mu}\in\ker\left[  \widehat{\Phi^{\dagger}\Phi}\right]
\end{equation}
for fixed $\kappa_{\mu}$, where $\left\{  \kappa_{\mu}\right\}  $ is a basis
for $\ker\left[  \widehat{\Phi^{\dagger}\Phi}\right]  .$%
\begin{align}
\left[  \widehat{\Phi^{\dagger}\Phi}\right]  \left[  \widehat{\Phi^{\dagger
}\Phi}\right]  ^{G}  &  =\sum_{1\leq j\leq r}\left\vert \mathbf{f}_{j}\right)
\left(  \mathbf{f}_{j}\right\vert =P_{\operatorname{Ran}\left[  \widehat
{\Phi^{\dagger}\Phi}\right]  }\\
\left[  \widehat{\Phi^{\dagger}\Phi}\right]  ^{G}\left[  \widehat
{\Phi^{\dagger}\Phi}\right]   &  =\sum_{1\leq j\leq r}\left\vert
\mathbf{e}_{j}\right)  \left(  \mathbf{e}_{j}\right\vert =P_{\ker\left[
\widehat{\Phi^{\dagger}\Phi}\right]  ^{\bot}}%
\end{align}
The generalized inverse gives
\begin{equation}
\sum_{1\leq j\leq r}\left[  \widehat{\Phi^{\dagger}\Phi}\right]  _{k_{1}%
k_{2}m_{1}m_{2}j}^{G}\gamma_{j}=b_{k_{1}k_{2}m_{1}m_{2}}%
\end{equation}
and
\begin{equation}
\left[  \widehat{\Phi^{\dagger}\Phi}\right]  ^{G}\left(  \gamma\right)
=\sum_{1\leq k_{1},k_{2},m_{1},m_{2}\leq r_{N}^{2}}b_{k_{1}k_{2}m_{1}m_{2}%
}\Lambda_{k_{1}k_{2}}\otimes\Lambda_{m_{1}m_{2}}.
\end{equation}
The question arises whether $\left[  \widehat{\Phi^{\dagger}\Phi}\right]
^{G}$ is a positive map, is $\left[  \widehat{\Phi^{\dagger}\Phi}\right]
^{G}\left(  \gamma\right)  $ even positive?. Let $a_{0}\left(  \gamma\right)
$ be a reference one dimensional projector for $\gamma$ then
\begin{equation}
a_{0}\left(  \gamma\right)  \otimes a_{0}\left(  \gamma\right)  =\left[
\widehat{\Phi^{\dagger}\Phi}\right]  ^{G}\left(  \gamma\right)  +\kappa\left(
a_{0}\left(  \gamma\right)  \right)
\end{equation}
and
\begin{equation}
\left[  \gamma\right]  _{N}=\left\{  a_{0}\left(  \gamma\right)  \otimes
a_{0}\left(  \gamma\right)  \oplus\ker\left[  \widehat{\Phi^{\dagger}\Phi
}\right]  \right\}  .
\end{equation}


The operator $a+\delta a$ belongs to the level set if
\begin{align}
&  \left(  a+\delta a\left\vert \widehat{\Phi^{\dagger}\Phi}\right.  \left(
a+\delta a\right)  \right)  -\gamma=0\\
&  \Leftrightarrow2\left(  a\left\vert \widehat{\Phi^{\dagger}\Phi}\right.
\delta a\right)  +\left(  \delta a\left\vert \widehat{\Phi^{\dagger}\Phi
}\right.  \delta a\right)  =0,
\end{align}
which is satisfied if $\delta a\in\ker\left\{  \widehat{\Phi^{\dagger}%
\Phi\left(  \mathbf{x}\right)  }\right\}  \quad\forall\mathbf{x}$ or $\delta
a\notin\ker\left\{  \widehat{\Phi^{\dagger}\Phi\left(  \mathbf{x}\right)
}\right\}  \quad\forall\mathbf{x}$ and
\begin{equation}
2\left(  a\left\vert \widehat{\Phi^{\dagger}\Phi}\right.  \delta a\right)
=-\left(  \delta a\left\vert \widehat{\Phi^{\dagger}\Phi}\right.  \delta
a\right)  .
\end{equation}
One should note that $a+\delta a\in\mathcal{B}_{2sa}\left(  \mathcal{H}%
^{N}\right)  $ and $\operatorname{Tr}\left\{  \left(  a+\delta a\right)
^{2}N\right\}  =N$ implies $\operatorname{Tr}\left\{  \left(  a+\delta
a\right)  ^{2}\right\}  =1$ so that
\begin{align}
2\operatorname{Tr}\left\{  a\delta a\right\}  +\operatorname{Tr}\left\{
\delta a^{2}\right\}   &  =0\\
\left(  a\left\vert \delta a\right.  \right)   &  =-\frac{1}{2}\left(  \delta
a\left\vert \delta a\right.  \right)  \neq0.
\end{align}
Which is equivalent to
\begin{equation}
\left(  2a+\delta a\left\vert \delta a\right.  \right)  =0.
\end{equation}
Let
\begin{equation}
\delta a=\alpha a+b;\quad b\bot a
\end{equation}
then
\begin{equation}
\left(  \left(  2+\alpha\right)  a+b\left\vert \alpha a+b\right.  \right)
=0=\left(  2+\alpha\right)  \alpha+\left(  b\left\vert b\right.  \right)
=\alpha^{2}+2\alpha+\left(  b\left\vert b\right.  \right)
\end{equation}
as $\left(  a\left\vert a\right.  \right)  =1.$ $\lceil$We can always choose
the reference $a\left(  \gamma\right)  $ to be a one dimensional projector so
that
\begin{equation}
a\left(  \gamma\right)  ^{2}=a\left(  \gamma\right)  .\rfloor
\end{equation}
Thus%
\begin{equation}
\alpha=-1\pm\sqrt{1-\left(  b\left\vert b\right.  \right)  }.
\end{equation}


\subsection{Example}

\subsection{Implicit Function Theorem}

We first consider the discrete constraint case with Lagrangian
\begin{equation}
\mathcal{L}\left(  Q,Q^{\dagger},\lambda,\gamma\right)  =\left\{  \sum_{1\leq
k,l,m\leq r_{N}}\bar{Q}_{kl}H_{km}Q_{ml}-\sum_{1\leq j\leq r}\lambda
_{j}\left\{  \bar{Q}_{kl}T_{kmj}^{1}Q_{ml}-\gamma_{j}\right\}  \right\}
\end{equation}
and collect the variables of the Lagrangian into the form
\begin{equation}
\mathcal{L}\left(  \varsigma,\gamma\right)  \equiv\mathcal{L}\left(  Q,\bar
{Q},\lambda,\gamma\right)  .
\end{equation}
The function that defines the implicit relationship is
\begin{align}
F  &  :\mathbb{C}^{2r_{^{N}}^{2}}\times\mathbb{R}^{2r}\rightarrow
\mathbb{C}^{2r_{^{N}}^{2}}\times\mathbb{R}^{r}\nonumber\\
F\left(  \varsigma,\gamma\right)   &  \equiv D_{\varsigma}\mathcal{L}\left(
\varsigma,\gamma\right)  =\left(  F_{1}\left(  \varsigma,\gamma\right)
,\cdots,F_{2r_{^{N}}^{2}+r}\left(  \varsigma,\gamma\right)  \right)
=\mathbf{0\in}\mathbb{C}^{2r_{^{N}}^{2}}\times\mathbb{R}^{r},
\end{align}
where
\begin{equation}
F_{\mu}\left(  \varsigma,\gamma\right)  =\left\{
\begin{array}
[c]{c}%
\sum\limits_{1\leq k\leq r_{N}}\bar{Q}_{kl}\left\{  \bar{H}_{mk}%
-\sum\limits_{1\leq j\leq r}\lambda_{j}\bar{T}_{mkj}^{1}\right\}  ;\quad
1\leq\mu\leq r_{^{N}}^{2},\quad\mu=\left(  ml\right)  \qquad\qquad\quad\\
\sum\limits_{1\leq k\leq r_{N}}Q_{kl}\left\{  H_{mk}-\sum\limits_{1\leq j\leq
r}\lambda_{j}T_{mkj}^{1}\right\}  ;\quad r_{^{N}}^{2}+1\leq\mu\leq2r_{^{N}%
}^{2},\quad\mu=\left(  ml\right)  +r_{^{N}}^{2}\\
\sum\limits_{1\leq k,l,m\leq r_{N}}\bar{Q}_{kl}T_{mkj}^{1}Q_{ml}-\gamma
_{j};\quad2r_{^{N}}^{2}+1\leq\mu\leq2r_{^{N}}^{2}+r,\quad\mu=j+2r_{^{N}}%
^{2}\quad\quad
\end{array}
\right.
\end{equation}
and
\begin{equation}
F_{\mu}\left(  \varsigma,\gamma\right)  =\left\{
\begin{array}
[c]{c}%
\partial_{Q_{ml}}\mathcal{L}\left(  Q,Q^{\dagger},\lambda,\gamma\right)
;\quad1\leq\mu\leq r_{^{N}}^{2},\quad\mu=\left(  ml\right)  \qquad\qquad
\qquad\\
\partial_{\bar{Q}_{ml}}\mathcal{L}\left(  Q,Q^{\dagger},\lambda,\gamma\right)
;\quad r_{^{N}}^{2}+1\leq\mu\leq2r_{^{N}}^{2},\quad\mu=\left(  ml\right)
+r_{^{N}}^{2}\quad\\
\partial_{\lambda_{j}}\mathcal{L}\left(  Q,Q^{\dagger},\lambda,\gamma\right)
;\quad2r_{^{N}}^{2}+1\leq\mu\leq2r_{^{N}}^{2}+r,\quad\mu=j+2r_{^{N}}^{2}%
\end{array}
\right.  .
\end{equation}
The derivatives of $F$ are
\begin{align}
\partial_{Q_{pq}}F_{\mu}\left(  \varsigma,\gamma\right)   &  =\left\{
\begin{array}
[c]{c}%
0;\quad1\leq\mu\leq r_{^{N}}^{2},\quad\mu=\left(  ml\right)  \qquad
\qquad\qquad\qquad\qquad\qquad\qquad\qquad\\
\left\{  H_{mp}-\sum\limits_{1\leq j\leq r}\lambda_{j}T_{mpj}^{1}\right\}
\delta_{lq};\quad r_{^{N}}^{2}+1\leq\mu\leq2r_{^{N}}^{2},\quad\mu=\left(
ml\right)  +r_{^{N}}^{2}\quad\\
\sum\limits_{1\leq m\leq r_{N}}\bar{Q}_{mq}T_{mp\mu}^{1};\quad2r_{^{N}}%
^{2}+1\leq\mu\leq2r_{^{N}}^{2}+r,\quad\mu=j+2r_{^{N}}^{2}\quad\quad\quad
\quad\quad
\end{array}
\right. \nonumber\\
\partial_{\bar{Q}_{pq}}F_{\mu}\left(  \varsigma,\gamma\right)   &  =\left\{
\begin{array}
[c]{c}%
\left\{  \bar{H}_{mp}-\sum\limits_{1\leq j\leq r}\lambda_{j}\bar{T}_{mpj}%
^{1}\right\}  \delta_{lq};\quad1\leq\mu\leq r_{^{N}}^{2},\quad\mu=\left(
ml\right)  \qquad\qquad\qquad\quad\quad\\
0;\quad r_{^{N}}^{2}+1\leq\mu\leq2r_{^{N}}^{2},\quad\mu=\left(  ml\right)
+r_{^{N}}^{2}\qquad\qquad\qquad\qquad\qquad\qquad\qquad\\
\sum\limits_{1\leq m\leq r_{N}}T_{pm\mu}^{1}Q_{mq};\quad2r_{^{N}}^{2}+1\leq
\mu\leq2r_{^{N}}^{2}+r,\quad\mu=j+2r_{^{N}}^{2}\qquad\qquad\qquad
\end{array}
\right. \nonumber\\
\partial_{\lambda_{p}}F_{\mu}\left(  \varsigma,\gamma\right)   &  =\left\{
\begin{array}
[c]{c}%
\sum\limits_{1\leq k\leq r_{N}}\bar{Q}_{kl}\left\{  \bar{H}_{mk}-\bar{T}%
_{mkp}^{1}\right\}  ;\quad1\leq\mu\leq r_{^{N}}^{2},\quad\mu=\left(
ml\right)  \quad\quad\quad\quad\quad\\
\sum\limits_{1\leq k\leq r_{N}}Q_{kl}\left\{  H_{mk}-T_{mkp}^{1}\right\}
;\quad r_{^{N}}^{2}+1\leq\mu\leq2r_{^{N}}^{2},\quad\mu=\left(  ml\right)
+r_{^{N}}^{2}\\
0;\quad2r_{^{N}}^{2}+1\leq\mu\leq2r_{^{N}}^{2}+r,\quad\mu=j+2r_{^{N}}^{2}%
\quad\quad\quad\quad\quad\quad\quad\quad\quad
\end{array}
\right. \nonumber\\
\partial_{\gamma_{p}}F_{\mu}\left(  \varsigma,\gamma\right)   &  =\left\{
\begin{array}
[c]{c}%
0;\quad1\leq\mu\leq r_{^{N}}^{2},\quad\mu=\left(  ml\right)  \quad\quad
\quad\quad\quad\\
0;\quad r_{^{N}}^{2}+1\leq\mu\leq2r_{^{N}}^{2},\quad\mu=\left(  ml\right)
+r_{^{N}}^{2}\\
\delta_{pj};\quad1\leq p,j\leq r;\quad\mu=j+2r_{^{N}}^{2}\quad\quad\quad
\end{array}
\right.
\end{align}
and
\begin{align}
\partial_{Q_{pq}}F_{\mu}\left(  \varsigma,\gamma\right)   &  =\left\{
\begin{array}
[c]{c}%
\partial_{Q_{pq}Q_{ml}}^{2}\mathcal{L}\left(  Q,Q^{\dagger},\lambda
,\gamma\right)  ;\quad1\leq\mu\leq r_{^{N}}^{2},\quad\mu=\left(  ml\right)
\qquad\qquad\quad\\
\partial_{Q_{pq}\bar{Q}_{ml}}^{2}\mathcal{L}\left(  Q,Q^{\dagger}%
,\lambda,\gamma\right)  ;\quad r_{^{N}}^{2}+1\leq\mu\leq2r_{^{N}}^{2},\quad
\mu=\left(  ml\right)  +r_{^{N}}^{2}\\
\partial_{Q_{pq}\lambda_{j}}^{2}\mathcal{L}\left(  Q,Q^{\dagger}%
,\lambda,\gamma\right)  ;\quad2r_{^{N}}^{2}+1\leq\mu\leq2r_{^{N}}^{2}%
+r,\quad\mu=j+2r_{^{N}}^{2}%
\end{array}
\right. \nonumber\\
\partial_{\bar{Q}_{pq}}F_{\mu}\left(  \varsigma,\gamma\right)   &  =\left\{
\begin{array}
[c]{c}%
\partial_{\bar{Q}_{pq}Q_{ml}}^{2}\mathcal{L}\left(  Q,Q^{\dagger}%
,\lambda,\gamma\right)  ;\quad1\leq\mu\leq r_{^{N}}^{2},\quad\mu=\left(
ml\right)  \qquad\qquad\quad\\
\partial_{\bar{Q}_{pq}\bar{Q}_{ml}}^{2}\mathcal{L}\left(  Q,Q^{\dagger
},\lambda,\gamma\right)  ;\quad r_{^{N}}^{2}+1\leq\mu\leq2r_{^{N}}^{2}%
,\quad\mu=\left(  ml\right)  +r_{^{N}}^{2}\\
\partial_{\bar{Q}_{pq}\lambda_{j}}^{2}\mathcal{L}\left(  Q,Q^{\dagger}%
,\lambda,\gamma\right)  ;\quad2r_{^{N}}^{2}+1\leq\mu\leq2r_{^{N}}^{2}%
+r,\quad\mu=j+2r_{^{N}}^{2}%
\end{array}
\right. \nonumber\\
\partial_{\lambda_{p}}F_{\mu}\left(  \varsigma,\gamma\right)   &  =\left\{
\begin{array}
[c]{c}%
\partial_{\lambda_{p}Q_{ml}}^{2}\mathcal{L}\left(  Q,Q^{\dagger}%
,\lambda,\gamma\right)  ;\quad1\leq\mu\leq r_{^{N}}^{2},\quad\mu=\left(
ml\right)  \qquad\qquad\quad\\
\partial_{\lambda_{p}\bar{Q}_{ml}}^{2}\mathcal{L}\left(  Q,Q^{\dagger}%
,\lambda,\gamma\right)  ;\quad r_{^{N}}^{2}+1\leq\mu\leq2r_{^{N}}^{2},\quad
\mu=\left(  ml\right)  +r_{^{N}}^{2}\\
\partial_{\lambda_{p}\lambda_{j}}^{2}\mathcal{L}\left(  Q,Q^{\dagger}%
,\lambda,\gamma\right)  ;\quad2r_{^{N}}^{2}+1\leq\mu\leq2r_{^{N}}^{2}%
+r,\quad\mu=j+2r_{^{N}}^{2}%
\end{array}
\right. \nonumber\\
\partial_{\gamma_{p}}F_{\mu}\left(  \varsigma,\gamma\right)   &  =\left\{
\begin{array}
[c]{c}%
\partial_{\gamma_{p}Q_{ml}}^{2}\mathcal{L}\left(  Q,Q^{\dagger},\lambda
,\gamma\right)  ;\quad1\leq\mu\leq r_{^{N}}^{2},\quad\mu=\left(  ml\right)
\qquad\qquad\quad\\
\partial_{\gamma_{p}\bar{Q}_{ml}}^{2}\mathcal{L}\left(  Q,Q^{\dagger}%
,\lambda,\gamma\right)  ;\quad r_{^{N}}^{2}+1\leq\mu\leq2r_{^{N}}^{2},\quad
\mu=\left(  ml\right)  +r_{^{N}}^{2}\\
\partial_{\gamma_{p}\lambda_{j}}^{2}\mathcal{L}\left(  Q,Q^{\dagger}%
,\lambda,\gamma\right)  ;\quad2r_{^{N}}^{2}+1\leq\mu\leq2r_{^{N}}^{2}%
+r,\quad\mu=j+2r_{^{N}}^{2}%
\end{array}
\right.
\end{align}
We want to find solutions $\left(  \varsigma,\gamma\right)  $ in a
neighborhood of a point $\left(  \varsigma_{0},\gamma_{0}\right)  ,$ where
$D_{\varsigma}\mathcal{L}\left(  \varsigma_{0},\gamma_{0}\right)
=\mathbf{0}.$ Expanding $F$ up to first order about the point $\left(
\varsigma_{0},\gamma_{0}\right)  $ gives
\begin{equation}
F\left(  \varsigma_{0}+\delta\varsigma,\gamma_{0}+\delta\gamma\right)
=\left(  D_{\varsigma}F\left(  \varsigma_{0},\gamma_{0}\right)  ,D_{\gamma
}F\left(  \varsigma_{0},\gamma_{0}\right)  \right)  \left(
\begin{array}
[c]{c}%
\delta\varsigma\\
\delta\gamma
\end{array}
\right)  ,
\end{equation}
where
\begin{align}
D_{\varsigma\gamma}F\left(  \varsigma_{0},\gamma_{0}\right)   &  =\left(
\begin{array}
[c]{cccccc}%
\partial_{\varsigma_{1}}F_{1} & \cdots & \partial_{\varsigma_{M}}F_{1} &
\partial_{\gamma_{1}}F_{1} & \cdots & \partial_{\gamma_{r}}F_{1}\\
\vdots & \vdots & \vdots & \vdots & \vdots & \vdots\\
\partial_{\varsigma_{1}}F_{M} & \cdots & \partial_{\varsigma_{M}}F_{M} &
\partial_{\gamma_{1}}F_{M} & \cdots & \partial_{\gamma_{r}}F_{M}%
\end{array}
\right) \nonumber\\
&  =\left(
\begin{array}
[c]{cccc}%
D_{QQ}^{2}\mathcal{L} & D_{\bar{Q}Q}^{2}\mathcal{L} & D_{\lambda Q}%
^{2}\mathcal{L} & D_{\gamma Q}^{2}\mathcal{L}\\
D_{Q\bar{Q}}^{2}\mathcal{L} & D_{\bar{Q}\bar{Q}}^{2}\mathcal{L} &
D_{\lambda\bar{Q}}^{2}\mathcal{L} & D_{\gamma\bar{Q}}^{2}\mathcal{L}\\
D_{Q\lambda}^{2}\mathcal{L} & D_{\bar{Q}\lambda}^{2}\mathcal{L} &
D_{\lambda\lambda}^{2}\mathcal{L} & D_{\gamma\lambda}^{2}\mathcal{L}%
\end{array}
\right) \nonumber\\
&  =\left(
\begin{array}
[c]{cccc}%
\mathbb{O}_{r_{N}^{2}\times r_{N}^{2}} & D_{\bar{Q}Q}^{2}\mathcal{L} &
D_{\lambda Q}^{2}\mathcal{L} & \mathbb{O}_{r\times r_{N}^{2}}\\
D_{Q\bar{Q}}^{2}\mathcal{L} & \mathbb{O}_{r_{N}^{2}\times r_{N}^{2}} &
D_{\lambda\bar{Q}}^{2}\mathcal{L} & \mathbb{O}_{r\times r_{N}^{2}}\\
D_{Q\lambda}^{2}\mathcal{L} & D_{\bar{Q}\lambda}^{2}\mathcal{L} &
\mathbb{O}_{r\times r} & \mathbb{I}_{r\times r}%
\end{array}
\right)
\end{align}
and $M=2r_{^{N}}^{2}+r$ so $D_{\varsigma\gamma}F\left(  \varsigma_{0}%
,\gamma_{0}\right)  $ is $M\times\left(  M+r\right)  $ matrix.

Eq. $\left(  \ref{nbhd}\right)  $ is satisfied in small enough nbhd
$\mathcal{N}_{\varsigma_{0}\gamma_{0}}$ of $\left(  \varsigma_{0},\gamma
_{0}\right)  $ if
\begin{equation}
\left(  D_{\varsigma}F\left(  \varsigma_{0},\gamma_{0}\right)  ,D_{\gamma
}F\left(  \varsigma_{0},\gamma_{0}\right)  \right)  \left(
\begin{array}
[c]{c}%
\delta\varsigma\\
\delta\gamma
\end{array}
\right)  =0\quad\forall\left(
\begin{array}
[c]{c}%
\varsigma\\
\gamma
\end{array}
\right)  \in\mathcal{N}_{\varsigma_{0}\gamma_{0}}.
\end{equation}
This gives
\begin{equation}
D_{\varsigma}F\left(  \varsigma_{0},\gamma_{0}\right)  \left(  \delta
\varsigma\right)  =-D_{\gamma}F\left(  \varsigma_{0},\gamma_{0}\right)
\left(  \delta\gamma\right)
\end{equation}
and
\begin{equation}
\delta\varsigma=-\left\{  \left[  D_{\varsigma}F\left(  \varsigma_{0}%
,\gamma_{0}\right)  \right]  ^{G}\circ\left[  D_{\gamma}F\left(  \varsigma
_{0},\gamma_{0}\right)  \right]  \right\}  \left(  \delta\gamma\right)
+\sigma,
\end{equation}
where $\left[  D_{\varsigma}F\left(  \varsigma_{0},\gamma_{0}\right)  \right]
^{G}$ is the generalized inverse of $D_{\varsigma}F\left(  \varsigma
_{0},\gamma_{0}\right)  $ and $\sigma\in\ker D_{\varsigma}F\left(
\varsigma_{0},\gamma_{0}\right)  .$ Thus in $\mathcal{N}_{\varsigma_{0}%
\gamma_{0}}$
\begin{align}
\varsigma &  =-\left\{  \left[  D_{\varsigma}F\left(  \varsigma_{0},\gamma
_{0}\right)  \right]  ^{G}\circ\left[  D_{\gamma}F\left(  \varsigma_{0}%
,\gamma_{0}\right)  \right]  \right\}  \left(  \gamma\right) \\
&  -\left\{  \left[  D_{\varsigma}F\left(  \varsigma_{0},\gamma_{0}\right)
\right]  ^{G}\circ\left[  D_{\gamma}F\left(  \varsigma_{0},\gamma_{0}\right)
\right]  \right\}  \left(  \gamma_{0}\right)  +\varsigma_{0}+\sigma
_{\varsigma_{0}\gamma_{0}}\\
&  \equiv f\left(  \gamma,\sigma_{\varsigma_{0}\gamma_{0}}\right)
\end{align}
If As the domains are linear spaces $\delta\varsigma=\varsigma-\varsigma_{0}$
and $\delta\gamma=\gamma-\gamma_{0}$ and $\mathcal{L}\left(  \varsigma
,\gamma\right)  $ is a quadratic function, solutions to $D\mathfrak{F}\left(
\delta\varsigma,\delta\gamma\right)  =0$ are the same as solutions to
$D\mathcal{L}\left(  \varsigma,\gamma\right)  =0$ and we need to solve the
linear equations which

If these minima are unique and then one could define kinetic energy, coulomb
and exchange - correlation functionals by eq. $\left(  \ref{HK}\right)  $ as
$.$
\begin{align}
\mathcal{F}_{K}\left(  \gamma\right)   &  =\operatorname{Tr}\left\{  KD_{\ast
}^{1}\left(  \gamma\right)  \right\}  ,\nonumber\\
\mathcal{F}_{C}\left(  \gamma\right)   &  =\int\frac{\gamma\left(
\mathbf{r,\xi}\right)  \gamma\left(  \mathbf{r}^{\prime}\mathbf{,\xi}\right)
}{\left\Vert \mathbf{r-r}^{\prime}\right\Vert }d\mathbf{r}d\mathbf{r}^{\prime
}d\mathbf{\xi},\nonumber\\
\mathcal{F}_{XC}\left(  \gamma\right)   &  =\operatorname{Tr}\left\{
V_{2}D_{\ast}^{2}\left(  \gamma\right)  \right\}  -\mathcal{F}_{C}\left(
\gamma\right)  ,
\end{align}
producing the decomposition
\begin{equation}
\mathcal{F}_{HK}\left(  \gamma\right)  =\mathcal{F}_{K}\left(  \gamma\right)
+\mathcal{F}_{C}\left(  \gamma\right)  +\mathcal{F}_{XC}\left(  \gamma\right)
.
\end{equation}
The total energy functional of a system that contains on average $N$ electrons
and experiences an environment described by a \emph{local} potential
$\mathrm{v}$ is
\begin{equation}
\mathcal{E}\left(  \gamma\mathbf{,}\mathrm{v}\right)  =\mathcal{F}_{HK}\left(
\gamma\right)  +f_{\gamma}\left(  \mathrm{v}\right)  ,
\end{equation}
where the explicit effect of the environment is described by the functional
\begin{equation}
f_{\gamma}\left(  \mathrm{v}\right)  =\int\mathrm{v}\left(  \mathbf{r,\xi
}\right)  \gamma\left(  \mathbf{r,\xi}\right)  d\mathbf{r}d\mathbf{\xi.}%
\end{equation}
The ground state energy $E_{GS}\left(  \mathrm{v}\right)  $ and space-spin
density $\gamma_{GS}$ $\left(  \mathrm{v}\right)  $ are then determined by
\begin{equation}
E_{GS}\left(  \mathrm{v}\right)  =\mathcal{E}\left(  \gamma_{GS}%
\mathbf{,}\mathrm{v}\right)  =\min_{\mathbf{\gamma\in}\mathcal{D}}%
\mathcal{E}\left(  \gamma\mathbf{,}\mathrm{v}\right)  ,
\end{equation}
where
\begin{equation}
\mathcal{D=}\left\{  \gamma\left\vert \int\gamma\left(  \mathbf{r,\xi}\right)
d\mathbf{r}d\mathbf{\xi}=N,\quad\gamma\left(  \mathbf{r,\xi}\right)
\geq0\right.  \right\}  .
\end{equation}
$\lceil$The domain $\mathcal{D}$ is a compact (closed in $L_{1}$ topology and
bounded) convex set thus a unique minimum exists if ...?$\rfloor$

The density $\gamma_{GS}$ $\left(  \mathrm{v}\right)  $ can be always
expressed in terms of the Natural General Spin Orbitals (NGSO's) $\left\{
\psi_{j}\right\}  $ and occupation numbers $\left\{  N_{j}\right\}  $ of
$D_{\ast}^{1}\left(  \gamma\right)  $ as%

\begin{equation}
\gamma_{GS}\left(  \mathrm{v}\right)  \left(  \mathbf{x}\right)  =D_{\ast}%
^{1}\left(  \gamma_{GS}\left(  \mathrm{v}\right)  \right)  \left(
\mathbf{x,x}\right)  =\sum_{1\leq j\leq r}N_{j}\left\vert \psi_{j}\left(
\mathbf{x}\right)  \right\vert ^{2},
\end{equation}
but such expansions are not unique, in particular Harimanns \cite{Harimann}
theorem ensures that there are always expansions of the form
\begin{equation}
\gamma_{GS}\left(  \mathrm{v}\right)  \left(  \mathbf{x}\right)  =D_{\ast}%
^{1}\left(  \gamma_{GS}\left(  \mathrm{v}\right)  \right)  \left(
\mathbf{x,x}\right)  =\sum_{1\leq j\leq N}\left\vert \varphi_{j}\left(
\mathbf{x}\right)  \right\vert ^{2}%
\end{equation}
in terms of square summable functions $\left\{  \left.  \varphi_{j}\right\vert
1\leq j\leq N\right\}  $.

\subsection{Kohn-Sham DFT}

\bigskip One can introduce a kinetic energy functional
\begin{equation}
\mathcal{F}_{T}\left(  \gamma\right)  =\operatorname{Tr}\left\{  \left(
K+\mathrm{w}\left(  \gamma\right)  \right)  D_{\#}^{1}\right\}  =\min
_{D^{1}\in\left[  \gamma\right]  _{\left\langle N\right\rangle 1}%
}\operatorname{Tr}\left\{  \left(  K+\mathrm{w}\left(  \gamma\right)  \right)
D^{1}\right\}  , \label{KEFunc}%
\end{equation}
where $\mathrm{w}\left(  \gamma\right)  $ is a potential that produces
$\gamma$ as a minimum.

$\lceil$As $K$ is not a bounded operator (thus does not define a continuous
functional of $D^{1}$) the $\min$ might not exist and should be replaced with
$\inf$ even though the domain $\left[  \gamma\right]  _{\left\langle
N\right\rangle 1}$ is a compact convex set. The resulting functional is convex
?
\begin{equation}
\mathcal{F}_{T}\left(  \gamma\right)  =\mathcal{F}_{T}\left(  \alpha_{1}%
\gamma_{1}+\alpha_{2}\gamma_{2}\right)  \leq\alpha_{1}\mathcal{F}_{T}\left(
\gamma_{1}\right)  +\alpha_{2}\mathcal{F}_{T}\left(  \gamma_{2}\right)  ?
\end{equation}
$\rfloor$ which can be expanded as
\begin{equation}
\mathcal{F}_{T}\left(  \gamma\right)  =\sum_{1\leq j\leq r}\left\langle
\varphi_{\#j}\left\vert K\varphi_{\#j}\right.  \right\rangle N_{j},
\end{equation}
where
\begin{align}
D_{\#}^{1}\left(  \mathbf{x,x}^{\prime}\right)   &  =\sum_{1\leq j\leq r}%
N_{j}\varphi_{\#j}\left(  \mathbf{x}\right)  \varphi_{\#j}\left(
\mathbf{x}^{\prime}\right)  ^{\ast}\nonumber\\
\gamma\left(  \mathbf{x}\right)   &  =\sum_{1\leq j\leq r}\left\vert
\varphi_{\#j}\left(  \mathbf{x}\right)  \right\vert ^{2}N_{j}. \label{gammaKS}%
\end{align}
For a given space-spin density $\gamma$ one can define a total energy
functional
\begin{equation}
\mathcal{E}_{NI}\left(  \gamma\mathbf{,}\mathrm{w}\right)  =\mathcal{F}%
_{T}\left(  \gamma\right)  +f_{\gamma}\left(  \mathrm{w}\right)
\end{equation}
that describes groups of on average $N$ noninteracting electrons in an
environment characterized by a local potential $\mathrm{w.}$ One can then
determine a local potential $\mathrm{w}_{\#}$ such that
\begin{equation}
E_{NIGS}\left(  \mathrm{w}_{\#}\right)  =\mathcal{E}_{NIGS}\left(  \gamma
_{GS}\left(  \mathrm{v}\right)  \mathbf{,}\mathrm{w}_{\#}\right)
=\min_{\mathbf{\gamma\in}\mathcal{D}}\mathcal{E}_{NI}\left(  \gamma
\mathbf{,}\mathrm{w}_{\#}\right)  .
\end{equation}
As the Hamiltonian describing noninteracting electrons is a one body
Hamiltonian the ground state density operator $D_{\#}\left(  \gamma
_{GS}\left(  \mathrm{v}\right)  \mathbf{,}\mathrm{w}_{\#}\right)  $ must be an
ensemble or a pure Independent Particle State (IPS)\ density operator. Thus
the space-spin density of the ground state of the interacting system in an
local external potential $\mathrm{v}$ is the same as space-spin density of the
ground state of a non interacting system determined by an environment
$\mathrm{w}_{\#}\mathrm{.}$

One can define a Kohn-Sham energy functional, $\mathcal{F}_{KS},$ for the
interacting system as
\begin{equation}
\mathcal{F}_{KS}\left(  \gamma\right)  =\mathcal{F}_{T}\left(  \gamma\right)
+\mathcal{F}_{C}\left(  \gamma\right)  +\mathcal{F}_{KSXC}\left(
\gamma\right)  ,
\end{equation}
where the Kohn-Sham exchange functional is defined as
\begin{equation}
\mathcal{F}_{KSXC}\left(  \gamma\right)  =\mathcal{F}_{XC}\left(
\gamma\right)  +\mathcal{F}_{K}\left(  \gamma\right)  -\mathcal{F}_{T}\left(
\gamma\right)  .
\end{equation}
The motivation for the introduction of $\mathcal{F}_{KS}$ is that one can
parameterize $\gamma$ in terms of $\left\{  \varphi_{\#j},N_{\#j}\right\}  $
(which are unique up to unitary transformation that mix orbitals that
correspond to degenerate occupation numbers ) through eq. $\left(
\ref{KEFunct}\right)  $ i.e.
\begin{equation}
\gamma\left(  \mathbf{x}\right)  =\sum_{1\leq j\leq r}\left\vert \varphi
_{\#j}\left(  \mathbf{x}\right)  \right\vert ^{2}N_{\#j}%
\end{equation}
where
\begin{subequations}
\label{Con}%
\begin{align}
&  \sum_{1\leq j\leq r}N_{\#j}=N\label{c1}\\
&  0\leq N_{\#j}\leq1;\quad1\leq j\leq r\label{c2}\\
&  \int\varphi_{\#j}\left(  \mathbf{x}\right)  \varphi_{\#k}\left(
\mathbf{x}\right)  ^{\ast}d\mathbf{x}=\delta_{jk};\quad1\leq j,k\leq
r\mathbf{.} \label{c3}%
\end{align}
The Kohn-Sham energy functional can then be expressed in terms of the
variables $\left\{  \left.  \varphi_{\#j},N_{\#j}\right\vert 1\leq j\leq
r\right\}  $ i.e.
\end{subequations}
\begin{equation}
\mathcal{F}_{KS}\left(  \mathbf{\varphi}_{\#}\mathbf{,N}_{\#}\right)
=\mathcal{F}_{KS}\left(  \gamma\left(  \mathbf{\varphi}_{\#}\mathbf{,N}%
_{\#}\right)  \right)
\end{equation}
and the minimization problem that determines the ground state energy and
space-spin density as
\begin{equation}
\min_{\left(  \mathbf{\varphi,N}\right)  \mathbf{\in}\mathcal{\Lambda}%
}\mathcal{E}_{KS}\left(  \mathbf{\varphi}_{\#}\mathbf{,N}_{\#}\mathbf{,}%
\mathrm{v}\right)  , \label{KSmin}%
\end{equation}
where
\begin{equation}
\mathcal{E}_{KS}\left(  \mathbf{\varphi}_{\#}\mathbf{,N}_{\#}\mathbf{,}%
\mathrm{v}\right)  =\mathcal{F}_{KS}\left(  \mathbf{\varphi}_{\#}%
\mathbf{,N}_{\#}\right)  +f_{\left(  \mathbf{\varphi}_{\#}\mathbf{,N}%
_{\#}\right)  }\left(  \mathrm{v}\right)
\end{equation}
and $\mathcal{\Lambda}$ is the set of feasible values for $\left(
\mathbf{\varphi}_{\#}\mathbf{,N}_{\#}\right)  $ given by eqs. $\left(
\ref{Con}\right)  .$ The Kohn-Sham constrained minimization problem can be
expressed in terms of the Lagrangian
\begin{align}
\mathcal{L}\left(  \mathbf{\varphi}_{\#}\mathbf{,N}_{\#},\mu
,\mathbf{\varepsilon,\eta,}\mathrm{v}\right)   &  =\mathcal{E}_{KS}\left(
\mathbf{\varphi}_{\#}\mathbf{,N}_{\#}\mathbf{,}\mathrm{v}\right)  +\sum_{1\leq
j\leq r}\eta_{j}\left(  1-N_{\#j}\right)  -\sum_{1\leq j\leq r}\xi_{j}%
N_{\#j}\\
&  -\sum_{1\leq j,k\leq r}\varepsilon_{jk}\left(  \left\langle \left.
\varphi_{\#j}\right\vert \varphi_{\#k}\right\rangle -\delta_{jk}\right)
+\mu\left(  \sum_{1\leq j\leq r}N_{\#j}-N\right)  .
\end{align}
This leads to the necessary stationarity conditions
\begin{subequations}
\begin{align}
\frac{\partial\mathcal{L}\left(  \mathbf{\varphi}_{\#}\mathbf{,N}_{\#}%
,\mu,\mathbf{\varepsilon,\eta,}\mathrm{v}\right)  }{\partial\varphi
_{\#j}^{\ast}}  &  =0\\
\frac{\partial\mathcal{L}\left(  \mathbf{\varphi}_{\#}\mathbf{,N}_{\#}%
,\mu,\mathbf{\varepsilon,\eta,}\mathrm{v}\right)  }{\partial\varphi_{\#j}}  &
=0\\
\frac{\partial\mathcal{L}\left(  \mathbf{\varphi}_{\#}\mathbf{,N}_{\#}%
,\mu,\mathbf{\varepsilon,\eta,}\mathrm{v}\right)  }{\partial N_{\#j}}  &  =0.
\end{align}
in particular implies that
\end{subequations}
\begin{equation}
\frac{\partial\mathcal{E}_{KS}\left(  \mathbf{\varphi}_{\#}\mathbf{,N}%
_{\#}\mathbf{,}\mathrm{v}\right)  }{\partial\varphi_{\#j}^{\ast}}-\sum_{1\leq
k\leq r}\varepsilon_{jk}\left\vert \varphi_{\#k}\right\rangle =0
\end{equation}
One can develop the Kohn-Sham energy functional as%
\begin{align}
&  \mathcal{E}_{KS}\left(  \mathbf{\varphi}_{\#}\mathbf{,N}_{\#}%
\mathbf{,}\mathrm{v}\right)  =\mathcal{F}_{KS}\left(  \mathbf{\varphi}%
_{\#}\mathbf{,N}_{\#}\right)  +f_{\left(  \mathbf{\varphi}_{\#}\mathbf{,N}%
_{\#}\right)  }\left(  \mathrm{v}\right) \nonumber\\
&  =\mathcal{F}_{T}\left(  \mathbf{\varphi}_{\#}\mathbf{,N}_{\#}\right)
+\mathcal{F}_{C}\left(  \mathbf{\varphi}_{\#}\mathbf{,N}_{\#}\right)
+\mathcal{F}_{KSXC}\left(  \mathbf{\varphi}_{\#}\mathbf{,N}_{\#}\right)
+f_{\left(  \mathbf{\varphi}_{\#}\mathbf{,N}_{\#}\right)  }\left(
\mathrm{v}\right) \nonumber\\
&  =\sum_{1\leq j\leq r}\left\langle \varphi_{\#j}\left\vert K\varphi
_{\#j}\right.  \right\rangle N_{\#j}+\int\frac{\gamma\left(  \mathbf{r,\xi
}\right)  \gamma\left(  \mathbf{r}^{\prime}\mathbf{,\xi}\right)  }{\left\Vert
\mathbf{r-r}^{\prime}\right\Vert }d\mathbf{r}d\mathbf{r}^{\prime}d\mathbf{\xi
}+\int\mathrm{v}\left(  \mathbf{r,\xi}\right)  \gamma\left(  \mathbf{r,\xi
}\right)  d\mathbf{r}d\mathbf{\xi}+\mathcal{F}_{KSXC}\left(  \mathbf{\varphi
}_{\#}\mathbf{,N}_{\#}\right)  ,
\end{align}
where the minimization over the kinetic functional gives the correspondence
\begin{equation}
\gamma\left(  \mathbf{x}\right)  =\sum_{1\leq j\leq r}\left\vert \varphi
_{\#j}\left(  \mathbf{x}\right)  \right\vert ^{2}N_{\#j}.
\end{equation}
Using
\begin{equation}
\frac{\partial}{\partial N_{\#j}}=\int\frac{\partial\gamma\left(
\mathbf{x}\right)  }{\partial N_{\#j}}\frac{\delta}{\delta\gamma\left(
\mathbf{x}\right)  }d\mathbf{x=}\int\varphi_{\#j}\left(  \mathbf{x}\right)
\varphi_{\#j}\left(  \mathbf{x}\right)  ^{\ast}\frac{\delta}{\delta
\gamma\left(  \mathbf{x}\right)  }d\mathbf{x} \label{DerNj}%
\end{equation}
one obtains Janaks Theorem
\begin{align}
\frac{\partial\mathcal{E}_{KS}\left(  \mathbf{\varphi}_{\#}\mathbf{,N}%
_{\#}\mathbf{,}\mathrm{v}\right)  }{\partial N_{\#j}}  &  =\left\langle
\varphi_{\#j}\left\vert K\varphi_{\#j}\right.  \right\rangle +\int\frac
{\delta}{\delta\gamma\left(  \mathbf{x}\right)  }\left\{  \mathcal{F}%
_{C}\left(  \mathbf{\varphi}_{\#}\mathbf{,N}_{\#}\right)  +\mathcal{F}%
_{KSXC}\left(  \mathbf{\varphi}_{\#}\mathbf{,N}_{\#}\right)  +f_{\left(
\mathbf{\varphi}_{\#}\mathbf{,N}_{\#}\right)  }\left(  \mathrm{v}\right)
\right\}  \frac{\partial\gamma\left(  \mathbf{x}\right)  }{\partial N_{\#j}%
}d\mathbf{x}\nonumber\\
&  =\left\langle \varphi_{\#j}\left\vert K\varphi_{\#j}\right.  \right\rangle
+\left\langle \varphi_{\#j}\left\vert \frac{\delta}{\delta\gamma\left(
\mathbf{x}\right)  }\left\{  \mathcal{F}_{C}\left(  \mathbf{\varphi}%
_{\#}\mathbf{,N}_{\#}\right)  +\mathcal{F}_{KSXC}\left(  \mathbf{\varphi}%
_{\#}\mathbf{,N}_{\#}\right)  +f_{\left(  \mathbf{\varphi}_{\#}\mathbf{,N}%
_{\#}\right)  }\left(  \mathrm{v}\right)  \right\}  \right.  \varphi
_{\#j}\right\rangle \nonumber\\
&  =\varepsilon_{jj}.
\end{align}


\section{AGP Energy Functionals\label{AGP}}

\subsection{Canonical expansion of an AGP state}

\bigskip The $2N$-electron AGP state, $\left\vert g^{N}\right\rangle $, is the
$N^{th}$ antisymmetric power of the $2$-electron state described by the
geminal
\begin{equation}
\left\vert g\right\rangle =\sum_{1\leq j\leq s}n_{j}^{\frac{1}{2}}\left\vert
\psi_{j}\psi_{j+s}\right\rangle , \label{geminal}%
\end{equation}
where $\left\{  \left.  n_{j}^{\frac{1}{2}}\right\vert 1\leq j\leq s\right\}
$ are the positive square roots of $\left\{  \left.  n_{j}\right\vert 1\leq
j\leq s\right\}  ,$ and is given by \cite{Weiner&ortiz04}
\begin{equation}
\left\vert g^{N}\right\rangle =\left(  S_{N}\right)  ^{-\frac{1}{2}}%
\sum_{1\leq j_{1}<\cdots<j_{N}\leq s}n_{j_{1}}^{\frac{1}{2}}\cdots n_{j_{s}%
}^{\frac{1}{2}}\left\vert \psi_{j_{1}}\psi_{j_{1}+s}\cdots\psi_{j_{N}}%
\psi_{j_{N}+s}\right\rangle . \label{agp}%
\end{equation}
The normalization factor, $\left(  S_{N}\right)  ^{-\frac{1}{2}}$, is given in
terms of the elementary symmetric function
\begin{equation}
S_{N}=\sum_{1\leq j_{1}<\cdots<j_{N}\leq s}n_{j_{1}}\cdots n_{j_{N}},
\end{equation}
and the geminals $\left\vert g\right\rangle $ represent normalized
two--electron states so that
\begin{equation}
\sum_{1\leq j\leq s}n_{j}=1.
\end{equation}
Note that
\begin{equation}
S_{0}=1.
\end{equation}
It is evident from eq.\ $\left(  \ref{geminal}\right)  $ that the geminal and
geminal power state are at the least invariant to the group
\begin{equation}
SU\left(  2\right)  ^{s}=\overset{s\text{-factors}}{\overbrace{\,SU\left(
2\right)  \times\cdots\times SU\left(  2\right)  }},
\end{equation}
that is formed by $2\times2$ \emph{special} unitary transformations of the
paired CGSOs $\left\{  \left.  \left\vert \psi_{j}\right\rangle ,\left\vert
\psi_{j+s}\right\rangle \right\vert 1\leq j\leq s\right\}  ,$ while if
degeneracies exist amongst the $\left\{  \left.  n_{j}\right\vert 1\leq j\leq
s\right\}  $ this invariance group is larger \cite{AGPCoherent}. The canonical
CSGOs are unique up to transformations belonging this invariance group.

\subsection{The AGP Energy Functional}

\section{Optimization Theory\label{Optimiz}}

\bigskip

\section{\bigskip AGP Lagrangian\label{Lagran}}

\section{\bigskip Generalized Janaks Theorem\label{General}}

\section{\bigskip Discussion\label{Discuss}}

The total energy can be expressed as
\begin{equation}
\mathcal{E}\left(  \mathbf{\varphi}\right)  =\frac{1}{2}\sum_{1\leq p\leq
2s}N\left(  \varphi_{p}\right)  \left\{  I\left(  \varphi_{p}\right)
+\left\langle \left.  \varphi_{p}\right\vert h_{1}\varphi_{p}\right\rangle
\right\}  =\frac{1}{2}%
%TCIMACRO{\tsum \limits_{1\leq p\leq2s}}%
%BeginExpansion
{\textstyle\sum\limits_{1\leq p\leq2s}}
%EndExpansion
N\left(  \varphi_{p}\right)  \kappa\left(  \varphi_{p}\right)
\end{equation}
where $\left\{  N\left(  \varphi_{p}\right)  \right\}  $ are the occupation
numbers of the orthogonal NGSO's $\left\{  \varphi_{p}\right\}  $ and the one
electron quantities $\left\{  \kappa\left(  \varphi_{p}\right)  \right\}  $
are
\begin{equation}
\kappa\left(  \varphi_{p}\right)  =I\left(  \varphi_{p}\right)  +\frac
{\left\langle \left.  \varphi_{p}\right\vert h_{1}\varphi_{p}\right\rangle
}{\left\langle \left.  \varphi_{p}\right\vert \varphi_{p}\right\rangle },
\end{equation}
in terms of the ionization potentials $\left\{  I\left(  \varphi_{p}\right)
\right\}  $ and the one electron Hamiltonian $h_{1}$
\begin{equation}
h_{1}=\sum_{1\leq j,k\leq2s}h_{1jk}\left(  \varphi\right)  \frac{\left\vert
\varphi_{j}\right\rangle \left\langle \varphi_{k}\right\vert }{\sqrt
{\left\langle \left.  \varphi_{j}\right\vert \varphi_{j}\right\rangle }%
\sqrt{\left\langle \left.  \varphi_{k}\right\vert \varphi_{k}\right\rangle }},
\end{equation}
where
\begin{align}
h_{1jk}\left(  \varphi\right)   &  =\frac{\left\langle \varphi_{j}\left\vert
h\right.  \varphi_{k}\right\rangle }{\sqrt{\left\langle \left.  \varphi
_{j}\right\vert \varphi_{j}\right\rangle }\sqrt{\left\langle \left.
\varphi_{k}\right\vert \varphi_{k}\right\rangle }}\label{h1}\\
&  =-\frac{1}{\sqrt{\left\langle \left.  \varphi_{j}\right\vert \varphi
_{j}\right\rangle }\sqrt{\left\langle \left.  \varphi_{k}\right\vert
\varphi_{k}\right\rangle }}\int\varphi_{j}^{\ast}\left(  \mathbf{r,\xi
}\right)  \left\{  \sum_{_{1\leq l\leq N_{nuc}}}\frac{Z_{l}}{\left\Vert
\mathbf{R}_{l}-\mathbf{r}\right\Vert }+\nabla_{\mathbf{r}}^{2}\right\}
\varphi_{k}\left(  \mathbf{r,\xi}\right)  d\mathbf{r}d\xi;1\leq j,k\leq2s.
\end{align}
Using the Fock operator the energy can be written as
\begin{equation}
\mathcal{E}\left(  \mathbb{D}^{2},\mathbf{\varphi}\right)  =\frac{1}%
{2}\left\{  \overset{}{\operatorname{Tr}\left\{  F\left(  \mathbb{D}%
^{2},\mathbf{\varphi}\right)  \right\}  +\mathcal{E}_{1}\left(  \mathbb{D}%
^{2},\mathbf{\varphi}\right)  }\right\}  ,
\end{equation}
where
\begin{equation}
\mathcal{E}_{1}\left(  \mathbb{D}^{2},\mathbf{\varphi}\right)
=\operatorname{Tr}\left\{  \mathbb{D}^{1}\left(  \mathbf{\varphi}\right)
\mathbb{H}_{1}\left(  \mathbf{\varphi}\right)  \right\}  .
\end{equation}
This gives
\begin{equation}
\mathcal{E}\left(  \mathbb{D}^{2},\mathbf{\varphi}\right)  =\frac{1}{2}%
\sum_{1\leq p\leq2s}\left\{  \epsilon_{p}+h_{1pp}\left(  \varphi\right)
N\left(  \varphi_{p}\right)  \right\}  ,
\end{equation}
where $\left\{  \epsilon_{p}\right\}  $ are the eigenvalues of the Fock
operator and $\left\{  \varphi_{p}\right\}  $ are NGSO's. One should note that
the GSO's that diagonalize the Fock operator are \emph{not }in general the NGSO's.

The energy of the AGP state is given by \emph{ }
\begin{align}
\mathcal{E}\left(  \mathbf{n,\varphi}\right)   &  =\frac{\mathcal{H}\left(
\mathbf{n,\varphi}\right)  }{S_{N}\left(  \mathbf{n}\right)  }\nonumber\\
&  =\frac{1}{S_{N}\left(  \mathbf{n}\right)  }\left\{  \sum_{1\leq j\leq
s}n_{j}S_{N-1}\left(  \jmath\right)  \left\{  \left\langle \varphi
_{j}\left\vert h_{1}\varphi_{j}\right.  \right\rangle +\left\langle
\varphi_{j+s}\left\vert h_{1}\varphi_{j+s}\right.  \right\rangle +\left\langle
\varphi_{j}\varphi_{j+s}||\varphi_{j}\varphi_{j+s}\right\rangle \right\}
\right. \nonumber\\
&  +\sum_{1\leq j,k\leq s}n_{j}^{\frac{1}{2}}n_{k}^{\frac{1}{2}}S_{N-1}\left(
\jmath k\right)  \left\langle \varphi_{j}\varphi_{j+s}||\varphi_{k}%
\varphi_{k+s}\right\rangle \left.  +\sum_{1\leq j<k\leq s}n_{i}n_{j}%
S_{N-2}\left(  \jmath k\right)  \left\langle \varphi_{j}\varphi_{k}%
|||\varphi_{j}\varphi_{k}\right\rangle \right\}  \label{energy}%
\end{align}
in terms of geminal occupation numbers and orthonormal CGSO's, while in terms
of the one electron group parameters it is
\begin{equation}
\mathcal{E}\left(  \mathbf{\eta},\mathbf{\lambda}\right)  =\frac
{\mathcal{H}\left(  \mathbf{\eta},\mathbf{\lambda}\right)  }{S_{N}\left(
\mathbf{\eta}\right)  }=\frac{\left\langle \left.  g\left(  \mathbf{\eta
},\mathbf{\lambda}\right)  ^{N}\right\vert Hg\left(  \mathbf{\eta
},\mathbf{\lambda}\right)  ^{N}\right\rangle }{S_{N}\left(  \mathbf{\eta
}\right)  }, \label{gp_funct}%
\end{equation}
where%
\begin{equation}
S_{N}\left(  \mathbf{\eta}\right)  =\left\langle \left.  k\left(
\mathbf{\eta}\right)  g_{0}^{N}\right\vert k\left(  \mathbf{\eta}\right)
g_{0}^{N}\right\rangle =\sum_{1\leq j_{1}<\cdots<j_{N}\leq s}e^{2\left(
\eta_{j_{1}}+\cdots+\eta_{j_{N}}\right)  }n_{0j_{1}}\cdots n_{0j_{N}}.
\end{equation}
One can define a Lagrangian
\begin{equation}
\mathcal{L}\left(  \mathbf{n,\varphi}\right)  =\mathcal{E}\left(
\mathbf{n,\varphi}\right)  -\sum_{1\leq j,k\leq2s}\varepsilon_{kj}\left(
\left\langle \left.  \varphi_{j}\right\vert \varphi_{k}\right\rangle
-\delta_{jk}\right)
\end{equation}
associated with the constrained variation (orthogonal with a fixed
normalization) of energy function $\mathcal{E}\left(  \mathbf{n,\varphi
}\right)  .$ At constrained stationary points
\begin{equation}
\frac{\partial\mathcal{L}\left(  \mathbf{n,\varphi}\right)  }{\partial n_{j}%
}=\frac{\partial\mathcal{E}\left(  \mathbf{n,\varphi}\right)  }{\partial
n_{j}}.
\end{equation}
This gives
\begin{align}
&  \frac{\partial\mathcal{L}\left(  \mathbf{n,\varphi}\right)  }{\partial
N_{j}}=\sum_{1\leq k\leq2s}\frac{\partial\mathcal{L}\left(  \mathbf{n,\varphi
}\right)  }{\partial n_{k}}\frac{\partial n_{k}}{\partial N_{j}}\\
&  \Rightarrow\frac{\partial\mathcal{L}\left(  \mathbf{n,\varphi}\right)
}{\partial N_{j}}=\sum_{1\leq k\leq2s}\frac{\partial\mathcal{L}\left(
\mathbf{n,\varphi}\right)  }{\partial n_{k}}\frac{\partial n_{k}}{\partial
N_{j}}=\sum_{1\leq k\leq2s}\frac{\partial\mathcal{L}\left(  \mathbf{n,\varphi
}\right)  }{\partial n_{k}}\left(  \frac{\partial N_{k}}{\partial n_{j}%
}\right)  ^{-1}%
\end{align}
i.e.%
\begin{equation}
\frac{\partial\mathcal{L}\left(  \mathbf{n,\varphi}\right)  }{\partial
\mathbf{N}}=\frac{\partial\mathcal{L}\left(  \mathbf{n,\varphi}\right)
}{\partial\mathbf{N}}=\frac{\partial\mathcal{L}\left(  \mathbf{n,\varphi
}\right)  }{\partial\mathbf{n}}\left(  \frac{\partial\mathbf{N}}%
{\partial\mathbf{n}}\right)  ^{-1}.
\end{equation}
At a constrained stationary point
\begin{equation}
\frac{\partial\mathcal{L}\left(  \mathbf{n,\varphi}\right)  }{\partial
\mathbf{N}}=\frac{\partial\mathcal{L}\left(  \mathbf{n,\varphi}\right)
}{\partial\mathbf{n}}=\mathbf{0.}%
\end{equation}
The occupation numbers of the AGP\ state are
\begin{align}
N_{j}  &  =\mathcal{N}_{j}\left(  \mathbf{n}\right)  =n_{j}\frac
{S_{N-1}\left(  \jmath\right)  }{S_{N}}=n_{j}\frac{\partial\ln S_{N}}{\partial
n_{j}};\quad1\leq j\leq s,\\
\mathbf{N}  &  =\mathcal{N}\left(  \mathbf{n}\right)  .
\end{align}%
\begin{equation}
\frac{\partial S_{N}}{\partial n_{j}}=S_{N-1}\left(  j\right)
\end{equation}%
\begin{align}
\sum_{j}n_{j}S_{N-1}\left(  j\right)   &  =NS_{N}\\
\sum_{k}n_{k}S_{N-2}\left(  jk\right)   &  =\left(  N-1\right)  S_{N-1}\left(
j\right) \\
\sum_{j<k}n_{j}n_{k}S_{N-2}\left(  jk\right)   &  =\frac{N\left(  N-1\right)
}{2}S_{N}%
\end{align}
The Jacobian is given by
\begin{equation}
\frac{\partial\mathcal{L}\left(  \mathbf{n,\varphi}\right)  }{\partial n_{j}%
}=\sum_{k}\frac{\partial\mathcal{L}\left(  \mathbf{n,\varphi}\right)
}{\partial N_{k}}\frac{\partial N_{k}}{\partial n_{j}}=\sum_{k}\frac
{\partial\mathcal{L}\left(  \mathbf{n,\varphi}\right)  }{\partial N_{k}}J_{kj}%
\end{equation}
The matrix elements of the Jacobian are
\begin{align}
\frac{\partial N_{k}}{\partial n_{j}}  &  =\frac{\partial}{\partial n_{j}%
}\left\{  n_{k}\frac{\partial\ln S_{N}}{\partial n_{k}}\right\}  =\delta
_{jk}\frac{\partial\ln S_{N}}{\partial n_{k}}+n_{k}\frac{\partial^{2}\ln
S_{N}}{\partial n_{j}\partial n_{k}}\\
&  =\delta_{jk}\frac{S_{N-1}\left(  \jmath\right)  }{S_{N}}-n_{k}\frac
{S_{N-1}\left(  \jmath\right)  S_{N-1}\left(  k\right)  }{S_{N}^{2}}%
+n_{k}\frac{S_{N-2}\left(  jk\right)  }{S_{N}}.\nonumber
\end{align}
For a $2N$-electron IPS the geminal has the form
\begin{equation}
\left\vert g\right\rangle =N^{-\frac{1}{2}}\sum_{1\leq j\leq N}\left\vert
\psi_{j}\psi_{j+s}\right\rangle
\end{equation}
which for a $8$-electron case is
\begin{equation}
\left\vert g\right\rangle =4^{-\frac{1}{2}}\sum_{1\leq j\leq4}\left\vert
\psi_{j}\psi_{j+s}\right\rangle =\frac{1}{2}\sum_{1\leq j\leq4}\left\vert
\psi_{j}\psi_{j+s}\right\rangle
\end{equation}
so that
\begin{equation}
n_{j}=\left\{
\begin{array}
[c]{c}%
N^{-1};1\leq j\leq N\\
\qquad0;\text{ }N+1\leq j\leq s
\end{array}
\right.
\end{equation}
which for a $6$-electron case is%
\begin{equation}
n_{j}=\left\{
\begin{array}
[c]{c}%
\frac{1}{4};1\leq j\leq4\\
0;\text{ }5\leq j\leq s
\end{array}
\right.
\end{equation}%
\begin{equation}
S_{N}=N^{-N}%
\end{equation}
which for a $8$-electron case is%
\begin{equation}
S_{4}=\frac{1}{256}.
\end{equation}
The first derivatives of the symmetric function are
\begin{equation}
S_{N-1}\left(  \jmath\right)  =\sum_{\substack{j\notin\left\{  j_{1}%
,\cdots,j_{N-1}\right\}  \\1\leq j_{1}<\cdots<j_{N-1}\leq N}}n_{j_{1}}\cdots
n_{j_{N-1}}=\left\{
\begin{array}
[c]{c}%
\binom{N}{N-1}N^{-\left(  N-1\right)  };\quad j\notin\left\{  1,\cdots
,N\right\} \\
\binom{N-1}{N-1}N^{-\left(  N-1\right)  };\quad j\in\left\{  1,\cdots
,N\right\}
\end{array}
\right.
\end{equation}
which for a $8$-electron $s=8$ case where $\left(  n_{1},n_{2},n_{3}%
,n_{4}\right)  =\left(  \frac{1}{4},\frac{1}{4},\frac{1}{4},\frac{1}%
{4}\right)  $ is%

\begin{align}
S_{3}\left(  8\right)   &  =n_{1}n_{2}n_{3}+n_{1}n_{2}n_{4}+n_{1}n_{2}%
n_{5}+n_{1}n_{2}n_{6}+n_{1}n_{2}n_{7}\\
&  +n_{1}n_{3}n_{4}+n_{1}n_{3}n_{5}+n_{1}n_{3}n_{6}+n_{1}n_{3}n_{7}\\
&  +n_{1}n_{4}n_{5}+n_{1}n_{4}n_{6}+n_{1}n_{4}n_{7}\\
&  +n_{1}n_{5}n_{6}+n_{1}n_{5}n_{7}+n_{1}n_{6}n_{7}\\
&  +n_{2}n_{3}n_{4}+n_{2}n_{3}n_{5}+n_{2}n_{3}n_{6}+n_{2}n_{3}n_{7}\\
&  +n_{2}n_{4}n_{5}+n_{2}n_{4}n_{6}+n_{2}n_{4}n_{7}\\
&  +n_{2}n_{5}n_{6}+n_{2}n_{5}n_{7}+n_{2}n_{6}n_{7}\\
&  +n_{3}n_{4}n_{5}+n_{3}n_{4}n_{6}+n_{3}n_{4}n_{7}\\
&  +n_{3}n_{5}n_{6}+n_{3}n_{5}n_{7}+n_{3}n_{6}n_{7}\\
&  +n_{4}n_{5}n_{6}+n_{4}n_{5}n_{7}+n_{4}n_{6}n_{7}+n_{5}n_{6}n_{7}%
\end{align}%
\begin{equation}
S_{4}\left(  \jmath\right)  =\sum_{\substack{j\notin\left\{  j_{1},j_{2}%
,j_{3},j_{4}\right\}  \\1\leq j_{1}<j_{2}<j_{3}<j_{4}\leq8}}n_{j_{1}}n_{j_{2}%
}n_{j_{3}}n_{j_{4}}=\left\{
\begin{array}
[c]{c}%
\frac{1}{16};\quad j\notin\left\{  1,2,3,4\right\} \\
\frac{1}{64};\quad j\in\left\{  1,2,3,4\right\}
\end{array}
\right.
\end{equation}%
\begin{equation}
S_{N-2}\left(  jk\right)  =\sum_{\substack{j,k\notin\left\{  j_{1}%
,\cdots,j_{N-2}\right\}  \\1\leq j_{1}<\cdots<j_{N-2}\leq N}}n_{j_{1}}\cdots
n_{j_{N-2}}=\left\{
\begin{array}
[c]{c}%
\binom{N}{N-2}N^{-\left(  N-2\right)  };\quad j,k\notin\left\{  1,\cdots
,N\right\} \\
\binom{N-1}{N-2}N^{-\left(  N-2\right)  };\quad j\notin\left\{  1,\cdots
,N\right\}  ,k\in\left\{  1,\cdots,N\right\} \\
\binom{N-1}{N-2}N^{-\left(  N-2\right)  };\quad k\notin\left\{  1,\cdots
,N\right\}  ,j\in\left\{  1,\cdots,N\right\} \\
\binom{N-2}{N-2}N^{-\left(  N-2\right)  };\quad j\in\left\{  1,\cdots
,N\right\}  ,k\in\left\{  1,\cdots,N\right\}
\end{array}
\right.
\end{equation}
which for a $8$-electron case is
\begin{align}
S_{2}\left(  7,8\right)  =n_{1}n_{2}  &  +n_{1}n_{3}+n_{1}n_{4}+n_{1}%
n_{5}+n_{1}n_{6}+n_{2}n_{3}+n_{2}n_{4}+n_{2}n_{5}\\
&  +n_{2}n_{6}+n_{3}n_{4}+n_{3}n_{5}+n_{3}n_{6}+n_{4}n_{5}+n_{4}n_{6}%
+n_{5}n_{6}%
\end{align}
\begin{equation}
S_{2}\left(  jk\right)  =\sum_{\substack{j,k\notin\left\{  j_{1}%
,\cdots,j_{N-2}\right\}  \\1\leq j_{1}<\cdots<j_{N-2}\leq N}}n_{j_{1}}\cdots
n_{j_{N-2}}=\left\{
\begin{array}
[c]{c}%
\frac{3}{8};\quad j,k\notin\left\{  1,\cdots,N\right\}  \qquad\qquad\qquad\\
\frac{3}{16};\quad j\notin\left\{  1,\cdots,N\right\}  ,k\in\left\{
1,\cdots,N\right\} \\
\frac{3}{16};\quad k\notin\left\{  1,\cdots,N\right\}  ,j\in\left\{
1,\cdots,N\right\} \\
\frac{1}{16};\quad j\in\left\{  1,\cdots,N\right\}  ,k\in\left\{
1,\cdots,N\right\}
\end{array}
\right.
\end{equation}


The diagonal
\begin{equation}
\frac{\partial N_{k}}{\partial n_{k}}=\left\{
\begin{array}
[c]{c}%
N;\quad k\notin\left\{  1,\cdots,N\right\} \\
0;\quad k\in\left\{  1,\cdots,N\right\}
\end{array}
\right.
\end{equation}%
\begin{equation}
\frac{\partial N_{k}}{\partial n_{j}}=\left\{
\begin{array}
[c]{c}%
-n_{k}\frac{N\left(  N+1\right)  }{2};\quad j\neq k\notin\left\{
1,\cdots,N\right\} \\
n_{k};\quad j\notin\left\{  1,\cdots,N\right\}  ,k\in\left\{  1,\cdots
,N\right\} \\
n_{k};\quad k\notin\left\{  1,\cdots,N\right\}  ,j\in\left\{  1,\cdots
,N\right\} \\
0;\quad j\in\left\{  1,\cdots,N\right\}  ,k\in\left\{  1,\cdots,N\right\}
\end{array}
\right.
\end{equation}


The occupation number variation nonunitary subgroup $G_{num}$ is given by
\begin{equation}
G_{num}=\left\{  \left.  k\left(  \mathbf{\eta}\right)  =\exp\left\{
\sum_{1\leq j\leq s}\eta_{j}\Omega_{j}\right\}  \right\vert \eta_{j}%
\in\mathbb{R}\right\}  , \label{occup}%
\end{equation}
where
\begin{equation}
\Omega_{j}=\left(  a_{j}^{\dagger}a_{j}+a_{j+s}^{\dagger}a_{j+s}\right)
;\quad1\leq j\leq s.
\end{equation}
The manifold of $2N$-electron AGP\ states $\mathcal{M}$ can be expressed as
\begin{equation}
\mathcal{M}\mathcal{=}\left\{  \left.  \overset{}{k\left(  \mathbf{\eta
}\right)  u\left(  \mathbf{\lambda}\right)  }\left\vert g_{0}^{N}\right\rangle
\right\vert \right.  \left.  \mathbf{\lambda}=\mathbf{\lambda}^{\dagger
}\right\}  .
\end{equation}
Explicitly the action on a reference geminal is given by
\begin{equation}
\left\vert \tilde{g}\left(  \mathbf{\eta},\mathbf{\lambda}\right)
\right\rangle =k\left(  \mathbf{\eta}\right)  u\left(  \mathbf{\lambda
}\right)  \left\vert g_{0}\right\rangle =\sum_{1\leq l\leq s}e^{\eta_{l}%
}c_{0l}\left\vert \varphi_{l}\left(  \mathbf{\lambda}\right)  \varphi
_{l+s}\left(  \mathbf{\lambda}\right)  \right\rangle ,
\end{equation}
where
\begin{equation}
\left\vert \varphi_{l}\left(  \mathbf{\lambda}\right)  \right\rangle
=\exp\left\{  i\sum_{1\leq j,k\leq2s}\lambda_{jk}a_{j}^{\dagger}a_{k}\right\}
\left\vert \varphi_{0l}\right\rangle .
\end{equation}
One can define a normalized geminal by
\begin{equation}
\left\vert g\left(  \mathbf{\eta},\mathbf{\lambda}\right)  \right\rangle
=\mathit{s}\left(  \mathbf{\eta}\right)  ^{-\frac{1}{2}}\left\vert \tilde
{g}\left(  \mathbf{\eta},\mathbf{\lambda}\right)  \right\rangle ,
\end{equation}
where
\begin{equation}
\mathit{s}\left(  \mathbf{\eta}\right)  =\left\langle \left.  \tilde{g}\left(
\mathbf{\eta},\mathbf{\lambda}\right)  \right\vert \tilde{g}\left(
\mathbf{\eta},\mathbf{\lambda}\right)  \right\rangle =\sum_{1\leq l\leq
s}e^{2\eta_{l}}c_{0l}^{2}=\sum_{1\leq l\leq s}e^{2\eta_{l}}n_{0l}%
\end{equation}
giving
\begin{align}
\left\vert g\left(  \mathbf{\eta},\mathbf{\lambda}\right)  \right\rangle  &
=\mathit{s}\left(  \mathbf{\eta}\right)  ^{-\frac{1}{2}}\sum_{1\leq l\leq
s}e^{\eta_{l}}c_{0l}\left\vert \varphi_{l}\left(  \mathbf{\lambda}\right)
\varphi_{l+s}\left(  \mathbf{\lambda}\right)  \right\rangle \\
&  =\sum_{1\leq l\leq s}e^{\eta_{l}}c_{0l}\mathit{s}\left(  \mathbf{\eta
}\right)  ^{-\frac{1}{2}}\left\vert \varphi_{l}\left(  \mathbf{\lambda
}\right)  \varphi_{l+s}\left(  \mathbf{\lambda}\right)  \right\rangle
=\sum_{1\leq l\leq s}c_{l}\left(  \mathbf{\eta}\right)  \left\vert \varphi
_{l}\left(  \mathbf{\lambda}\right)  \varphi_{l+s}\left(  \mathbf{\lambda
}\right)  \right\rangle ,
\end{align}
where
\begin{equation}
c_{l}\left(  \mathbf{\eta}\right)  =e^{\eta_{l}}c_{0l}\mathit{s}\left(
\mathbf{\eta}\right)  ^{-\frac{1}{2}}%
\end{equation}
leading to
\begin{equation}
n_{l}\left(  \mathbf{\eta}\right)  =e^{2\eta_{l}}n_{0l}\mathit{s}\left(
\mathbf{\eta}\right)  ^{-1}.
\end{equation}
This gives
\begin{equation}
\sum_{1\leq l\leq s}n_{l}\left(  \mathbf{\eta}\right)  =\mathit{s}\left(
\mathbf{\eta}\right)  ^{-1}\sum_{1\leq l\leq s}e^{2\eta_{l}}n_{0l}=\left(
\sum_{1\leq l\leq s}e^{2\eta_{l}}n_{0l}\right)  ^{-1}\sum_{1\leq l\leq
s}e^{2\eta_{l}}n_{0l}=1
\end{equation}
which shows that $\left\{  n_{l}\left(  \mathbf{\eta}\right)  \right\}  $
satisfy the geminal occupation number positivity and normalization
constraints. We can define%
\begin{equation}
\left\vert \varphi_{l}\left(  \mathbf{\eta,\lambda}\right)  \right\rangle
=\exp\eta_{l}\exp\left\{  i\sum_{1\leq j,k\leq2s}\lambda_{jk}a_{j}^{\dagger
}a_{k}\right\}  \left\vert \varphi_{0l}\right\rangle ,
\end{equation}
so that
\begin{equation}
\left\langle \left.  \varphi_{l}\left(  \mathbf{\eta,\lambda}\right)
\right\vert \varphi_{l}\left(  \mathbf{\eta,\lambda}\right)  \right\rangle
=e^{2\eta_{l}}\left\langle \left.  \varphi_{0l}\right\vert \varphi
_{0l}\right\rangle =e^{2\eta_{l}}%
\end{equation}
and
\begin{equation}
\left\vert g\left(  \mathbf{\eta},\mathbf{\lambda}\right)  \right\rangle
=\mathit{s}\left(  \mathbf{\eta}\right)  ^{-\frac{1}{2}}\sum_{1\leq l\leq
s}c_{0l}\left\vert \varphi_{l}\left(  \mathbf{\eta,\lambda}\right)
\varphi_{l+s}\left(  \mathbf{\eta,\lambda}\right)  \right\rangle .
\end{equation}


We can express the geminal occupation numbers (and hence the occupation
numbers of the FORDM) and the normalization of the CGSO's in terms of
$\mathbf{\eta}$
\begin{align}
n_{l}\left(  \mathbf{\eta}\right)   &  =e^{2\eta_{l}}n_{0l}\mathit{s}\left(
\mathbf{\eta}\right)  ^{-1}\nonumber\\
s_{l}  &  =\left\langle \left.  \varphi_{l}\left(  \mathbf{\eta,\lambda
}\right)  \right\vert \varphi_{l}\left(  \mathbf{\eta,\lambda}\right)
\right\rangle =e^{2\eta_{l}}%
\end{align}
giving
\begin{equation}
\eta_{l}=\frac{1}{2}\ln\left\langle \left.  \varphi_{l}\left(  \mathbf{\eta
,\lambda}\right)  \right\vert \varphi_{l}\left(  \mathbf{\eta,\lambda}\right)
\right\rangle \equiv\frac{1}{2}\ln\left\langle \left.  \varphi_{l}\left(
\mathbf{\lambda}\right)  \right\vert \varphi_{l}\left(  \mathbf{\lambda
}\right)  \right\rangle =\frac{1}{2}\ln s_{l}%
\end{equation}
Hence%
\begin{equation}
\mathcal{E}\left(  \mathbf{n}_{\min}\mathbf{,\varphi}_{\min}\right)
=\mathcal{E}\left(  \mathbf{\eta}_{\min}\mathbf{,\lambda}_{\min}\right)
=\mathcal{E}\left(  \mathbf{\varphi}_{\min}\right)
\end{equation}
where
\begin{align}
\mathcal{E}\left(  \mathbf{n}_{\min}\mathbf{,\varphi}_{\min}\right)   &
=\min_{\mathbf{n\in\mathcal{S},\varphi\in}\mathcal{O}\left(  \mathbf{1}%
_{s}\right)  }\mathcal{E}\left(  \mathbf{n,\varphi}\right) \\
\mathcal{E}\left(  \mathbf{\eta}_{\min}\mathbf{,\lambda}_{\min}\right)   &
=\min_{\mathbf{\eta,\lambda}}\mathcal{E}\left(  \mathbf{\eta,\lambda}\right)
\\
\mathcal{E}\left(  \mathbf{\varphi}_{\min}\right)   &  =\min_{\mathbf{s\in
}\mathbb{R}_{+}^{s}}\min_{\mathbf{\varphi\in}\mathcal{O}\left(  \mathbf{s}%
\right)  }\mathcal{E}\left(  \mathbf{\varphi}\right)
\end{align}
and
\begin{align}
\mathcal{S}  &  =\left\{  \left.  \mathbf{n}\right\vert n_{j}\geq0,1\leq j\leq
s;\sum_{1\leq j\leq s}n_{j}=1\right\} \\
\mathcal{O}\left(  \mathbf{s}\right)   &  =\left\{  \left.  \mathbf{\varphi
}\right\vert \left\langle \left.  \varphi_{j}\right\vert \varphi
_{k}\right\rangle =s_{j}\delta_{jk},1\leq j,k\leq s\right\}  .
\end{align}%
\begin{equation}
\mathbf{\lambda}=\mathbf{\lambda}^{\dagger},\mathbf{\eta\in}\mathbb{R}^{s}%
\end{equation}%
\begin{align}
\mathcal{L}\left(  \mathbf{\varphi,s}\right)   &  =\mathcal{E}\left(
\mathbf{n}_{0}\mathbf{,\varphi}\right)  -\sum_{1\leq j,k\leq2s}\varepsilon
_{kj}\left(  \left\langle \left.  \varphi_{j}\right\vert \varphi
_{k}\right\rangle -s_{j}\delta_{jk}\right) \\
\frac{\partial\mathcal{L}\left(  \mathbf{\varphi,s}\right)  }{\partial s_{j}}
&  =\varepsilon_{jj}.
\end{align}%
\begin{equation}
\mathcal{L}\left(  \mathbf{n\left(  \mathbf{\eta}\right)  ,\varphi}\right)
=\mathcal{L}\left(  \mathbf{\varphi,s}\left(  \mathbf{\eta}\right)  \right)
\end{equation}


\section{Sensitivity}

Consider the optimization problem
\begin{equation}
\min_{\mathbf{x\in}\mathcal{D}\left(  \mathbf{c}\right)  }f\left(
\mathbf{x}\right)  ,
\end{equation}
where
\begin{equation}
\mathcal{D}\left(  \mathbf{c}\right)  \mathcal{=}\left\{  \mathbf{x}\left\vert
g\left(  \mathbf{x}\right)  =\mathbf{c}\right.  \right\}  .
\end{equation}
This can be solved using the Lagrangian
\begin{equation}
\mathcal{L}\left(  \mathbf{x,\lambda}\right)  =f\left(  \mathbf{x}\right)
-\mathbf{\lambda}^{t}\left(  g\left(  \mathbf{x}\right)  -\mathbf{c}\right)  .
\end{equation}
$\left\{  \mathbf{x}_{m}\left(  \mathbf{\lambda}\left(  \mathbf{c}\right)
\right)  \right\}  $ are constrained stationary points if
\begin{equation}
\nabla_{\mathbf{x}}\mathcal{L}\left(  \mathbf{x}_{m}\left(  \mathbf{\lambda
}\left(  \mathbf{c}\right)  \right)  \mathbf{,\lambda}\left(  \mathbf{c}%
\right)  \right)  =\nabla_{\mathbf{x}}f\left(  \mathbf{x}_{m}\left(
\mathbf{\lambda}\left(  \mathbf{c}\right)  \right)  \right)  -\mathbf{\lambda
}\left(  \mathbf{c}\right)  ^{t}\nabla_{\mathbf{x}}g\left(  \mathbf{x}%
_{m}\left(  \mathbf{\lambda}\left(  \mathbf{c}\right)  \right)  \right)
=\mathbf{0} \label{sc1}%
\end{equation}
and
\begin{equation}
g\left(  \mathbf{x}_{m}\left(  \mathbf{\lambda}\left(  \mathbf{c}\right)
\right)  \right)  -\mathbf{c=0}%
\end{equation}
for
\begin{equation}
\mathbf{\lambda}\left(  \mathbf{c}\right)  \mathbf{\in\Lambda}\left(
\mathbf{c}\right)  ,
\end{equation}
where $\mathbf{\Lambda}\left(  \mathbf{c}\right)  $ is the set of Lagrange
parameters for which eq. $\left(  \ref{sc1}\right)  $ has solutions. $\lceil
$The procedure is: find a $\lambda\left(  \mathbf{c}\right)  $ then a set of
solutions $\Gamma\left(  \lambda\left(  \mathbf{c}\right)  \right)  =\left\{
\mathbf{x}_{m}\left(  \lambda\left(  \mathbf{c}\right)  \right)  \right\}  $
for that $\lambda\left(  \mathbf{c}\right)  \rfloor$. The sets
$\mathbf{\Lambda}\left(  \mathbf{c}\right)  $ and $\Gamma\left(
\lambda\left(  \mathbf{c}\right)  \right)  $ may be empty.

The stationary points are minima if
\begin{equation}
\nabla_{\mathbf{x}}^{2}f\left(  \mathbf{x}_{m}\right)  \geq0\text{ in the
feasible subspace}%
\end{equation}
they are non degenerate minima if $\nabla_{\mathbf{x}}^{2}f\left(
\mathbf{x}_{m}\right)  $ is positive definite, but degenerate minima if
$\nabla_{\mathbf{x}}^{2}f\left(  \mathbf{x}_{m}\right)  $ is only positive.

\subsection{Dual Problem}

\bigskip One can consider the Lagrangian $\mathcal{L}\left(  \mathbf{x,\lambda
}\right)  $ to be function of $\mathbf{x}$ and $\mathbf{\lambda}$ then the
stationarity conditions are
\begin{subequations}
\begin{align}
\nabla_{\mathbf{x}}\mathcal{L}\left(  \mathbf{x,\lambda}\right)   &
=\nabla_{\mathbf{x}}f\left(  \mathbf{x}\right)  -\mathbf{\lambda}^{t}%
\nabla_{\mathbf{x}}g\left(  \mathbf{x}\right)  =\mathbf{0}\label{d1}\\
\nabla_{\mathbf{\lambda}}\mathcal{L}\left(  \mathbf{x,\lambda}\right)   &
=-\left(  g\left(  \mathbf{x}\right)  -\mathbf{c}\right)  =\mathbf{0}
\label{d2}%
\end{align}
one should note that the condition eq. $\left(  \ref{d2}\right)  $ is the
constraint. We denote solutions to these equations as $\left\{  \left(
\mathbf{x}_{s}\left(  \mathbf{c}\right)  \mathbf{,\lambda}_{s}\left(
\mathbf{c}\right)  \right)  \right\}  $. The Hessian of the Lagrangian is
\end{subequations}
\begin{equation}
\left(
\begin{array}
[c]{cc}%
\nabla_{\mathbf{x}}^{2}\mathcal{L}\left(  \mathbf{x,\lambda}\right)  &
\nabla_{\mathbf{x\lambda}}^{2}\mathcal{L}\left(  \mathbf{x,\lambda}\right) \\
\nabla_{\mathbf{\lambda x}}^{2}\mathcal{L}\left(  \mathbf{x,\lambda}\right)  &
\nabla_{\mathbf{\lambda}}^{2}\mathcal{L}\left(  \mathbf{x,\lambda}\right)
\end{array}
\right)  ,
\end{equation}
where
\begin{align}
\nabla_{\mathbf{x\lambda}}^{2}\mathcal{L}\left(  \mathbf{x,\lambda}\right)
&  =-\nabla_{\mathbf{x}}g\left(  \mathbf{x}\right) \nonumber\\
\nabla_{\mathbf{\lambda x}}^{2}\mathcal{L}\left(  \mathbf{x,\lambda}\right)
&  =-\nabla_{\mathbf{x}}g\left(  \mathbf{x}\right)  ^{t}%
\end{align}
and
\begin{equation}
\nabla_{\mathbf{\lambda}}^{2}\mathcal{L}\left(  \mathbf{x,\lambda}\right)
=\mathbb{O}.
\end{equation}
If $\mathbf{x}_{m}\left(  \mathbf{c}\right)  $ is a local minimum then%
\begin{align}
&  \left(
\begin{array}
[c]{cc}%
\delta\mathbf{x} & \mathbf{0}%
\end{array}
\right)  \left(
\begin{array}
[c]{cc}%
\nabla_{\mathbf{x}}^{2}\mathcal{L}\left(  \mathbf{x}_{m}\left(  \mathbf{c}%
\right)  \mathbf{,\lambda}_{m}\left(  \mathbf{c}\right)  \right)  &
-\nabla_{\mathbf{x}}g\left(  \mathbf{x}_{m}\left(  \mathbf{c}\right)  \right)
\\
-\nabla_{\mathbf{x}}g\left(  \mathbf{x}_{m}\left(  \mathbf{c}\right)  \right)
^{t} & \mathbb{O}%
\end{array}
\right)  \left(
\begin{array}
[c]{c}%
\delta\mathbf{x}\\
\mathbf{0}%
\end{array}
\right)  \geq0\forall\delta\mathbf{x\backepsilon}\delta\mathbf{x}^{t}%
\nabla_{\mathbf{x}}g\left(  \mathbf{x}_{m}\left(  \mathbf{c}\right)  \right)
=0\nonumber\\
&  \qquad\qquad\qquad\qquad\qquad\qquad\qquad\qquad\qquad\Updownarrow
\nonumber\\
&  \qquad\qquad\qquad\qquad\qquad\qquad\delta\mathbf{x}^{t}\nabla_{\mathbf{x}%
}^{2}\mathcal{L}\left(  \mathbf{x}_{m}\left(  \mathbf{c}\right)
\mathbf{,\lambda}_{m}\left(  \mathbf{c}\right)  \right)  \delta\mathbf{x}%
\geq0\mathbf{\backepsilon}\delta\mathbf{x}^{t}\nabla_{\mathbf{x}}g\left(
\mathbf{x}_{m}\left(  \mathbf{c}\right)  \right)  =0
\end{align}
and $\mathbf{\lambda}_{m}\left(  \mathbf{c}\right)  $ is a local maximum
\begin{align}
&  \left(
\begin{array}
[c]{cc}%
\mathbf{0} & \delta\mathbf{\lambda}%
\end{array}
\right)  \left(
\begin{array}
[c]{cc}%
\nabla_{\mathbf{x}}^{2}\mathcal{L}\left(  \mathbf{x}_{m}\left(  \mathbf{c}%
\right)  \mathbf{,\lambda}_{m}\left(  \mathbf{c}\right)  \right)  &
\nabla_{\mathbf{x\lambda}}^{2}\mathcal{L}\left(  \mathbf{x}_{m}\left(
\mathbf{c}\right)  \mathbf{,\lambda}_{m}\left(  \mathbf{c}\right)  \right) \\
\nabla_{\mathbf{\lambda x}}^{2}\mathcal{L}\left(  \mathbf{x}_{m}\left(
\mathbf{c}\right)  \mathbf{,\lambda}_{m}\left(  \mathbf{c}\right)  \right)  &
\mathbb{O}%
\end{array}
\right)  \left(
\begin{array}
[c]{c}%
\mathbf{0}\\
\delta\mathbf{\lambda}%
\end{array}
\right)  \leq0\nonumber\\
&  \qquad\qquad\qquad\qquad\qquad\qquad\qquad\qquad\qquad\Updownarrow
\nonumber\\
&  \qquad\qquad\qquad\qquad\qquad\qquad\delta\mathbf{\lambda}^{t}%
\nabla_{\mathbf{x}}g\left(  \mathbf{x}_{m}\left(  \mathbf{c}\right)  \right)
\delta\mathbf{\lambda}\geq0,
\end{align}
i.e. $\left(  \mathbf{x}_{m}\left(  \mathbf{c}\right)  \mathbf{,\lambda}%
_{m}\left(  \mathbf{c}\right)  \right)  $ is a saddle point.

The Lagrangian can be expanded about an arbitrary point $\left(  Q_{a},\bar
{Q}_{a},\lambda_{a}\right)  $ as
\begin{align}
\mathcal{L}\left(  Q_{a}+\delta Q,\bar{Q}_{a}+\delta\bar{Q},\lambda_{a}%
+\delta\lambda\right)   &  =\mathcal{L}\left(  Q_{a},\bar{Q}_{a},\lambda
_{a}\right)  +D_{Q\bar{Q}\lambda}\mathcal{L}\left(  Q_{a},\bar{Q}_{a}%
,\lambda_{a}\right)  \left(
\begin{array}
[c]{c}%
\delta Q\\
\delta\bar{Q}\\
\delta\lambda
\end{array}
\right) \nonumber\\
&  +\frac{1}{2}\left(
\begin{array}
[c]{ccc}%
\delta Q & \delta\bar{Q} & \delta\lambda
\end{array}
\right)  D_{Q\bar{Q}\lambda}^{2}\mathcal{L}\left(  Q_{a},\bar{Q}_{a}%
,\lambda_{a}\right)  \left(
\begin{array}
[c]{c}%
\delta Q\\
\delta\bar{Q}\\
\delta\lambda
\end{array}
\right)
\end{align}
and
\begin{align}
D_{Q\bar{Q}\lambda}\mathcal{L}\left(  Q_{a}+\delta Q,\bar{Q}_{a}+\delta\bar
{Q},\lambda_{a}+\delta\lambda\right)   &  =D_{Q\bar{Q}\lambda}\mathcal{L}%
\left(  Q_{a},\bar{Q}_{a},\lambda_{a}\right) \\
&  +\left(
\begin{array}
[c]{ccc}%
\delta Q & \delta\bar{Q} & \delta\lambda
\end{array}
\right)  D_{Q\bar{Q}\lambda}^{2}\mathcal{L}\left(  Q_{a},\bar{Q}_{a}%
,\lambda_{a}\right)  .
\end{align}
Thus $\left(  Q_{a}+\delta Q,\bar{Q}_{a}+\delta\bar{Q},\lambda_{a}%
+\delta\lambda\right)  $ is a stationary point at a fixed $\gamma$ if
\begin{align}
D_{Q\bar{Q}\lambda}\mathcal{L}  &  \left(  Q_{a}+\delta Q,\bar{Q}_{a}%
+\delta\bar{Q},\lambda_{a}+\delta\lambda\right)  =\label{NR}\\
&  \left\{  D_{Q\bar{Q}\lambda}\mathcal{L}\left(  Q_{a},\bar{Q}_{a}%
,\lambda_{a}\right)  +\left(
\begin{array}
[c]{ccc}%
\delta Q & \delta\bar{Q} & \delta\lambda
\end{array}
\right)  D_{Q\bar{Q}\lambda}^{2}\mathcal{L}\left(  Q_{a},\bar{Q}_{a}%
,\lambda_{a}\right)  \right\}  =0.
\end{align}
Multiplying from the left by the generalized inverse $\left[  D_{Q\bar
{Q}\lambda}^{2}\mathcal{L}\left(  Q_{a},\bar{Q}_{a},\lambda_{a}\right)
\right]  ^{G}$ gives
\begin{align}
0  &  =D_{Q\bar{Q}\lambda}\mathcal{L}\left(  Q_{a},\bar{Q}_{a},\lambda
_{a}\right)  \left[  D_{Q\bar{Q}\lambda}^{2}\mathcal{L}\left(  Q_{a},\bar
{Q}_{a},\lambda_{a}\right)  \right]  ^{G}\\
&  +\left(
\begin{array}
[c]{ccc}%
\delta Q & \delta\bar{Q} & \delta\lambda
\end{array}
\right)  D_{Q\bar{Q}\lambda}^{2}\mathcal{L}\left(  Q_{a},\bar{Q}_{a}%
,\lambda_{a}\right)  \left[  D_{Q\bar{Q}\lambda}^{2}\mathcal{L}\left(
Q_{a},\bar{Q}_{a},\lambda_{a}\right)  \right]  ^{G}\\
&  \Leftrightarrow-D_{Q\bar{Q}\lambda}\mathcal{L}\left(  Q_{a},\bar{Q}%
_{a},\lambda_{a}\right)  \left[  D_{Q\bar{Q}\lambda}^{2}\mathcal{L}\left(
Q_{a},\bar{Q}_{a},\lambda_{a}\right)  \right]  ^{G}=\left(
\begin{array}
[c]{ccc}%
\delta Q & \delta\bar{Q} & \delta\lambda
\end{array}
\right)  P_{\left\{  \ker D_{Q\bar{Q}\lambda}^{2}\mathcal{L}\left(  Q_{a}%
,\bar{Q}_{a},\lambda_{a}\right)  \right\}  ^{c}}%
\end{align}
Thus
\begin{align}
&  \left(
\begin{array}
[c]{ccc}%
\delta Q & \delta\bar{Q} & \delta\lambda
\end{array}
\right)  P_{\left\{  \ker D_{Q\bar{Q}\lambda}^{2}\mathcal{L}\left(  Q_{a}%
,\bar{Q}_{a},\lambda_{a}\right)  \right\}  ^{c}}D_{Q\bar{Q}\lambda}%
^{2}\mathcal{L}\left(  Q_{a},\bar{Q}_{a},\lambda_{a}\right) \\
&  =\left(
\begin{array}
[c]{ccc}%
\delta Q & \delta\bar{Q} & \delta\lambda
\end{array}
\right)  D_{Q\bar{Q}\lambda}^{2}\mathcal{L}\left(  Q_{a},\bar{Q}_{a}%
,\lambda_{a}\right)  P_{\left\{  \operatorname{Ran}D_{Q\bar{Q}\lambda}%
^{2}\mathcal{L}\left(  Q_{a},\bar{Q}_{a},\lambda_{a}\right)  \right\}  }\\
&  =-D_{Q\bar{Q}\lambda}\mathcal{L}\left(  Q_{a},\bar{Q}_{a},\lambda
_{a}\right)  P_{\left\{  \operatorname{Ran}D_{Q\bar{Q}\lambda}^{2}%
\mathcal{L}\left(  Q_{a},\bar{Q}_{a},\lambda_{a}\right)  \right\}  }%
\end{align}
If a stationary point
\begin{equation}
\left(  Q_{\ast},\bar{Q}_{\ast},\lambda_{\ast}\right)  =\left(  Q_{a}+\delta
Q,\bar{Q}_{a}+\delta\bar{Q},\lambda_{a}+\delta\lambda\right)
\end{equation}
exists and is unique then
\begin{equation}
\left(
\begin{array}
[c]{ccc}%
\delta\bar{Q} & \delta Q & \delta\lambda
\end{array}
\right)  =-D_{Q\bar{Q}\lambda}\mathcal{L}\left(  Q_{a},\bar{Q}_{a},\lambda
_{a}\right)  \left[  D_{Q\bar{Q}\lambda}^{2}\mathcal{L}\left(  Q_{a},\bar
{Q}_{a},\lambda_{a}\right)  \right]  ^{-1}.
\end{equation}
The stationary point is a constrained minimum if
\begin{equation}
D_{Q}^{2}\mathcal{L}\left(  Q_{\ast},\bar{Q}_{\ast},\lambda_{\ast}\right)  >0.
\end{equation}
Many stationary points exist if $\ker D_{Q\bar{Q}\lambda}^{2}\mathcal{L}%
\left(  Q_{a},\bar{Q}_{a},\lambda_{a}\right)  \neq\left\{  0\right\}  $ and
the vectors
\begin{equation}
\left(  Q_{a}+\delta Q+q,\bar{Q}_{a}+\delta\bar{Q}+\bar{q},\lambda_{a}%
+\delta\lambda+\varkappa\right)
\end{equation}
are solutions to eq. $\left(  \ref{NR}\right)  ,$ where $\left(
\begin{array}
[c]{ccc}%
q & \bar{q} & \varkappa
\end{array}
\right)  \in\ker$ $\left\{  D^{2}\mathcal{L}\left(  Q_{a},\bar{Q}_{a}%
,\lambda_{a}\right)  \right\}  $ and
\begin{equation}
\left(
\begin{array}
[c]{ccc}%
\delta\bar{Q} & \delta Q & \delta\lambda
\end{array}
\right)  =-D\mathcal{L}\left(  Q_{a},\bar{Q}_{a},\lambda_{a}\right)  \left[
D^{2}\mathcal{L}\left(  Q_{a},\bar{Q}_{a},\lambda_{a}\right)  \right]  ^{G},
\end{equation}
where
\begin{equation}
\left[  D_{Q\bar{Q}\lambda}^{2}\mathcal{L}\left(  Q_{a},\bar{Q}_{a}%
,\lambda_{a}\right)  \right]  ^{G}%
\end{equation}
is the generalized inverse of $D_{Q\bar{Q}\lambda}^{2}\mathcal{L}\left(
Q_{a},\bar{Q}_{a},\lambda_{a}\right)  .$ If $D\mathcal{L}\left(  Q_{a},\bar
{Q}_{a},\lambda_{a}\right)  \notin\operatorname{Ran}\left\{  D^{2}%
\mathcal{L}\left(  Q_{a},\bar{Q}_{a},\lambda_{a}\right)  \right\}  $ then
there is no solution to eq. $\left(  \ref{NR}\right)  .$ This can only occur
if $D^{2}\mathcal{L}\left(  Q_{a},\bar{Q}_{a},\lambda_{a}\right)  $ is not surjective.

$\lceil$We can also express one particle positive trace class operators as
\begin{equation}
D^{1}=Q_{1}Q_{1}^{\dagger},
\end{equation}
where
\begin{equation}
Q_{1}\in\mathcal{B}_{2}\left(  \mathcal{H}^{1}\right)
\end{equation}
so that
\begin{equation}
\gamma\left(  \mathbf{x}\right)  =D^{1}\left(  \mathbf{x,x}\right)  =\int
Q_{1}\left(  \mathbf{x,y}\right)  \bar{Q}_{1}\left(  \mathbf{y,x}\right)
d\mathbf{y}%
\end{equation}
and
\begin{equation}
\int Q_{1}\left(  \mathbf{x,y}\right)  \bar{Q}_{1}\left(  \mathbf{y,x}\right)
d\mathbf{x}d\mathbf{y}=N\mathbf{.}%
\end{equation}
Further
\begin{equation}
Q_{1}=\sum_{1\leq k,l\leq M_{1}}Q_{1kl}\left\vert \varphi_{k}\right\rangle
\left\langle \varphi_{l}\right\vert
\end{equation}
and
\begin{align}
Q_{1}\left(  \mathbf{x,y}\right)   &  =\sum_{1\leq k,l\leq M_{1}}%
Q_{1kl}\varphi_{k}\left(  \mathbf{x}\right)  \bar{\varphi}_{l}\left(
\mathbf{y}\right) \\
\sum_{1\leq k,l\leq M_{1}}\left\vert Q_{1kl}\right\vert ^{2}  &  =N\nonumber\\
\int Q_{1}\left(  \mathbf{x,y}\right)  \bar{Q}_{1}\left(  \mathbf{y,x}\right)
d\mathbf{y}  &  \mathbf{=}\int\sum_{1\leq k,l\leq M_{1}}Q_{1kl}\varphi
_{k}\left(  \mathbf{x}\right)  \bar{\varphi}_{l}\left(  \mathbf{y}\right)
\sum_{1\leq r,s\leq M_{1}}\bar{Q}_{1rs}\bar{\varphi}_{r}\left(  \mathbf{y}%
\right)  \bar{\varphi}_{s}\left(  \mathbf{x}\right)  d\mathbf{y}\\
&  =\sum_{1\leq k,l,r,s\leq M_{1}}Q_{1kl}\bar{Q}_{1rs}\varphi_{k}\left(
\mathbf{x}\right)  \bar{\varphi}_{s}\left(  \mathbf{x}\right)  \int
\bar{\varphi}_{l}\left(  \mathbf{y}\right)  \bar{\varphi}_{r}\left(
\mathbf{y}\right)  d\mathbf{y}\\
&  =\sum_{1\leq k,l,s\leq M_{1}}Q_{1kl}\bar{Q}_{1ls}\varphi_{k}\left(
\mathbf{x}\right)  \bar{\varphi}_{s}\left(  \mathbf{x}\right)  \rfloor
\end{align}



\end{document}